\documentclass[11pt,a4paper]{article}

% --- Paquets pour la rigueur scientifique ---
\usepackage[utf8]{inputenc}
\usepackage[T1]{fontenc}
\usepackage[french]{babel}
\usepackage{amsmath, amssymb, amsfonts, amsthm}
\usepackage{geometry}
\geometry{margin=2.5cm}
\usepackage{hyperref}
\usepackage{booktabs}
\usepackage{graphicx}
\usepackage{enumitem}
\usepackage{xfrac}
\setcounter{tocdepth}{1}
% --- Métadonnées ---

\title{L’émergence de la réalité physique dans le cadre \\ du Logo Dynamique Projectif (LDP)}

\author{Cédric Laubscher\\
\small chercheur indépendant}

\date{2 février 2026}

\begin{document}

\maketitle

\begin{abstract}

Ce texte propose un cadre fondationnel, le \emph{Logo Dynamique Projectif} (LDP), visant à décrire l’émergence de la réalité physique à partir d’axiomes strictement relationnels, sans postuler \emph{a priori} l’existence d’un espace, d’un temps ou de particules. Le réel y est formalisé comme un réseau de fermetures logiques minimales, constitué d’entités pulsatoires et de leurs relations, régi par un principe d’optimisation de cohérence interne plutôt que par un ensemble de lois dynamiques externes.

Sous ces contraintes axiomatiques, une structure minimale \(\left(4,6\right)\) est déduite de conditions d’autonomie, de reproductibilité et de non‑hiérarchicité, et sert de brique élémentaire pour la description de l’électron. Le proton est, quant à lui, modélisé comme une architecture hiérarchique combinant trois régimes stationnaires locaux, une mer relationnelle dynamique et une surface active de taille finie. Dans le formalisme introduit ultérieurement, cette structure minimale \(\left(4,6\right)\) est représentée par un graphe complet à quatre sommets, doté de signes d’accord/désaccord sur ses six arêtes, pour lequel il existe une configuration de signes rendant tous les triangles logiquement cohérents, ainsi qu’une pulsation binaire globale \(s \leftrightarrow -s\) préservant cette cohérence.

Cette configuration réalise explicitement les axiomes~1 à~3 au niveau élémentaire et est identifiée au régime stationnaire de l’électron au repos. L’architecture proposée permet de réinterpréter plusieurs constantes adimensionnées comme des rapports de cohérence interne : \(\hbar\) est lue comme la granularité minimale de l’action, la constante de structure fine \(\alpha\) comme la fraction du budget de cohérence projetée sur la surface, via une surface active \(R_{\mathrm{surf}}^{\left(\varphi\right)}\) qui fait intervenir le rapport de masses \(\mu\), le contenu de valence \(r_{\mathrm{valence}}\) et le nombre d’or \(\varphi\). De même, \(k_B\) est interprétée comme une granularité thermique associée à la fraction dynamique et à une fuite de cohérence \(\varepsilon\), tandis qu’un couplage gravitationnel effectif est proposé comme traduction métrique de cette fuite minimale.

Plusieurs accords numériques sont obtenus dans des marges d’incertitude contrôlées, sans introduction de paramètres continus supplémentaires. Toutefois, certaines constructions demeurent conjecturales ou partiellement ajustées sur les données expérimentales, en particulier pour \(k_B\), la gravitation et la cosmologie. Le texte s’achève sur une esquisse spéculative de scénarios de genèse et de recyclage des fermetures, où des épisodes de type Big Bang et des singularités gravitationnelles extrêmes sont réinterprétés comme des événements de contrainte dans un champ logique pré‑physique. L’ensemble doit être compris comme une proposition de cadre fondationnel en cours d’élaboration, plutôt que comme une théorie pleinement établie.

\end{abstract}

\newpage

% Génère la table des matières
\tableofcontents

\newpage

\section*{Introduction}
\addcontentsline{toc}{section}{Introduction}

Ce travail adopte un postulat méthodologique distinct : plutôt que de rechercher de nouvelles lois physiques définies sur un espace-temps préalablement donné, il propose de remonter à un niveau plus fondamental et de reconstruire la réalité physique à partir d’axiomes strictement relationnels. Le \emph{Logo Dynamique Projectif} (LDP) constitue la formalisation de cette approche en introduisant un ensemble restreint de contraintes sur les modalités d’organisation des relations, de manière à permettre l’émergence de cycles fermés et reproductibles. Sous ces contraintes (autonomie logique des structures, stabilité de la référence, absence de hiérarchie a priori imposée, optimisation de la cohérence) émerge une structure minimale, notée \(\left(4,6\right)\), qui joue le rôle de brique élémentaire du modèle.

Une part substantielle du manuscrit est consacrée, d’une part, à la dérivation de cette structure minimale et à l’examen détaillé de ses propriétés, et, d’autre part, à la construction d’architectures hiérarchiques plus complexes, identifiées aux nucléons, en particulier au proton. Ces architectures relationnelles permettent de réinterpréter plusieurs constantes adimensionnées (\(\hbar\), \(\alpha\), \(k_B\), un couplage gravitationnel effectif) comme des rapports de cohérence interne mettant en jeu des entiers combinatoires, le rapport de masses \(\mu\), un paramètre de fuite \(\varepsilon\) ainsi que le nombre d’or \(\varphi\).

Cette démarche conduit également à une reformulation relationnelle du temps propre, de la gravitation et de certains phénomènes thermiques, et se prolonge par une esquisse de cadre cosmologique dans lequel les fermetures stables émergent et se recyclent comme manifestations d’événements de contrainte au sein d’un champ logique pré-physique. Néanmoins, le travail ne revendique pas le statut de théorie achevée et close : certains résultats sont rigoureusement démontrés à partir des axiomes, tandis que d’autres demeurent conjecturaux ou partiellement ajustés aux données expérimentales. Plusieurs domaines, en particulier la cosmologie détaillée et les phénomènes hors équilibre, ne sont abordés qu’à titre exploratoire. Le LDP doit ainsi être considéré comme une proposition de cadre fondationnel en cours d’élaboration, appelée à être complétée, précisée et soumise à l’épreuve de la dérivation de prédictions nouvelles et falsifiables.

\newpage

\section{Limites des concepts physiques usuels}

Toute théorie physique repose sur un ensemble de concepts définis au sein d’un domaine de validité précisément circonscrit. Lorsqu’une telle théorie est extrapolée jusqu’à des régimes extrêmes, au-delà desquels ses concepts fondamentaux cessent de pouvoir être spécifiés de manière cohérente, la difficulté ne relève plus uniquement des limitations expérimentales, mais met en cause la consistance logique de ces concepts eux-mêmes.

Dans ces régimes, les dispositifs de mesure peuvent devenir inopérants, mais l’enjeu ne réside plus dans la seule précision des mesures : c’est la possibilité même d’attribuer un contenu conceptuel non ambigu aux notions d’espace, de temps, d’énergie ou d’état physique qui est questionnée. Le problème se déplace alors du plan empirique vers le plan proprement conceptuel.

L’objectif n’est donc pas de procéder à de simples corrections ou ajustements marginaux des théories existantes, mais d’interroger les conditions de possibilité de leur formulation. Il s’agit de remonter en deçà de toute physique déjà constituée, vers un cadre conceptuel plus élémentaire, au sein duquel les grandeurs physiques usuelles ne sont plus posées comme allant de soi.

Dans une telle perspective, il devient nécessaire de suspendre l’usage des grandeurs métriques ainsi que des structures mathématiques héritées des théories établies. L’espace, le temps, l’énergie et, plus généralement, l’ensemble des grandeurs mesurables ne peuvent plus être postulés comme fondements : ils doivent, au contraire, être traités comme des structures émergentes possibles, dont il convient d’élucider les conditions d’apparition. La question de l’existence doit, dans ce contexte, être reformulée dans un cadre dépourvu de toute métrique a priori.

À ce stade, aucune hypothèse formelle précise n’a encore été introduite. La démarche demeure exclusivement d’ordre logique : il s’agit d’identifier les conditions minimales qu’une structure doit satisfaire pour pouvoir être dite exister de manière persistante, indépendamment de tout postulat métrique préalable. Les axiomes destinés à formaliser ces conditions seront explicités dans la suite du texte.

\section{Reformuler la question de l’existence}

Une fois toute métrique et toute possibilité de mesure suspendues, la problématique de l’existence se reconfigure de manière radicale : comment distinguer « quelque chose » du néant lorsqu’aucune propriété mesurable ne peut être invoquée comme critère de différenciation ?

En l’absence d’espace, de temps, d’énergie ou de toute autre grandeur physique ou abstraite, aucune caractéristique intrinsèque ne peut être mobilisée pour établir une distinction ontologique. Toute tentative de définir une entité comme « ce qui occupe un lieu », « ce qui dure » ou « ce qui possède une énergie » devient alors dépourvue de pertinence, dans la mesure où ces notions présupposent déjà un cadre métrique et conceptuel qui, par hypothèse, n’est pas disponible. Une telle configuration conduit nécessairement à une situation d’indifférenciation absolue : en l’absence de tout critère discriminant, rien ne permet de distinguer une entité d’une autre, ni même de distinguer l’existence du néant.

Dans ce contexte, une existence strictement instantanée ne saurait constituer une existence au sens fort. Une entité dépourvue de toute forme de persistance serait logiquement indiscernable du néant : si aucun dispositif, même conceptuel, ne permet de la retrouver, de la réidentifier ou de la relier à une occurrence ultérieure, elle ne peut être distinguée d’une absence totale. Le problème fondamental ne consiste donc pas à énumérer ou caractériser « ce qui existe », mais à déterminer les conditions minimales sous lesquelles quelque chose peut exister de manière persistante, c’est-à-dire se maintenir comme distinct du néant.

Sous ces contraintes, ni un état strictement immuable ni un changement unique et non récurrent ne peuvent constituer un fondement satisfaisant de l’existence. Un état immuable, dépourvu de toute possibilité de différence interne ou externe, reste indiscernable d’un néant uniformément identique à lui-même. De même, un changement unique, qui ne se reproduirait jamais, ne peut être identifié qu’au prix d’une contradiction : en l’absence de répétition ou de récurrence, il ne laisse aucune trace logique ou structurelle permettant de l’individualiser. La seule configuration cohérente consiste alors dans l’émergence d’une différence susceptible de se reproduire, c’est-à-dire d’une structure répétable assurant les conditions minimales d’identification, de distinction et, par conséquent, d’existence.

\section{Alternance minimale et entité informationnelle}

Dans ce cadre pré-métrique, l’existence ne peut être fondée ni sur une substance, ni sur une localisation, ni sur une durée, dans la mesure où ces notions présupposent déjà la disponibilité d’un espace, d’un temps ou d’une énergie. Ce qui peut soutenir une distinction persistante ne peut dès lors être qu’une structure logique minimale : une différence qui se manifeste et se reproduit.

Sous ces contraintes, une alternance logique minimale s’impose comme la forme la plus élémentaire d’existence cohérente. Cette alternance ne se déploie ni temporellement ni spatialement : en l’absence de toute métrique, il est impossible de lui attribuer une période, une fréquence, une trajectoire ou une amplitude. Elle ne possède aucune intensité définissable ; elle ne peut être caractérisée que par le fait même « qu’il se passe quelque chose plutôt que rien » et par le caractère répétable de ce « quelque chose ».

Cette alternance minimale peut alors être définie comme une entité informationnelle élémentaire. Elle ne l’est pas en vertu du transport d’un contenu sémantiquement déterminé, mais parce qu’elle réalise la forme la plus simple de distinction reproductible dans un cadre dépourvu de toute métrique. L’« information », en ce sens, ne désigne pas un message, mais la condition de possibilité même d’une différence qui se maintient.

\section{Multiplicité minimale et rôle de la relation}

Dans un cadre conceptuel où aucune métrique ni aucune propriété intrinsèque ne peut être définie, une entité informationnelle élémentaire ne peut se différencier d’une autre qu’en vertu de son occurrence, c’est-à-dire du fait même de « se produire ». Deux entités de ce type sont donc, par construction, indiscernables si l’on ne dispose d’aucun autre critère de différenciation que leur existence en tant que telles.

Dans ces conditions, une entité informationnelle strictement isolée demeure logiquement indéfinissable. Pour qu’une distinction soit opératoire, il ne suffit pas qu’un événement survienne : il est nécessaire qu’une différence soit repérable relativement à une autre entité, à un état, ou à une structure de référence. Toute opération de distinction requiert ainsi un référentiel, explicite ou implicite.

Dans ce cadre, la relation ne constitue pas un mécanisme surajouté entre des entités préalablement constituées ; elle représente la condition même de possibilité sous laquelle une entité informationnelle peut être dite exister de manière définissable. En l’absence de relation, aucune entité ne peut être isolée conceptuellement : elle se confond alors avec un fond indifférencié assimilable au néant. Par conséquent, l’existence d’une entité informationnelle ne peut être envisagée que dans le contexte d’une multiplicité minimale, au sein de laquelle la relation assure la fonction de support de la distinction.

Toutefois, une relation unique demeure insuffisante pour fonder une distinction stable et persistante. Une seule liaison ne permet pas de stabiliser la référence nécessaire à la reproduction de la différence : toute modification de l’un des termes ou de la liaison elle-même annule immédiatement la possibilité d’identification. Il devient alors nécessaire de postuler une structure relationnelle plus riche, impliquant au moins un réseau minimal de relations, afin de garantir la pérennité et la robustesse de la distinction.

\section{Contraintes sur une structure relationnelle élémentaire}

L’analyse précédente permet de dégager un ensemble de contraintes logiques auxquelles doit satisfaire toute structure relationnelle élémentaire admissible.

\paragraph{Reproductibilité de la distinction.}
La différence qui fonde l’existence doit être intrinsèquement reproductible. Une structure qui ne permettrait pas la réitération de l’opération distinctive ne peut pas soutenir une existence persistante, dans la mesure où l’instance considérée ne pourrait être reconduite ni confirmée au sein de la structure.

\paragraph{Stabilité de la référence.}
La référence à partir de laquelle la distinction est définie doit présenter une stabilité suffisante. Une structure qui dépendrait d’un élément externe, susceptible de s’abolir sans laisser de trace dans la structure elle-même, serait incapable de garantir cette stabilité référentielle, et, par conséquent, l’identité de la distinction.

\paragraph{Autonomie logique.}
La structure élémentaire doit être logiquement autosuffisante. Elle ne peut pas requérir le recours à un « arrière-plan » indéterminé pour assurer sa cohérence, faute de quoi la distinction qu’elle porte demeurerait conditionnée par un extérieur non explicité et ne pourrait être entièrement spécifiée à partir des seules relations internes.

Ces contraintes impliquent qu’une structure élémentaire soit fermée : l’ensemble des références nécessaires à la définition, à la reproduction et au maintien de la distinction doit être interne à la structure. Une organisation ouverte, dépendant de repères externes non intégrés, contreviendrait à l’exigence d’autonomie logique.

Enfin, parmi toutes les structures fermées satisfaisant ces conditions, seules peuvent être retenues au niveau fondamental celles qui les réalisent avec un nombre minimal de relations. La minimalité de la structure n’est pas posée comme une hypothèse arbitraire : elle résulte du processus même de sélection imposé par l’ensemble des contraintes précédentes.

\subsection*{Axiomes fondamentaux du cadre logique}

Les contraintes précédemment introduites ne constituent pas des conditions techniques adjointes a posteriori. Elles peuvent être formulées sous la forme de quatre axiomes structurants. Dans ce qui suit, une configuration logique est modélisée par un graphe fini \(G = (V,E)\), dont les sommets représentent des entités informationnelles et les arêtes des relations binaires d’accord/désaccord entre ces entités. À chaque arête \(\{i,j\} \in E\) est associé un signe \(s_{ij} \in \{+1,-1\}\), ce qui permet de représenter l’alternance minimale sous la forme d’une dynamique discrète portant sur les signes. Les cycles de longueur 3 (triangles) jouent le rôle de motifs élémentaires de fermeture logique.

\subsection*{Axiome 1 – Pulsation binaire minimale}

L’existence se manifeste par une alternance élémentaire accord/désaccord, répétable et définissable indépendamment de toute métrique préalablement donnée. Formellement, il existe une dynamique discrète
\[
F : \Sigma^E \rightarrow \Sigma^E,\qquad \Sigma = \{+1,-1\}
\]
telle qu’il existe au moins un état de signes \(s^\star\) satisfaisant
\[
F(s^\star) = -s^\star,\qquad F(-s^\star) = s^\star.
\]
Cette pulsation binaire ne présuppose aucune métrique temporelle : elle exprime uniquement l’existence d’un cycle logique de longueur 2 dans l’espace des configurations.

\subsection*{Axiome 2 – Cohérence triangulaire}

Toute structure stable requiert une fermeture logique : les violations de cette fermeture (triangles ouverts, ruptures de symétrie non compensées) diminuent la stabilité. Pour une configuration \(\sigma = (G,s)\), on considère l’ensemble des triplets \(\{i,j,k\} \subset V\) :
\begin{itemize}
  \item le triangle \(\{i,j,k\}\) est dit \emph{présent} si les trois arêtes \(\{i,j\}, \{j,k\}, \{i,k\}\) appartiennent à \(E\) ;
  \item il est dit \emph{cohérent} si, lorsqu’il est présent, le produit des signes vérifie
  \[
  s_{ij}\, s_{jk}\, s_{ik} = +1
  \]
  (boucle sans rupture non compensée) ;
  \item il est dit \emph{violé} si au moins une des trois arêtes est absente, ou si le produit vaut \(-1\).
\end{itemize}

On définit alors \(\nu(\sigma)\) comme la fraction de triangles violés :
\[
\nu(\sigma) = \frac{\#\{\text{triangles violés de } G\}}{\binom{|V|}{3}},
\]
avec la convention \(\nu(\sigma) = 0\) si \(|V| < 3\). Les structures stables sont celles qui minimisent la quantité \(\nu(\sigma)\).

\subsection*{Axiome 3 – Complétude minimale}

Une structure fermée doit satisfaire un critère de complétude interne : l’ensemble des références nécessaires à la définition et au maintien de la distinction doit être entièrement contenu dans la structure elle-même, sans recours à un arrière-plan implicite ou à des paramètres externes. 

On appellera \emph{structure stationnaire élémentaire} toute configuration \(\sigma = \left(G,s\right)\) qui satisfait simultanément les conditions suivantes :
\begin{itemize}
  \item le graphe \(G\) est connexe et complet (toute paire de sommets est reliée par une arête) ;
  \item il existe une affectation de signes \(s\) telle que \(\nu\left(\sigma\right) = 0\), c’est-à-dire que l’ensemble des triangles de \(G\) est intégralement cohérent ;
  \item la cohérence triangulaire constitue un unique bloc : le sous-graphe induit par l’union de tous les triangles cohérents est connexe et contient au moins deux triangles partageant une arête commune ;
  \item les sommets de \(G\) sont structurellement équivalents, ce qui se traduit par une symétrie interne globale et l’absence de toute hiérarchie intrinsèque entre les sommets ;
  \item la dynamique \(F\) admet une pulsation binaire \(s \leftrightarrow -s\) qui préserve la cohérence triangulaire, de sorte que l’inversion globale des signes ne modifie pas la structure de cohérence.
\end{itemize}

Une structure stationnaire élémentaire est dite \emph{minimale} lorsqu’aucune structure de cardinalité strictement inférieure ne satisfait l’ensemble de ces conditions. Dans ce cadre, la première structure minimale admissible apparaît pour \(n = 4\), et est réalisée par le graphe complet à quatre sommets et six arêtes, noté \(\left(4,6\right)\).

\subsection*{Axiome 4 – Optimisation logique}

Parmi l’ensemble des configurations admissibles, seules sont susceptibles de persister dynamiquement celles qui maximisent une fonction d’optimisation \(\Omega\) combinant les critères suivants :
\begin{itemize}
  \item la minimisation des violations \(\nu\left(\sigma\right)\) (réduction du nombre de triangles incohérents) ;
  \item la maximisation de la connectivité relative \(\rho\left(\sigma\right)\) ;
  \item la réalisation conjointe de la fermeture et de la minimalité structurelle.
\end{itemize}

On supposera que la fonction d’optimisation \(\Omega\left(\sigma\right)\) dépend uniquement des grandeurs \(\nu\left(\sigma\right)\), \(\rho\left(\sigma\right)\) et d’un indicateur de minimalité, et qu’elle vérifie les propriétés suivantes :
\begin{itemize}
  \item elle est strictement décroissante en \(\nu\left(\sigma\right)\) (un nombre plus faible de triangles violés induit une plus grande stabilité structurelle) ;
  \item elle est croissante en \(\rho\left(\sigma\right)\) (une connectivité relative plus élevée correspond à une exploitation plus efficiente du budget relationnel interne) ;
  \item elle est croissante lorsqu’une configuration réalise le statut de structure stationnaire élémentaire minimale.
\end{itemize}

Dans ce cadre, la structure \(\left(4,6\right)\) émerge comme le premier maximum local de \(\Omega\) au niveau fondamental, et joue, à ce titre, le rôle de brique élémentaire du modèle.

\section{Exclusion des configurations insuffisantes}

La définition de cette structure relationnelle minimale procède par l’exclusion systématique des configurations qui ne satisfont pas l’ensemble des contraintes spécifiées.

\paragraph{Une seule relation.}
Une relation isolée ne peut générer une distinction reproductible sans recours à un repère externe. La référence ne peut y être stabilisée de manière autonome. Une telle configuration contrevient ainsi à l’exigence d’autonomie logique et doit, à ce titre, être écartée.

\paragraph{Deux relations.}
Un système constitué de deux relations peut instaurer une forme d’alternance, mais celle-ci demeure dépendante d’un repère implicite extérieur. La structure qui en résulte reste soit ouverte, soit dégénérée : elle ne garantit ni la clôture, ni la stabilité de la référence. Une telle configuration se révèle donc insuffisante pour satisfaire simultanément les exigences de reproductibilité, de stabilité et d’autonomie.

\paragraph{Trois relations.}
Une organisation fondée sur trois relations présente une complexité combinatoire accrue, mais impose inévitablement un rôle structurel privilégié à l’une d’entre elles. Cette asymétrie introduit une hiérarchie implicite, incompatible avec l’exigence d’une structure élémentaire non hiérarchisée. La référence n’y est pas distribuée de façon homogène, de sorte que la distinction reste dépendante d’un élément particulier. Une telle configuration ne peut dès lors constituer un fondement minimal.

Sur la base de ces exclusions, il apparaît qu’aucune structure comportant moins de quatre relations ne peut satisfaire simultanément les conditions de reproductibilité de la distinction, de stabilité de la référence, d’autonomie logique et d’absence de hiérarchie interne.

Cette analyse conceptuelle se traduit, dans le langage des graphes, par un seuil combinatoire simple : pour \(n = 1,2\), il n’existe aucun triangle ; pour \(n = 3\), il n’existe qu’un seul triangle, qui ne peut former un réseau recouvrant. Dans tous ces cas, il est impossible de réaliser une cohérence triangulaire distribuée, caractérisée par la présence d’au moins deux triangles partageant une arête et par une union connexe des triangles cohérents. Une structure stationnaire élémentaire, au sens défini précédemment, ne peut donc émerger qu’à partir de quatre entités.

\section{Première structure admissible et introduction des graphes logiques}

Le premier cas non exclu par les contraintes précédemment établies correspond à une structure relationnelle comportant au moins quatre relations internes. Il s’agit dès lors d’identifier et de caractériser explicitement cette structure minimale.

À cette fin, il est nécessaire d’introduire un langage formel permettant de représenter des entités indifférenciées ainsi que leurs relations, sans postuler a priori de métrique ni de structure physique sous‑jacente. Les graphes logiques fournissent un tel formalisme : les entités informationnelles y sont représentées par les sommets, et les relations qui les lient par les arêtes.

Dans ce cadre, une structure relationnelle fermée est formalisée par un graphe dont la cohérence interne ne dépend d’aucun sommet ni d’aucune arête extérieure. Le régime de stabilité maximale, entendu comme le cas de distribution homogène de la référence, correspond alors à une structure dans laquelle chaque entité est reliée à toutes les autres, c’est‑à‑dire à un graphe complet.

L’introduction de ce formalisme rend possible le comptage explicite du nombre d’entités et de relations nécessaires pour satisfaire l’ensemble des contraintes posées. Une structure admissible doit comporter au minimum quatre relations, mais ce nombre de relations ne peut être considéré indépendamment du nombre d’entités : une structure fermée et non hiérarchisée impose en effet que les relations soient distribuées de manière homogène, de sorte qu’aucune entité ne joue un rôle structurel privilégié.

Avec moins de quatre entités, cette exigence de distribution homogène conduit inévitablement à des configurations dégénérées ou asymétriques, incompatibles avec les conditions de fermeture et d’autonomie logique. À partir de quatre entités, en revanche, il devient possible de construire un graphe complet : chaque entité est reliée à l’ensemble des autres, sans dépendance externe et sans hiérarchie implicite.

Pour un régime stationnaire local comportant \(n\) entités, le nombre de relations internes est donné par
\[
r(n) = \frac{n(n-1)}{2}.
\]
Pour \(n = 4\), on obtient \(r(4) = 6\), ce qui fixe la structure minimale par le couple \((4,6)\).

Cette configuration \((4,6)\) constitue le premier cas satisfaisant simultanément l’ensemble des contraintes et définit, à ce titre, la structure relationnelle minimale admissible. Dans le formalisme ainsi introduit, cette configuration \((4,6)\) n’est pas seulement la première à mettre en œuvre une distribution homogène des relations : elle est également la première pour laquelle il existe une affectation de signes rendant tous les triangles cohérents (\(\nu = 0\)), une pulsation binaire globale \(s \leftrightarrow -s\) préservant cette cohérence, ainsi qu’un réseau de triangles recouvrants couvrant l’intégralité du graphe. Elle réalise ainsi, de manière explicite, les axiomes~1 à~3 au niveau élémentaire.

\section{Propriétés internes de la structure \texorpdfstring{\(\left(4,6\right)\)}{(4,6)}}

La structure relationnelle minimale associée au couple \(\left(4,6\right)\) présente un ensemble de propriétés internes qui découlent directement de sa définition formelle. En tant que graphe logique complet à quatre entités, elle se caractérise d’abord par une distribution homogène des relations : chaque entité est connectée à toutes les autres, et aucune relation ne se voit conférer de rôle privilégié.

Cette symétrie interne annule toute hiérarchie implicite. La référence nécessaire à la distinction n’est plus localisée sur un élément particulier, mais distribuée de manière équivalente sur l’ensemble de la structure. La possibilité même de distinguer « quelque chose » du néant repose ainsi sur la cohérence globale du réseau de relations, plutôt que sur les propriétés d’un point isolé.

Cette distribution homogène induit une forme de stabilité logique spécifique pour la structure \(\left(4,6\right)\). Toute opération de suppression ou de modification d’une relation rompt immédiatement la clôture du graphe et réintroduit soit une dépendance externe, soit une asymétrie interne, ce qui rend la distinction non reproductible. Réciproquement, l’ajout d’une relation supplémentaire est impossible sans introduire une redondance, puisque toutes les paires d’entités sont déjà connectées.

Dans le formalisme triangulaire, la structure \(\left(4,6\right)\) contient quatre triangles distincts dont l’union recouvre l’intégralité du graphe et qui s’intersectent deux à deux. Il existe des affectations de signes sur les six arêtes, notamment l’affectation uniforme \(s_{ij}=+1\), pour lesquelles l’ensemble de ces triangles est cohérent, de sorte que \(\nu\left(\sigma\right)=0\) pour au moins un état. La dynamique \(F\) peut alors être spécifiée de manière à réaliser une pulsation globale \(s \leftrightarrow -s\) : chaque cycle de pulsation envoie un état cohérent sur un autre état cohérent, sans création de triangle violé. La structure \(\left(4,6\right)\) satisfait ainsi simultanément les contraintes de cohérence triangulaire parfaite, de symétrie interne et de pulsation minimale.

La structure \(\left(4,6\right)\) peut dès lors être interprétée comme un point fixe logique : elle est à la fois minimale et saturée relativement aux exigences de reproductibilité, d’autonomie et de non-hiérarchie. Aucune structure plus simple ne peut satisfaire simultanément ces conditions, et toute complexification locale altère son caractère minimal sans introduire de gain de cohérence fondamentale.

\section{Régime stationnaire logique et persistance}

Le franchissement du seuil défini par la structure \(\left(4,6\right)\) ne correspond ni à l’émergence d’une nouvelle entité au sens ontologique ou substantiel, ni à l’introduction d’un mécanisme additionnel. Il signale l’instauration d’un régime contraint dans lequel l’alternance informationnelle, jusque-là seulement possible, se trouve nécessairement maintenue par la structure relationnelle elle-même.

Plus précisément, la présence d’un graphe complet \(\left(4,6\right)\) associée à une affectation de signes cohérente assure l’existence d’une orbite de période 2 pour la dynamique \(F\) sur l’espace des configurations de signes : l’alternance accord/désaccord \(s \leftrightarrow -s\) devient un cycle logique intrinsèquement imposé par la structure. Ce régime stationnaire ne requiert aucune métrique temporelle : il exprime exclusivement la possibilité, imposée par la clôture triangulaire, de maintenir indéfiniment la distinction au moyen d’une pulsation interne.

La fermeture complète de la structure minimale implique, en effet, une circulation interne de la distinction : la différence constitutive de l’existence ne peut ni s’interrompre ni se dissiper sans entraîner une rupture de la cohérence globale du graphe. La persistance qui en découle n’est pas contingente ; elle est d’ordre structurel. Elle résulte directement de la saturation relationnelle : tant que la structure \(\left(4,6\right)\) est préservée, la distinction qu’elle supporte ne peut disparaître sans engendrer une contradiction interne.

Au-delà de ce seuil, la distinction ne se limite plus à être simplement reproductible au sens minimal ; elle est stabilisée par l’organisation fermée qui la supporte. Cette stabilisation ne se fonde sur aucun repère externe, aucune métrique préalable, ni aucune référence explicite au temps ou à l’espace. Elle est garantie exclusivement par la cohérence interne de la structure.

On peut dès lors qualifier ce régime de \emph{régime stationnaire logique} : un état dans lequel la distinction est maintenue sans apport extérieur, par la seule redistribution continue des relations qui constituent la structure fermée. La circulation interne impose une régularité intrinsèque : la distinction ne peut être maintenue que sous la forme d’une alternance structurée, contrainte par la fermeture de l’ensemble.

Cette alternance se manifeste uniquement comme une périodicité logique, c’est-à-dire comme la répétition nécessaire d’un motif de distinction au sein de la structure fermée, en l’absence de toute référence à une métrique temporelle ou spatiale.

\section{Fonction d’optimisation logique et sélection des configurations}

Le franchissement du seuil défini par la structure \(\left(4,6\right)\) ne suffit pas, à lui seul, à déterminer une configuration unique : un grand nombre d’organisations relationnelles plus complexes demeurent, en principe, logiquement possibles. Le problème se reformule alors en termes de sélection : selon quels critères certaines configurations peuvent-elles se maintenir de manière durable, tandis que d’autres ne subsistent que transitoirement ou sont éliminées du domaine des états accessibles ?

Dans le cadre conceptuel adopté ici, ce processus de sélection n’est pas attribué à une dynamique imposée de l’extérieur. Il est plutôt interprété comme l’effet d’un principe d’optimisation interne : parmi l’ensemble des configurations admissibles, seules sont stables celles qui maximisent un certain compromis entre cohérence, connectivité et minimalité, sous les contraintes de fermeture précédemment établies.

Pour expliciter ce principe, on considère une configuration \(\sigma\) caractérisée par les paramètres suivants :
\begin{itemize}
  \item \(n\left(\sigma\right)\) : le nombre d’entités informationnelles impliquées ;
  \item \(r\left(\sigma\right)\) : le nombre de relations internes effectivement actives ;
  \item \(\nu\left(\sigma\right)\) : le nombre de violations de la cohérence triangulaire, au sens de l’axiome~2 (triangles ouverts, fermetures incompatibles, ruptures de symétrie non compensées) ;
  \item \(\rho\left(\sigma\right) = r\left(\sigma\right)/r_{\max}\left(n\right)\) : la densité relationnelle normalisée, où \(r_{\max}\left(n\right) = n\left(n-1\right)/2\) désigne le nombre de relations d’un graphe complet à \(n\) entités ;
  \item un type de régime \(\tau\left(\sigma\right)\), stationnaire ou dynamique, selon que la structure réalise localement une fermeture saturée ou une connectivité partielle.
\end{itemize}

Plus précisément, \(\nu\left(\sigma\right)\) est défini à partir des triangles du graphe logique sous-jacent à \(\sigma\). Pour une configuration donnée \(\sigma\), on note \(G(\sigma) = (V,E)\) le graphe dont les sommets représentent les entités informationnelles et dont les arêtes codent les relations actives entre celles-ci. Un triplet \((i,j,k)\) de sommets constitue un \emph{triangle cohérent} si, d’une part, les trois arêtes \((i,j)\), \((j,k)\) et \((i,k)\) appartiennent à \(E\) et si, d’autre part, l’alternance accord/désaccord associée à ces relations satisfait la contrainte de fermeture de l’axiome~2, c’est-à-dire qu’aucune « rupture » non compensée n’apparaît dans le cycle ainsi formé.

À l’inverse, on qualifie de \emph{triangle violé} tout triplet \((i,j,k)\) pour lequel au moins l’une des trois arêtes est absente, ou bien pour lequel la combinaison des signes accord/désaccord induit une rupture non compensée au sein de la boucle. La quantité \(\nu(\sigma)\) est alors définie comme le nombre total de triangles violés dans \(G(\sigma)\), normalisé par le nombre total de triangles possibles \(\binom{n(\sigma)}{3}\) :
\[
\nu(\sigma)
 \;=\;
 \frac{\#\{\text{triangles violés dans } G(\sigma)\}}{\binom{n(\sigma)}{3}}.
\]

Un exemple élémentaire est fourni par le graphe à trois sommets. Lorsque ces trois sommets sont reliés de manière cohérente (trois arêtes présentes, avec une configuration d’accords / désaccords compatible avec la contrainte de fermeture), on obtient \(\nu(\sigma)=0\). Si une seule arête est absente, le triplet \((i,j,k)\) constitue alors un triangle ouvert et l’on a \(\nu(\sigma)=1\). 

Dans une structure minimale de type \((4,6)\), tous les triangles sont présents et vérifient la contrainte de fermeture, ce qui impose \(\nu(\sigma)=0\) et rend compte de sa stabilité logique maximale. Le cas à trois sommets montre ainsi qu’un triangle isolé peut être localement cohérent \(\bigl(\nu(\sigma)=0\bigr)\), mais, comme discuté en section~6, il ne réalise pas une complétude triangulaire minimale : il ne contient qu’un seul cycle et ne permet pas de distribuer la référence au sein d’un réseau de triangles recouvrants. La première structure stationnaire élémentaire au sens des axiomes est donc bien la configuration \((4,6)\).

Dans ce qui suit, on supposera que la fonction d’optimisation \(F\), associée à une configuration \(\sigma\), dépend uniquement de \(\nu(\sigma)\), de la densité relationnelle \(\rho(\sigma)\) et d’un indicateur de minimalité \(m(\sigma)\), selon
\[
F(\sigma)
= w_{\mathrm{coh}}\, f_{\mathrm{coh}}\bigl(\nu(\sigma)\bigr)
+ w_{\mathrm{conn}}\, f_{\mathrm{conn}}\bigl(\rho(\sigma)\bigr)
+ w_{\mathrm{min}}\, m(\sigma),
\]
où les coefficients \(w_{\mathrm{coh}}, w_{\mathrm{conn}}, w_{\mathrm{min}} > 0\) sont fixés une fois pour toutes, et où les formes fonctionnelles de \(f_{\mathrm{coh}}\) et \(f_{\mathrm{conn}}\) sont choisies de manière à satisfaire les axiomes~2 et~4.

On introduit alors une fonction d’optimisation \(\Omega(\sigma)\), faisant office de critère de sélection logique :
\[
\Omega(\sigma)
= w_{\mathrm{coh}}\!\left(1 - \frac{\nu(\sigma)}{n(\sigma)^2}\right)
+ w_{\mathrm{conn}}\, \rho(\sigma)
+ w_{\mathrm{min}}\, \delta_{\text{minimalité}}(\sigma),
\]
où \(w_{\mathrm{coh}}, w_{\mathrm{conn}}, w_{\mathrm{min}}\) désignent des poids strictement positifs, et où \(\delta_{\text{minimalité}}(\sigma) = 1\) si, et seulement si, la configuration réalise une fermeture sans recours à une structure plus riche, et vaut \(0\) dans le cas contraire. Le premier terme favorise les configurations minimisant les violations de cohérence \(\nu(\sigma)\) (au sens de l’axiome~2), le second terme privilégie une connectivité suffisante pour garantir la circulation de la distinction, et le troisième terme sélectionne les organisations satisfaisant au mieux le principe de minimalité imposé par la structure \((4,6)\).

Dans l’état actuel du modèle, on adopte pour \(f_{\mathrm{coh}}\) et \(f_{\mathrm{conn}}\) des formes linéaires normalisées :
\[
f_{\mathrm{coh}}\left(\nu\right) = 1-\nu,
\qquad
f_{\mathrm{conn}}\left(\rho\right) = \rho,
\]
de telle sorte que la contribution de cohérence décroît linéairement avec le taux de triangles violés, tandis que la contribution de connectivité croît linéairement avec la densité relationnelle.

Un choix naturel des poids consiste à poser
\[
w_{\mathrm{coh}} = 1,
\qquad
w_{\mathrm{conn}} = 1,
\qquad
w_{\mathrm{min}} = 1,
\]
ce qui conduit à
\[
\Omega\left(\sigma\right)
= \left(1-\nu\left(\sigma\right)\right)
+ \rho\left(\sigma\right)
+ \delta_{\text{minimalité}}\!\left(\sigma\right).
\]
Ce choix ne prétend pas être unique : il représente la forme la plus élémentaire compatible avec les axiomes~2 et~4, et il est suffisant pour discriminer, au niveau purement combinatoire, les architectures minimales de type \(\left(4, 6\right)\) et les structures composites candidates pour les nucléons. Un développement mathématique ultérieur pourra raffiner cette fonction d’optimisation ; toutefois, les conclusions qualitatives présentées ici ne dépendent pas de la valeur précise des poids, pourvu que ceux-ci demeurent positifs et du même ordre de grandeur.

Le principe de sélection logique s’énonce alors de manière directe : une configuration \(\sigma\) est admissible en tant que régime persistant si et seulement si elle réalise un maximum local (et, au niveau fondamental, global) de \(\Omega\left(\sigma\right)\) sous les contraintes de fermeture et d’autonomie formulées dans les sections précédentes. Les configurations qui ne satisfont pas ces contraintes, ou qui présentent une valeur de \(\Omega\) significativement plus faible, restent envisageables en tant qu’états transitoires, mais ne peuvent pas constituer des briques stables de la hiérarchie relationnelle.

Appliqué aux structures élémentaires, ce principe conduit à un résultat particulièrement restrictif. Pour des configurations strictement stationnaires, la minimisation forte de \(\nu\left(\sigma\right)\), conjointement à la réalisation d’une fermeture complète, impose une densité relationnelle maximale \(\rho\left(\sigma\right)=1\) et exclut l’existence de toute hiérarchie interne. La première configuration non exclue par cet ensemble de contraintes est précisément la structure \(\left(4, 6\right)\), dont la saturation relationnelle et la symétrie interne maximisent \(\Omega\left(\sigma\right)\) au niveau fondamental. Aucune structure comportant un nombre inférieur d’entités ou de relations ne satisfait simultanément les exigences de reproductibilité, de stabilité et d’autonomie, et toute complexification locale détruit le caractère minimal sans produire de gain corrélatif en termes de cohérence.

Pour des configurations composites, le même principe reste opérant, mais le compromis entre cohérence et connectivité présente un degré de subtilité accru. Des régimes stationnaires locaux (structurellement analogues à \(\left(4, 6\right)\)) peuvent coexister avec des régions dynamiques de plus faible densité relationnelle, à condition que l’ensemble réalise une fermeture globale et que la valeur de \(\Omega\left(\sigma\right)\) demeure maximale parmi l’ensemble des architectures admissibles. C’est dans ce cadre formel que seront examinées ultérieurement les configurations associées au proton et au neutron : elles y apparaîtront comme des optima de la fonction d’optimisation, soumis à des contraintes combinatoires supplémentaires relatives au partage entre relations stationnaires et relations dynamiques.

À partir de ce point, il devient pertinent de rechercher des grandeurs quantitatives associées à ces régimes stationnaires optimaux, non en vue d’introduire une métrique externe, mais afin d’identifier des invariants logiques susceptibles, dans un second temps, de recevoir une interprétation en termes de masse, de charge ou de fréquence propre lorsque ces structures sont effectivement réalisées physiquement. La section suivante prolonge donc la démarche en examinant la manière dont la pulsation interne d’une structure stationnaire minimale peut faire l’objet d’une lecture quantitative, sans rompre avec le cadre pré-métrique initial.

\section{Correspondance minimale avec l’électron}

Dans le cadre actuel de la physique, l’électron constitue le premier candidat empirique compatible avec le régime stationnaire logique associé à la structure \(\left(4,6\right)\). Observé comme entité élémentaire strictement stable, identique en tout point de l’univers et dépourvue de structure interne accessible expérimentalement, il satisfait les critères de stabilité et d’universalité qui découlent du cadre relationnel.

Par ailleurs, en théorie quantique des champs, l’électron est associé à une pulsation intrinsèque déterminée par sa masse au repos. Cette pulsation, généralement introduite sous la forme de la fréquence de Compton, ne dépend ni de la trajectoire de l’électron ni de son mouvement dans un référentiel particulier : elle caractérise l’entité elle-même, indépendamment de toute description géométrique spécifique.

Dans ce contexte, il ne s’agit pas de postuler que la structure \(\left(4,6\right)\) « est » un électron au sens ontologique fort, mais plutôt de soutenir qu’un régime stationnaire logique de ce type admet une première réalisation physique cohérente sous la forme de l’électron au repos. Cette identification ne requiert aucune hypothèse additionnelle sur la nature de l’électron : elle repose sur le fait que \(\left(4,6\right)\) constitue la première structure stationnaire élémentaire au sens des axiomes, et que sa pulsation logique \(s \leftrightarrow -s\) fournit un candidat naturel pour rendre compte de l’alternance de Compton associée à la masse de l’électron, via la relation \(E = \hbar \,\omega\). Dans ce cadre, \(\hbar\) quantifie la granularité minimale de l’action attachée à un cycle complet de cette pulsation logique. La correspondance proposée est minimale : elle consiste à identifier la pulsation logique interne du régime stationnaire avec la pulsation physique associée à la masse de l’électron, sans postuler de mécanisme supplémentaire.

La fréquence de Compton de l’électron au repos peut alors être interprétée comme la première manifestation quantifiée admissible du régime stationnaire logique mis en évidence au seuil de persistance. Elle ne constitue pas une preuve du cadre logique proposé, mais en fournit la première interprétation physique cohérente : une grandeur mesurable qui se laisse lire comme l’empreinte d’une périodicité interne imposée par une structure relationnelle fermée.

Dans cette perspective, la constante de Planck réduite \(\hbar\) cesse d’apparaître comme un simple paramètre empirique et peut être interprétée comme la granularité minimale de l’action associée à un régime stationnaire logique élémentaire. Elle exprime le fait qu’une pulsation interne telle que celle de l’électron ne peut pas être décomposée en portions d’action arbitrairement petites : chaque cycle de cohérence transporte une quantité minimale d’action fixée par \(\hbar\).

Plus précisément, la relation usuelle \(E = \hbar \omega\), appliquée à la pulsation propre de l’électron, devient la formulation quantitative de l’énergie de cohérence requise pour le maintien du régime stationnaire correspondant. Dans le cadre relationnel, \(\hbar\) mesure ainsi le « pas » minimal par lequel la cohérence interne se referme sur elle-même, plutôt qu’une propriété intrinsèque attachée à un champ ou à une particule préalablement donnée.

\section{Énergie comme coût de cohérence}

Dans cette perspective, l’énergie associée à la pulsation de l’électron ne doit pas être appréhendée comme un attribut fondamental surajouté à une entité préalablement constituée. Dans cette interprétation, la relation \(E = \hbar \omega\) explicite, dans le registre quantitatif, que le coût de cohérence associé à un régime stationnaire élémentaire se trouve distribué en quanta d’action de taille \(\hbar\), chaque quantum correspondant à un cycle complet de la pulsation interne. Cette relation peut ainsi être comprise comme une mesure du coût de cohérence requis pour le maintien du régime stationnaire logique correspondant.

Aussi longtemps que la structure \(\left(4,6\right)\) demeure préservée, la distinction qu’elle institue ne peut être abolie sans générer une contradiction logique. Le maintien de cette cohérence interne impose une contrainte sur l’architecture relationnelle globale ; lorsque cette contrainte est effectivement réalisée dans un substrat physique, elle se manifeste sous la forme d’une énergie de repos bien définie.

À ce stade, aucune métrique spatiale n’est requise, aucune temporalité fondamentale n’est postulée et aucune hiérarchie interne supplémentaire n’est introduite. L’énergie est définie comme la traduction quantitative d’un régime de cohérence, plutôt que comme une grandeur primitive indépendante de la structure relationnelle qui la sous-tend.

À partir du moment où la pulsation stationnaire reçoit une interprétation quantitative, l’énergie ne peut plus être considérée comme une grandeur strictement globale. La possibilité même de l’existence de plusieurs régimes stationnaires implique que le coût de cohérence associé à chacun d’eux soit distinguable : chaque régime stationnaire élémentaire doit pouvoir être identifié par une valeur propre, sans que cette identification dépende d’un arrière-plan métrique préalablement constitué.

\section{Possibilité de localisation logique de l’énergie}

Dès lors que l’énergie est interprétée comme le coût de cohérence requis pour le maintien d’un régime stationnaire, elle ne peut plus être appréhendée comme une grandeur strictement globale, inséparable d’un arrière-plan uniforme et indifférencié. L’existence de plusieurs régimes stationnaires distincts impose, en effet, que ce coût puisse être différencié et attribué spécifiquement à chacun de ces régimes.

Le maintien d’un régime stationnaire déterminé suppose l’imposition d’une contrainte associée à une organisation particulière de la structure relationnelle sous-jacente. Une telle contrainte ne peut être dissoute dans un cadre homogène et non structuré sans perdre toute signification fonctionnelle : elle est, par nécessité, corrélée à la structure qui porte et réalise le régime stationnaire. Dans cette perspective, l’énergie devient logiquement localisable : elle se trouve associée à une organisation déterminée, bien qu’elle ne soit pas encore localisée au sens géométrique du terme.

À ce stade, la localisation considérée demeure purement conceptuelle. Elle n’implique ni position spatiale, ni extension métrique, ni notion de distance. Elle exprime uniquement le fait que la cohérence du système n’est pas indéfiniment distribuable : pour qu’un régime stationnaire soit maintenu, une fraction finie du « budget de cohérence » du cadre relationnel doit être spécifiquement allouée à une structure particulière.

\section{Coexistence de plusieurs régimes stationnaires}

La possibilité d’associer une énergie propre à un régime stationnaire ouvre la voie à la coexistence de plusieurs régimes distincts au sein d’un même cadre relationnel. Toutefois, cette coexistence n’est pas triviale : chaque régime impose des contraintes internes de cohérence, et ces contraintes doivent demeurer mutuellement compatibles pour que plusieurs régimes puissent persister simultanément.

La coexistence de plusieurs régimes stationnaires introduit ainsi une exigence supplémentaire : chaque régime doit être en mesure de préserver sa cohérence intrinsèque en présence d’autres régimes également persistants. Cette exigence ne relève pas encore d’une interaction au sens physique strict ; elle n’implique ni échange d’énergie ou d’information, ni force, ni mécanisme de propagation. Elle exprime plutôt une contrainte de compatibilité logique : deux régimes stationnaires ne peuvent coexister de manière stable que si la dynamique interne de la distinction qui leur est propre n’entre pas en conflit avec celle de l’autre.

La persistance d’un régime stationnaire ne dépend donc plus uniquement de sa pulsation interne ou du coût de cohérence qui lui est associé. Elle dépend également de la capacité de sa structure interne à se maintenir sans compromettre la cohérence des régimes voisins. Deux régimes caractérisés par des pulsations et des énergies comparables peuvent néanmoins se révéler incompatibles si leurs organisations internes ne permettent pas une coexistence stable.

Il s’ensuit que la stabilité interne acquiert un caractère structurel irréductible, qui ne peut être ramené à une unique grandeur scalaire. La compatibilité ou l’incompatibilité entre régimes stationnaires devient alors une propriété de leurs organisations internes, et non une simple fonction de leurs valeurs d’énergie.

\section{Vers une propriété différentielle : la charge}

La coexistence de plusieurs régimes stationnaires au sein d’un même cadre relationnel soulève une problématique inédite : comment des structures fermées peuvent-elles conserver leur cohérence interne lorsqu’elles cohabitent, sans que leurs régimes respectifs ne se neutralisent mutuellement ? Un électron, un proton ou un neutron réalisent chacun une organisation relationnelle stable ; toutefois, leur compatibilité mutuelle n’est pas triviale dès lors qu’aucun arrière-plan métrique n’est disponible pour les « séparer » ou les « isoler ».

Dans le cadre axiomatique introduit précédemment, cette compatibilité se formule en termes de cohérence triangulaire (axiome~2) et d’optimisation logique (axiome~4). Deux régimes stationnaires sont dits compatibles si la superposition de leurs structures n’engendre pas de violations supplémentaires de la cohérence triangulaire qui ne puissent être compensées par une fermeture globale. En d’autres termes, la composition de deux structures fermées doit demeurer admissible au sens de la fonction d’optimisation : le nombre de violations \(\nu\left(\sigma\right)\) ne peut augmenter que dans des proportions telles que la valeur de \(\Omega\left(\sigma\right)\) reste maximale parmi l’ensemble des configurations accessibles.

Dans une perspective strictement relationnelle, la « charge » apparaît alors comme la signature interne de cette compatibilité différenciée. La pulsation logique associée à un régime stationnaire ne se réduit pas au simple fait « qu’il se passe quelque chose plutôt que rien » ; elle possède une structure fine, déterminée par la manière dont les accords et les désaccords se distribuent le long des cycles constitutifs de la structure fermée. Deux régimes peuvent ainsi présenter une stabilité globale équivalente, tout en différant quant à la façon dont leurs alternances internes s’articulent lorsqu’ils sont mis en relation.

On peut définir, de façon abstraite, la charge \(q\) d’un régime stationnaire comme une asymétrie interne entre les contributions constructives et non constructives de ses relations. Les relations qui renforcent la cohérence triangulaire jouent ici le rôle de contributions « en phase », tandis que les relations qui tendent à introduire des ruptures non compensées jouent le rôle de contributions « en opposition ». La charge mesure alors, dans un sens purement logique, l’excès relatif de l’un de ces deux types de contributions lorsqu’un régime stationnaire est mis en composition avec un autre.

Dans ce cadre, il est possible de sélectionner un régime de référence dont la charge est normalisée à \(-1\). La structure minimale \(\left(4,6\right)\), identifiée dans la suite du texte au régime stationnaire de l’électron, constitue alors un candidat naturel : sa clôture complète et sa symétrie interne en font une unité de comparaison privilégiée. Un régime stationnaire composite dont la pulsation interne compense exactement l’asymétrie de ce régime de référence réalisera une charge \(+1\), tandis qu’un régime dont l’organisation interne demeure neutre vis-à-vis de cette asymétrie sera associé à une charge nulle.

Cette définition n’érige pas la charge en propriété adjointe a posteriori aux structures. Elle la fait au contraire émerger comme la conséquence de deux contraintes fondamentales : d’une part, la nécessité de préserver la cohérence triangulaire lors de la composition de régimes stationnaires (axiome~2) ; d’autre part, le principe de sélection, par optimisation, des configurations qui satisfont cette contrainte avec un coût minimal au sens de \(\Omega\left(\sigma\right)\) (axiome~4). Une structure strictement symétrique, ne présentant aucune asymétrie interne exploitable, se trouverait incapable d’interagir de manière différenciée et resterait, de ce fait, neutre relativement à cette propriété.

Dans les sections suivantes, cette notion de charge sera appliquée à des architectures plus concrètes. La distribution asymétrique des relations au sein des sous-structures internes d’un proton, combinée à l’existence d’une mer relationnelle dynamique, fournira une réalisation explicite de la charge \(+1\) comme signature d’un déséquilibre minimal, requis pour assurer la compatibilité avec la structure \(\left(4,6\right)\) associée à l’électron. La charge apparaît ainsi comme le coût logique nécessaire pour qu’une structure composite puisse interagir de manière stable avec une structure élémentaire, sans compromis sur la cohérence globale du cadre relationnel.

\section{Rôle de l’observation dans la construction du cadre}

L’analyse a, jusqu’à présent, été délibérément conduite dans un registre strictement logique, en suspendant toute référence explicite aux observations physiques. Néanmoins, la structure même des contraintes imposées au cadre conceptuel n’est pas indépendante de l’expérience : elle est, au contraire, informée et orientée par un ensemble de faits empiriques robustes.

En particulier, l’exigence d’un régime stationnaire minimal reflète l’existence, au sein de la physique établie, d’entités strictement stables, telles que l’électron, dont les propriétés demeurent invariantes dans l’ensemble des contextes d’observation accessibles. De même, la nécessité de distinguer des régimes stationnaires coexistant sans se confondre formalise le constat empirique selon lequel plusieurs particules peuvent être présentes simultanément sans perdre leur identité propre ni leur caractère distinct.

L’observation ne joue donc pas le rôle d’une source directe de postulats, mais celui d’un filtre rétrospectif : parmi l’ensemble des cadres logiquement possibles, ne sont retenus que ceux qui admettent une correspondance minimale avec des structures effectivement mises en évidence expérimentalement. C’est parce que la physique expérimentale manifeste l’existence d’entités stables, localisables et différenciables que la construction logique se trouve contrainte de faire émerger des régimes stationnaires, des énergies localisables et des charges.

En ce sens, le rôle de l’observation est fondamentalement double. D’une part, elle rend le travail de formalisation possible en révélant les invariants phénoménologiques que le cadre logique doit être en mesure de reconstruire. D’autre part, elle fournit le principe de sélection entre les constructions logiques envisageables : seules celles qui permettent une mise en correspondance cohérente et systématique avec le domaine empirique observé peuvent être retenues comme candidates plausibles pour décrire les modalités d’émergence de la réalité physique.

\section{Nécessité d’une séparation pré-métrique}

L’introduction de la charge en tant que propriété différentielle interne permet de caractériser, au sein d’un même cadre relationnel, la compatibilité ou l’incompatibilité de régimes stationnaires coexistant. Toutefois, tant que ce cadre demeure entièrement indifférencié, rien ne s’oppose, du point de vue strictement logique, à ce que des régimes mutuellement incompatibles se superposent sans qu’aucune distinction formelle ne puisse être établie.

Une telle situation entrerait en conflit avec l’exigence de persistance : si deux régimes stationnaires incompatibles pouvaient effectivement se confondre, la cohérence interne de chacun serait compromise, et aucune structure stable ne pourrait être maintenue. La coexistence de régimes stationnaires porteurs de charges distinctes requiert donc l’existence d’une forme minimale de séparation entre configurations incompatibles.

Cette séparation ne peut pas encore être interprétée en termes d’éloignement spatial, puisqu’aucune structure métrique n’a été définie. Elle ne peut pas davantage être conçue comme l’effet d’une interaction ou d’une dynamique, en l’absence de toute notion de force ou de propagation. Il s’agit d’une séparation pré-métrique, définie uniquement par des contraintes de cohérence : certaines combinaisons de régimes sont exclues, non parce qu’elles « se repoussent », mais parce qu’elles ne peuvent pas coexister au sein d’une même organisation relationnelle sans engendrer de contradiction.

\section{Vers une métrique minimale}

Dès lors que plusieurs régimes stationnaires coexistent, différenciés par leur charge et séparés par des contraintes de cohérence, il devient nécessaire de disposer d’un dispositif formel permettant de représenter ces discontinuités sans en altérer le statut logique.

Un tel dispositif ne génère pas la séparation ; il en propose une représentation explicite. Il ne constitue pas une structure physique additionnelle, mais un outil formel minimal visant à rendre comparables des relations déjà déterminées par les contraintes de compatibilité. Il s’agit plus précisément d’introduire une métrique minimale, qui n’instaure pas encore un espace géométrique continu, mais qui permet de comparer et d’ordonner les séparations entre différents régimes stationnaires.

Cette métrique minimale n’introduit ni géométrie au sens classique, ni continuité, ni dynamique. Elle se limite à discriminer les situations dans lesquelles deux régimes stationnaires peuvent partager un même voisinage logique de celles où ils doivent être formellement disjoints. Elle fournit un langage formel pour exprimer qu’une configuration est « plus contrainte » qu’une autre, sans pour autant associer de distances ou de durées au sens physique du terme.

Il est crucial de souligner que cette métrique minimale ne possède aucun statut ontologique fondamental. Elle n’est pas postulée comme une structure a priori du réel, mais introduite comme une réponse nécessaire à un problème de représentation : en son absence, il serait impossible de décrire de manière structurée un ensemble de régimes stationnaires différenciés par leurs charges et leurs relations de compatibilité.

\section{Structures composites et hiérarchie minimale}

La métrique minimale constitue un cadre de représentation permettant de décrire la coexistence de plusieurs régimes stationnaires, distincts par leur charge et séparés par des contraintes de cohérence. Néanmoins, l’observation physique met en évidence des entités dont la masse au repos excède largement celle de l’électron, tout en présentant une stabilité de même ordre de grandeur. De telles entités ne peuvent être assimilées à de simples agrégats d’électrons indépendants : elles manifestent une organisation interne structurée, caractérisée par une cohérence globale irréductible à la somme des contributions élémentaires.

Une cohérence de cette nature ne peut être rendue compte par une superposition naïve de régimes stationnaires élémentaires, chacun conservant une clôture propre intégrale. Si chaque sous-structure réalisait une fermeture maximale du type \(\left(4,6\right)\), la densité relationnelle globale atteindrait un seuil incompatible avec l’émergence d’une hiérarchie interne : la structure composite deviendrait alors logiquement et fonctionnellement indécomposable.

L’existence d’une organisation composite stable requiert par conséquent une architecture dans laquelle la saturation relationnelle est non uniforme : certaines régions atteignent une fermeture maximale, tandis que d’autres restent en état de connectivité partielle. Cette distribution hétérogène de la densité relationnelle définit ce que l’on désigne comme la hiérarchie minimale d’une structure composite.

Dans une structure minimale telle que la configuration \(\left(4,6\right)\), la fermeture est dite parfaite : l’ensemble des relations nécessaires à la cohérence interne est entièrement contenu dans la structure, et aucune composante de la circulation de cohérence ne demeure sans support. Une telle configuration ne présente aucune « fuite logique » ; elle réalise une saturation relationnelle maximale compatible avec la stabilité du régime stationnaire.

À l’inverse, une structure composite hiérarchisée se trouve, par construction, dans une situation différente. Pour éviter une saturation globale, elle doit maintenir certaines régions en régime stationnaire local, tandis que d’autres demeurent en régime dynamique. Cette dissymétrie interne interdit la réalisation d’une fermeture parfaite à l’échelle globale : une fraction de la circulation de cohérence ne peut pas être intégralement confinée dans les sous-structures closes et doit nécessairement transiter par l’environnement relationnel collectif.

On peut, dans ce contexte, introduire la notion de \emph{fuite de cohérence} : une fraction minimale de la circulation nécessaire au maintien de la structure ne se referme pas totalement à l’intérieur de la hiérarchie interne, mais se propage à travers le cadre relationnel qui supporte la structure. Cette fuite ne doit pas être interprétée comme une perte au sens énergétique ; elle exprime plutôt l’impossibilité logique, pour une structure hiérarchique composite, de réaliser simultanément une fermeture maximale locale et une fermeture maximale globale.

\section{Le proton comme structure composite minimale : une architecture hiérarchique à deux niveaux}

Le proton ne peut être modélisé comme une simple juxtaposition de trois occurrences indépendantes de la structure minimale \(\left(4,6\right)\). Il doit plutôt être décrit comme une architecture hiérarchique à deux niveaux, constituée (i) de régimes stationnaires localisés, assimilables aux quarks de valence, et (ii) d’une mer relationnelle dynamique assurant leur couplage.

Les sous-structures stationnaires agissent comme des noyaux de cohérence locale : chacune tend à maximiser sa fermeture interne, de manière analogue à la structure \(\left(4,6\right)\), mais à une échelle composite plus complexe et plus riche. À l’inverse, la mer relationnelle se caractérise par une connectivité minimale : les relations y demeurent non saturées, de sorte qu’elles n’optimisent pas la fermeture interne mais garantissent la circulation et la propagation de la cohérence entre les quarks de valence.

Cette dissymétrie émerge des contraintes imposées par les axiomes de cohérence et d’optimisation. En effet, si l’ensemble des sous-structures était simultanément saturé selon le schéma \(\left(4,6\right)\), la densité relationnelle globale deviendrait telle qu’aucun transfert de cohérence entre ces sous-structures ne serait possible. L’édifice protonique requiert par conséquent une hiérarchie minimale : une fermeture maximale à l’échelle locale, une fermeture seulement partielle à l’échelle globale, et une mer relationnelle jouant le rôle de médiateur assurant la cohérence du système dans son ensemble.

\subsection{Contraintes combinatoires sur les quarks de valence}

Les régimes stationnaires locaux doivent satisfaire la condition de complétude minimale imposée par l’axiome~3 : le nombre d’entités élémentaires constituant un quark doit être un multiple de \(4\), de manière à permettre la construction de sa structure interne comme une superposition de tétrades logiques. On obtient ainsi, pour un quark générique :
\[
n_{\mathrm{quark}} \in \{4,8,12,16,20,24,28,\ldots\}.
\]

Dans le cas d’une structure stationnaire fermée, chaque quark est représenté par un graphe complet \(K_n\), caractérisé par :
\[
r_{\mathrm{quark}}
= \frac{n_{\mathrm{quark}}\left(n_{\mathrm{quark}}-1\right)}{2},
\]
où \(r_{\mathrm{quark}}\) désigne le nombre de relations internes (arêtes) au sein du graphe.

L’exigence de complétude interne est par ailleurs combinée à une contrainte d’asymétrie minimale : afin de permettre l’émergence d’une charge nette, les trois quarks d’un triplet ne peuvent pas présenter des architectures strictement identiques. Deux d’entre eux doivent partager la même structure interne (quarks up), tandis que le troisième (quark down) doit réaliser une structure légèrement distincte.

Cette brisure contrôlée de la symétrie interne constitue la condition nécessaire pour que la configuration globale maximise la fonction d’optimisation \(\Omega\left(\sigma\right)\), tout en générant une charge \(+1\) compatible avec la charge \(-1\) de la structure \(\left(4,6\right)\) associée à l’électron.

\subsection{Justification du choix \(n_u = 24\) et \(n_d = 28\)}

Sous ces contraintes, on considère un ensemble de configurations candidates pour lesquelles les quarks de valence vérifient :
\[
n_u, n_d \in \{4,8,12,16,20,24,28,\ldots\},
\]
avec deux quarks up identiques et un quark down distinct. Pour chaque paire \(\left(n_u,n_d\right)\), la structure composite associée est évaluée à l’aide de la fonction d’optimisation \(\Omega\left(\sigma\right)\) introduite à la section~10, en tenant compte simultanément : du nombre de violations de cohérence triangulaire \(\nu\left(\sigma\right)\), de la densité relationnelle, ainsi que du critère de minimalité hiérarchique.

À titre d’exemple, les choix \(\left(n_u,n_d\right)=\left(12,16\right)\) ou \(\left(16,20\right)\) conduisent effectivement à des quarks de valence obtenus par superposition de tétrades logiques (3 et 4 tétrades respectivement), mais ces configurations se révèlent insatisfaisantes sur plusieurs aspects :
\begin{itemize}
  \item la densité relationnelle globale demeure trop proche du régime de saturation, ce qui inhibe la dynamique de la mer relationnelle : la valeur de \(\rho\left(\sigma\right)\) est excessive et la contribution \(f_{\mathrm{conn}}\left(\rho\right)\) ne compense pas l’augmentation concomitante de \(\nu\left(\sigma\right)\) ;
  \item la distribution des charges effectives ne permet pas de réaliser une configuration \(u,u,d\) de charge nette \(+1\) sans introduire des asymétries marquées, ce qui accroît \(\nu\left(\sigma\right)\) au-delà du seuil compatible avec une structure stable.
\end{itemize}

À l’inverse, des paires plus élevées telles que \(\left(n_u,n_d\right)=\left(28,32\right)\) ou au-delà offrent une marge suffisante pour l’existence d’une mer relationnelle substantielle, mais au prix d’un nombre total de relations \(r_{\mathrm{total}}\) nettement supérieur à \(1{,}1\cdot 10^{4}\). Ces architectures sont alors exclues par le critère de minimalité, pour lequel \(\delta_{\text{minimalité}}\!\left(\sigma\right)=0\), de sorte que la valeur de \(\Omega\left(\sigma\right)\) est nettement inférieure à celle obtenue pour \(\left(24,28\right)\).

Dans ce cadre, la paire \(\left(n_u,n_d\right)=\left(24,28\right)\) apparaît comme la première combinaison maximisant \(\Omega\left(\sigma\right)\) sous les contraintes combinatoires imposées : multiplicités multiples de \(4\), asymétrie minimale nécessaire à la production d’une charge nette \(+1\), présence d’une mer relationnelle significative, et fraction stationnaire d’ordre \(\sim 9\) compatible avec la part d’énergie de valence dans la masse du proton.

Au terme de ce balayage combinatoire, la paire \(\left(n_u=24,n_d=28\right)\) constitue ainsi la première configuration satisfaisant simultanément : (i) une complétude interne de chaque quark (multiplicités multiples de \(4\)), (ii) une asymétrie minimale suffisante pour engendrer une charge totale \(+1\), et (iii) une valeur maximale de \(\Omega\left(\sigma\right)\) parmi l’ensemble des configurations testées. Les valeurs
\[
n_u = 24 = 6\times 4,\qquad
n_d = 28 = 7\times 4
\]
doivent alors être interprétées comme le résultat d’un processus de sélection sous contraintes plutôt que comme des paramètres arbitraires : elles correspondent à des superpositions de six et sept tétrades logiques, respectivement, qui optimisent la cohérence interne de la structure protonique.

De ce fait, les nombres de relations stationnaires associés aux quarks up et down s’écrivent :
\[
r_u = \frac{24\times 23}{2} = 276,\qquad
r_d = \frac{28\times 27}{2} = 378,
\]
ce qui prépare l’analyse détaillée des contributions de valence et de la mer relationnelle développée dans les sous-sections suivantes.

\subsection{Quantification de la mer relationnelle}

La contribution de la mer de quarks est modélisée par un graphe \emph{sparse} de connectivité minimale, caractérisé par la règle de comptage suivante :
\[
R_{\mathrm{mer}} = 2\, n_{\mathrm{mer}}.
\]
Pour le proton, l’application de cette règle conduit à :
\[
R_{\mathrm{mer}} \simeq 10\,087.
\]
Cette contribution est interprétée comme la composante dynamique (non stationnaire) de la cohérence relationnelle globale.

\subsection{Total relationnel et décomposition stationnaire/dynamique}

Le nombre total de relations internes décrivant l’architecture relationnelle du proton est donné par :
\[
r_{\mathrm{total}} = r_{\mathrm{valence}} + R_{\mathrm{mer}} = 930 + 10\,087 = 11\,017.
\]
On en déduit immédiatement les fractions relatives suivantes :
\[
\text{Fraction stationnaire : } f_{\mathrm{stat}} = \frac{930}{11\,017} \simeq 9\%\!,
\qquad
\text{Fraction dynamique : } f_{\mathrm{dyn}} = \frac{10\,087}{11\,017} \simeq 91\%\!.
\]

Ces valeurs sont qualitativement cohérentes avec le constat empirique selon lequel \(\sim 91\%\) de l’énergie du proton proviennent de la dynamique interne (champ de gluons et mer de quarks), et non des seules masses des quarks de valence. 

Dans le cadre LDP, le partage \(\approx \frac{9}{91}\) émerge d’une architecture hiérarchique déterminée par les axiomes~2–4 (minimisation des violations, maximisation de la connectivité, contrainte de minimalité) et par une fonction d’optimisation dont les poids sont choisis pour refléter la dominance des régimes dynamiques, plutôt que par un ajustement direct à une valeur expérimentale spécifique.

La bonne compatibilité avec les estimations issues de la QCD peut dès lors être interprétée comme une validation \emph{a posteriori} de cette architecture dérivée du cadre LDP, plutôt que comme la conséquence d’un \emph{fit} paramétrique imposé aux données. 

Les entiers considérés ne sont pas introduits comme des paramètres libres, mais sont sélectionnés sous une double contrainte :
\begin{itemize}
  \item compatibilité avec les axiomes (multiplicité de \(4\), utilisation de graphes complets ou \emph{sparse} minimaux, minimisation de \(\nu\)) ;
  \item reproduction de rapports sans dimension issus de la QCD (partage \(\approx \frac{9}{91}\), ordres de grandeur des rayons et des masses réduites).
\end{itemize}

Des simulations effectuées pour d’autres choix d’architectures relationnelles indiquent qu’elles requièrent soit des tailles de graphes considérablement plus grandes, soit une accumulation de violations incompatible avec l’axiome~4, ce qui renforce le caractère non arbitraire de l’architecture retenue.

Dans le formalisme graphique introduit précédemment, ces structures sont représentées par des graphes signés de grande taille, au sein desquels trois fermetures stationnaires locales jouent le rôle effectif de quarks de valence. Les entiers \(930\), \(10\,087\) et \(11\,017\) dénombrent respectivement les relations stationnaires associées à ces fermetures, les relations dynamiques de la « mer relationnelle » et le total des relations internes. Leur sélection résulte d’une optimisation conjointe de la cohérence triangulaire (minimisation de \(\nu\)), de la densité relationnelle \(\rho\) et de la minimalité hiérarchique, sous les contraintes imposées par les axiomes~2 à~4.

\subsection{Fuite de cohérence}

La hiérarchie ne peut pas satisfaire simultanément la condition de fermeture maximale pour chacun de ses sous-systèmes et pour la structure globale prise dans son ensemble. Il en résulte une fuite de cohérence minimale : une fraction extrêmement réduite de la circulation relationnelle nécessaire au maintien de la stabilité du proton ne se referme pas entièrement dans ses \(11\,017\) relations internes, mais se propage au-delà, dans le cadre relationnel environnant.

On introduit alors un paramètre de fuite \(\varepsilon\), qui quantifie la déviation relative entre une configuration de fermeture hiérarchique idéale et la fermeture effectivement réalisée par la structure protonique. Ce paramètre, très faible pour une structure composite minimale telle que le proton, caractérise une violation logique minimale de la fermeture parfaite, suffisante pour instaurer un couplage entre la hiérarchie interne et le cadre relationnel externe.

Dans le langage des axiomes, \(\varepsilon\) peut être interprété comme la fraction résiduelle de cohérence qui ne peut être capturée par une minimisation stricte de \(\nu\left(\sigma\right)\) restreinte aux seuls graphes internes au proton : il mesure l’écart entre l’optimum associé à une fermeture globale idéale (dans laquelle tous les triangles, à tous les niveaux hiérarchiques, seraient simultanément cohérents) et l’optimum effectivement atteignable sous les contraintes de minimalité et de hiérarchie imposées par \(\Omega\left(\sigma\right)\).

Dans ce qui suit, ce paramètre jouera un rôle central : il permettra d’établir une relation quantitative entre la structure interne du proton et les effets de courbure de la métrique émergente, et de définir une constante de couplage gravitationnel effective.

\section{Surface active du proton et couplage avec l’électron}

La description relationnelle du proton distingue deux contributions principales : \(930\) relations stationnaires associées aux trois quarks de valence, et une mer de quarks dynamique totalisant environ \(10\,087\) relations. Toutes ces relations ne jouent cependant pas un rôle équivalent vis-à-vis d’un électron : lors du couplage, seule une fraction de la structure protonique intervient effectivement dans l’interaction.

Pour rendre compte de ce phénomène, on introduit la notion de \emph{surface active} du proton. On définit comme surface active l’enveloppe minimale de relations qui médiatisent effectivement le couplage logique avec un électron modélisé par la structure \(\left(4,6\right)\). Cette surface ne se réduit ni à l’ensemble complet des relations internes du proton, ni à la seule contribution des relations de valence ; elle correspond au sous-ensemble de relations qui portent la signature de charge \(+1\) et qui sont accessibles au régime stationnaire caractérisant l’électron.

Dans le formalisme des graphes signés, ce couplage est assuré par les triangles mixtes impliquant simultanément des sommets de la structure protonique et de la structure \(\left(4,6\right)\) associée à l’électron. Les relations appartenant à la surface active sont précisément celles qui interviennent dans des triangles mixtes cohérents, c’est-à-dire qui, une fois combinées avec les signes internes de \(\left(4,6\right)\), ne génèrent pas de violation supplémentaire de la cohérence triangulaire. La charge \(+1\) effective du proton vis-à-vis de l’électron se traduit alors par une asymétrie globale en faveur de ces contributions constructives au sein de l’ensemble des triangles mixtes accessibles.

Pour quantifier cette surface active, il est pertinent de distinguer la contribution des deux quarks up (\(u\)) et de celle du quark down (\(d\)). Les quarks up, porteurs d’une charge électrique \(+2/3\), contribuent chacun à la cohérence de surface par l’intermédiaire de leurs \(276\) relations stationnaires, tandis que le quark down, de charge \(-1/3\), contribue au moyen de ses \(378\) relations. Une estimation naturelle du « contenu de charge de surface » consiste alors à pondérer ces contributions par leurs charges respectives, ce qui conduit à considérer une combinaison du type :
\[
R_{\mathrm{brut}} \sim 2\, r_u + r_d
= 2\cdot 276 + 378
= 930,
\]
laquelle coïncide numériquement avec le nombre total de relations de valence.

Dans le formalisme des fermetures, les blocs internes \(u\) et \(d\) ne jouent pas le même rôle que la brique élémentaire \(\left(4,6\right)\). Ils ne constituent pas des structures stationnaires élémentaires autonomes, mais doivent être interprétés comme des fermetures stationnaires locales : ce sont des sous-graphes densément connectés, de taille \(24\) pour \(u\) et \(28\) pour \(d\), au sein desquels la cohérence triangulaire est fortement maximisée et où une pulsation interne compatible avec les axiomes peut être réalisée. Leur cohérence globale ne peut toutefois être garantie que dans le cadre de l’architecture hiérarchique complète du proton, qui inclut la mer relationnelle ainsi que les contraintes de fermeture globale.

Les blocs internes \(u\) (24 entités, 276 relations stationnaires) et \(d\) (28 entités, 378 relations stationnaires) sont ainsi formellement décrits comme des fermetures stationnaires locales au sein du proton : des domaines densément connectés où la cohérence triangulaire est fortement maximisée et où une pulsation interne peut être réalisée, mais qui ne satisfont pas, pris isolément, l’ensemble des conditions de complétude minimale définies pour la brique \(\left(4,6\right)\). Leur stabilité logique est donc conditionnée par l’organisation hiérarchique globale du proton (valence, mer, surface, fuite).

Cependant, l’ensemble des relations de valence ainsi définies n’est pas également efficace pour le couplage avec l’extérieur : une fraction d’entre elles est principalement mobilisée pour assurer le maintien de la cohérence interne des quarks et ne se projette pas directement sur l’interface avec l’électron. En tenant compte de cette hiérarchie interne — incluant les quarks localement fermés, la mer de quarks et les contraintes de fermeture globale — on obtient une surface effective réduite, de l’ordre de :
\[
R_{\mathrm{surface}} \simeq 500
\]
relations actives. Ce nombre effectif caractérise, en termes combinatoires, la taille du sous-ensemble de relations contribuant effectivement à la formation de triangles mixtes compatibles avec la structure \(\left(4,6\right)\). Il est inférieur au contenu brut de valence \(\left(930\right)\), dans la mesure où une partie des relations stationnaires des quarks \(u\) et \(d\) demeure principalement engagée dans la stabilisation de la cohérence interne des blocs et ne se projette pas directement sur l’interface logique avec l’électron.

Cette distinction entre nombre total de relations et nombre de relations de surface joue un rôle conceptuel central dans la suite. La constante de structure fine \(\alpha\) pourra être réinterprétée comme un indicateur de l’efficacité relative de ce couplage, en comparant la surface active du proton, le régime stationnaire de l’électron et le rapport de masse \(\mu = m_p/m_e\). La section suivante développera cette perspective afin de proposer une relecture relationnelle de \(\alpha\), en montrant comment l’ordre de grandeur de \(\alpha^{-1} \simeq 137\) peut émerger de l’architecture interne du proton et du caractère minimal de l’électron.

\section{Longueur d’onde de Compton du proton et fermeture de la forme}

La structure relationnelle du proton a, jusqu’à présent, été décrite sans recours explicite à une métrique spatiale préalable. Néanmoins, le proton est caractérisé expérimentalement comme une entité dotée d’un rayon effectif et d’une longueur d’onde de Compton \(\lambda_p\), associée à sa masse au repos. Dans le cadre présent, cette longueur d’onde n’est pas introduite comme une grandeur indépendante, mais comme la première traduction métrique d’un régime stationnaire global corrélé à la fermeture de la forme.

Soit \(R_p\) un rayon logique effectif du proton. Les mesures de diffusion conduisent typiquement à \(R_p \simeq 0{,}84\times 10^{-15}\,\mathrm{m}\). Une onde stationnaire se propageant sur la surface peut être modélisée, au premier ordre, comme parcourant un grand cercle. En supposant qu’une onde unique boucle la forme en un demi-tour, la longueur d’arc parcourue est \(\pi R_p\). La longueur d’onde de Compton de référence pour le proton est \(\lambda_p \simeq 1{,}32\times 10^{-15}\,\mathrm{m}\).

Le rapport \(\lambda_p/(\pi R_p)\) est alors voisin de \(0{,}5\), ce qui suggère qu’il ne circule pas une, mais deux ondes de bouclage sur la surface logique du proton, chacune effectuant un demi-tour au cours d’une période. Le choix d’un mode suivant un grand cercle doit être compris comme un modèle minimal de mode de surface, parmi d’autres configurations géométriquement admissibles. Il ne prétend pas rendre compte de l’ensemble des solutions compatibles avec une forme fermée, mais fournit un scénario simple permettant d’examiner la cohérence entre la fermeture logique et sa première traduction métrique :
\[
2\pi R_p \simeq 2\times 1{,}32\times 10^{-15}\,\mathrm{m},
\]
ce qui, pour \(R_p \simeq 0{,}84\times 10^{-15}\,\mathrm{m}\), conduit à un accord numérique meilleur que \(1\,\%\). L’écart relatif entre le rayon déduit de \(\lambda_p\) par la relation \(\lambda_p \simeq 2\pi R_p\) et le rayon mesuré peut être interprété comme une légère « respiration » de la forme, liée à la dynamique interne de la mer de quarks.

Dans cette interprétation, la « longueur d’onde de Compton » n’introduit pas une nouvelle grandeur indépendante : elle reprend la relation usuelle \(\lambda = h/(mc)\) (ou \(\lambda = \hbar/(mc)\) selon la convention adoptée) appliquée à la masse au repos du proton, et sert uniquement de pont entre la fréquence interne et une échelle métrique effective. L’accord numérique à mieux que \(1\,\%\) doit donc être compris comme un test de plausibilité du modèle de surface proposé, plutôt que comme une dérivation unique et définitive de la géométrie protonique.

Cette cohérence numérique ne constitue pas une preuve au sens strict, mais fournit un indicateur selon lequel la longueur d’onde de Compton du proton peut être interprétée comme la manifestation géométrique d’une onde stationnaire de surface, contrainte par le régime stationnaire global. Dans ce cadre, l’espace ne peut plus être considéré comme un simple contenant neutre ; il apparaît plutôt comme la modalité selon laquelle une structure de fermeture logique stable se projette dans une représentation métrique.

\section{Temps propre interprété comme comptage de cycles}

L’interprétation de la longueur d’onde de Compton du proton comme la manifestation d’une onde stationnaire de surface permet d’attribuer un statut physique analogue à la fréquence de Compton associée à sa masse au repos. Cette fréquence ne décrit pas un mouvement se déroulant dans un temps préalablement donné ; elle quantifie plutôt le nombre de cycles de pulsation nécessaires au maintien d’un régime stationnaire global.

\subsection{Temps propre et densité relationnelle}

Dans un cadre relationnel, le temps propre associé à une structure n’est pas conçu comme une coordonnée externe, mais comme une mesure du nombre de cycles de cohérence que cette structure doit accomplir pour préserver son intégrité organisationnelle. Afin d’articuler cette notion avec une description plus géométrique, il est utile d’introduire la densité relationnelle \(\rho\), définie comme le nombre de relations actives par unité de volume logique.

On peut alors écrire, de manière schématique, que l’élément de temps propre \(\mathrm{d}\tau\) associé à une région caractérisée par une densité relationnelle \(\rho\) est relié à un pas de pulsation de référence \(\mathrm{d}t_0\) par une relation de la forme :
\[
\mathrm{d}\tau \sim \sqrt{\rho}\,\mathrm{d}t_0.
\]

Dans une région où la densité relationnelle est élevée (milieu matériel dense, forte concentration de structures hiérarchiquement organisées), la valeur effective de \(\mathrm{d}\tau\) par unité de \(\mathrm{d}t_0\) est réduite : la structure doit consacrer une fraction plus importante de ses cycles internes au maintien de sa cohérence, ce qui se manifeste, pour un observateur externe, sous la forme d’une dilatation du temps propre. À l’inverse, dans une région de faible densité relationnelle (proche du « vide logique »), la structure mobilise un nombre moindre de cycles pour conserver sa cohérence, et le temps propre s’écoule alors plus rapidement relativement au même pas de référence \(\mathrm{d}t_0\).

Cette dépendance qualitative reproduit la tendance générale des effets de dilatation temporelle bien établis : le temps propre est ralenti dans les régions de forte densité de matière ou d’énergie, et accéléré dans les régions relativement « vides ». La différence avec la relativité générale réside dans le fait que la courbure de la métrique émergente est ici interprétée comme la trace, dans une description effective, de la distribution de la densité relationnelle \(\rho\) et du coût de cohérence qui lui est associé, plutôt que comme une structure fondamentale autonome.

Dans le cas de l’électron, la fréquence de Compton \(\nu_e\) fournit une cadence intrinsèque extrêmement élevée : chaque période correspond à une ré-affirmation complète de la cohérence interne de la structure minimale \(\left(4,6\right)\). Le temps propre associé à l’électron peut alors être défini comme le comptage de ces cycles, c’est-à-dire le nombre d’itérations par lesquelles la structure confirme sa stabilité dynamique.

Cette définition ne modifie pas la notion de temps propre telle qu’elle est employée en relativité restreinte et générale : elle en propose une reformulation en termes de cadence interne, en mobilisant la relation \(E = \hbar \omega\) appliquée à la masse au repos. Dans le référentiel où l’électron est au repos, le « comptage de cycles » de sa pulsation de Compton coïncide ainsi avec la mesure standard de son temps propre, ce qui assure la compatibilité avec la description relativiste usuelle.

Pour un nucléon composite, comme le proton, la situation est plus riche. La fréquence globale associée à sa masse au repos définit une cadence moyenne, mais cette cadence résulte de l’articulation hiérarchique entre :
\begin{itemize}
  \item les pulsations des régimes stationnaires locaux (quarks de valence) ;
  \item la dynamique plus lâche du régime collectif (mer de quarks) ;
  \item et le régime stationnaire global qui impose l’onde de surface de Compton.
\end{itemize}

Dans ce cadre, la constante de structure fine \(\alpha\) peut être interprétée comme un rapport entre la cohérence interne totale du proton et la fraction de cette cohérence qui demeure disponible à sa surface pour des couplages externes. La surface logique sur laquelle se manifeste l’onde de Compton globale ne doit pas être assimilée à une simple frontière géométrique : elle constitue un support fini de relations dynamiquement actives, issues de la hiérarchie interne (quarks valence et mer de quarks), qui restent accessibles au couplage avec des structures extérieures, telles que l’électron.

La valeur adimensionnée \(\alpha \approx 1/137\) caractérise dès lors la proportion du « budget de cohérence » du proton qui se projette dans ces relations de surface, relativement au niveau de saturation interne requis pour assurer le maintien de sa fermeture hiérarchique. Dans une perspective relationnelle, \(\alpha\) quantifie ainsi la « finesse » du couplage électromagnétique, non comme une constante imposée de manière exogène, mais comme la signature quantitative de la part de cohérence de surface qu’un proton peut allouer à ses interactions sans compromettre sa stabilité interne.

Le temps propre d’un proton peut, dans ce contexte, être envisagé comme le décompte des cycles de cette pulsation globale, chaque cycle intégrant la contribution de l’ensemble des sous-structures qui assurent la cohérence hiérarchique. Deux structures de même masse au repos peuvent ainsi partager une fréquence globale identique, tout en différant par la densité de pulsations internes qu’elles mobilisent pour maintenir leur cohérence.

Dans cette perspective, le temps métrique de la physique établie apparaît comme une moyenne effective sur ces décomptes hiérarchiques, plutôt que comme un paramètre fondamental indépendant. La dilatation des durées mise en évidence par la relativité peut alors être réinterprétée comme une modification des conditions de cohérence imposées à ces régimes stationnaires lorsqu’ils sont soumis à des configurations relationnelles extrêmes, plutôt que comme une déformation d’un temps abstrait. Il s’agit d’une relecture conceptuelle, et non d’une remise en question des résultats quantitatifs de la relativité restreinte et de la relativité générale : les effets de dilatation temporelle restent décrits par les mêmes relations mathématiques, mais sont envisagés ici comme des variations du budget de cohérence disponible pour les régimes stationnaires, plutôt que comme des anomalies affectant un temps fondamental pré-donné.

\section{La constante de structure fine comme signature de cohérence de surface}

La constante de structure fine \(\alpha \simeq 1/137\) apparaît, dans la formulation standard de l’électrodynamique quantique, comme un paramètre adimensionné caractérisant l’intensité du couplage électromagnétique. Dans le cadre relationnel développé ici, elle peut cependant être réinterprétée comme la mesure d’un rapport de cohérence entre l’organisation interne du proton et la fraction de cette organisation qui demeure accessible à sa surface pour des couplages avec l’environnement.

Dans une structure minimale de type \(\left(4,6\right)\), telle que celle associée à l’électron dans ce formalisme, une telle constante n’a tout simplement pas de pertinence : la fermeture y est considérée comme parfaite, de sorte qu’il n’existe ni hiérarchie interne, ni surface logique distincte du volume de cohérence global. Ce n’est qu’au franchissement d’un seuil hiérarchique, avec l’émergence d’une structure composite telle que le proton (articulant régimes stationnaires locaux, « mer » relationnelle et fuite partielle de cohérence), qu’une interface relationnelle se différencie d’un cœur saturé, rendant alors définissable un rapport entre cohérence interne et cohérence de surface. Dans cette perspective, \(\alpha\) peut être vue comme un effet dérivé de ce changement de régime : elle quantifie la proportion de cohérence que la structure composite peut laisser accessible à sa surface sans compromettre sa stabilité interne.

\subsection{Cohérence interne et surface logique du proton}

Le proton, précédemment décrit comme une structure composite minimale, est modélisé comme la combinaison de trois régimes stationnaires locaux de type quark, auxquels s’ajoute une mer relationnelle collective assurant la cohésion globale du système sans atteindre un état de saturation complète. Le nombre total de relations internes, obtenu par un comptage relationnel détaillé, est de l’ordre de
\[
r_{\mathrm{total}} \simeq 1{,}1\cdot 10^{4},
\]
ce qui quantifie le coût global de cohérence nécessaire au maintien de l’architecture hiérarchique du proton.

Ces relations ne sont toutefois pas toutes équivalentes du point de vue des couplages externes. Une fraction d’entre elles demeure confinée au sein des fermetures locales associées aux quarks de valence et à la mer de quarks, tandis qu’un sous-ensemble forme une surface logique active : un ensemble fini de relations par lesquelles le proton peut établir des couplages stables avec des structures externes, en particulier avec un régime stationnaire de type électron.

Ce sous-ensemble de relations de surface ne correspond pas à un bord géométrique au sens classique, mais à une « peau relationnelle » résiduelle, c’est‑à‑dire la part de structure relationnelle restant disponible après saturation des contraintes de cohérence interne. Les analyses combinatoires fondées sur la structure interne du proton conduisent à un nombre effectif de relations de surface de l’ordre de
\[
R_{\mathrm{surf}} \simeq 500
\]
qui sera utilisé par la suite comme caractéristique quantitative de la surface logique active.

\subsection{Rapport de masse et échelle de cohérence}

Le rapport entre la masse du proton et celle de l’électron,
\[
\mu = \frac{m_p}{m_e} \simeq 1836,
\]
reflète, dans le cadre relationnel adopté ici, l’écart de « coût de cohérence » entre, d’une part, un régime stationnaire composite (le proton) et, d’autre part, un régime stationnaire élémentaire minimal (l’électron). Il ne s’agit donc pas d’un simple paramètre empirique : ce rapport encode la manière dont deux types de fermetures hiérarchiques, caractérisées par des degrés de complexité interne très différents, se manifestent dans le langage quantitatif de la masse au repos.

Dans cette description, l’électron est associé à une structure minimale \(\left(4,6\right)\), dont la fermeture parfaite impose une pulsation propre directement corrélée à sa fréquence de Compton. À l’inverse, le proton met en œuvre une fermeture hiérarchique composite, comportant une fraction stationnaire de l’ordre de quelques pourcents et une fraction dynamique d’environ \(91\%\), attribuée à la mer relationnelle collective. Le rapport \(\mu\) quantifie alors l’écart entre ces deux régimes de cohérence, une fois celui-ci exprimé dans l’unité commune définie par \(\hbar\) et \(c\).

\subsection{Définition relationnelle de \texorpdfstring{\(\alpha\)}{alpha}}

Dans la section précédente, la surface active du proton a été estimée à environ \(R_{\mathrm{surf}} \simeq 500\) relations, représentant l’enveloppe minimale par laquelle le proton peut établir un couplage stable avec un électron. Cette grandeur permet désormais de réinterpréter la constante de structure fine comme un rapport de cohérence entre cette surface et l’architecture interne globale du proton.

L’objectif n’est pas ici de déterminer \(\alpha\) de manière univoque, mais de mettre en évidence qu’une même architecture protonique, déjà contrainte par la QCD et par les axiomes du LDP, conduit naturellement à une valeur proche de \(\alpha\) sans recours à un ajustement continu de paramètres libres.

Dans ce cadre, la constante de structure fine peut être exprimée sous la forme du rapport sans dimension
\[
\alpha^{-1} \simeq \frac{R_{\mathrm{surf}}^{2}}{\mu},
\]
où \(R_{\mathrm{surf}}\) désigne le nombre de relations actives de surface du proton et \(\mu\) le rapport de masse proton/électron.

En prenant \(R_{\mathrm{surf}} \simeq 500\) et \(\mu \simeq 1836{,}15\), on obtient
\[
\alpha^{-1} \simeq \frac{500^{2}}{1836{,}15} \approx 136,
\]
ce qui est compatible, à mieux que l’ordre du pour cent, avec la valeur expérimentale \(\alpha^{-1} \simeq 137{,}036\).

Cette estimation doit être comparée à la valeur mesurée \(\alpha_{\mathrm{exp}} \simeq 1/137{,}036\). L’écart relatif, de l’ordre du pour cent, est cohérent avec les approximations introduites tant dans le comptage des relations de surface que dans la modélisation hiérarchique interne du proton. Il suggère que l’interprétation de \(\alpha\) comme un ratio de cohérence de surface est, au minimum, quantitativement crédible.

Cette relation ne fournit pas une prédiction indépendante de \(\alpha\), dans la mesure où elle mobilise la valeur empirique du rapport de masse \(\mu\). Elle propose plutôt une relecture de la constante de structure fine en termes de cohérence de surface, numériquement compatible avec la valeur observée, sans établir pour l’instant que cette valeur serait la seule possible dans le cadre axiomatique considéré.

La proximité de la valeur obtenue pour \(\alpha^{-1}\) (erreur relative \(\sim 10^{-3}\)) ne résulte pas d’un ajustement continu, mais d’une combinaison entière fixée par l’architecture relationnelle (environ \(500\) relations actives sur \(930\)) et par le rapport de masse \(m_p/m_e\), sans introduction de nouveaux paramètres libres.

\subsection{Interprétation : finesse du couplage et stabilité de l’hydrogène}

L’interprétation relationnelle de \(\alpha\) ne consiste pas à introduire une nouvelle constante fondamentale, mais à requalifier une constante déjà présente dans le formalisme physique établi en tant que mesure d’une fraction de cohérence de surface. Une valeur de \(\alpha\) significativement plus élevée impliquerait qu’une portion plus importante du « budget » de cohérence interne du proton serait disponible à son interface : le couplage avec l’électron deviendrait alors excessivement fort, ce qui compromettrait la stabilité des états liés, en particulier celle de l’atome d’hydrogène. À l’inverse, une valeur de \(\alpha\) beaucoup plus faible correspondrait à une interface logique insuffisamment active pour autoriser la formation d’états liés stables.

La valeur observée de \(\alpha\) peut dès lors être interprétée comme le compromis unique entre deux contraintes antagonistes :
\begin{itemize}
  \item préserver la clôture hiérarchique interne du proton, en allouant l’essentiel de la cohérence à la structuration de sa compositionalité interne ;
  \item maintenir une « peau » relationnelle minimale mais suffisante pour permettre la stabilisation d’un régime stationnaire de type électronique à distance finie.
\end{itemize}

Cette interprétation est compatible avec le fait que l’énergie de liaison de l’atome d’hydrogène, calculée dans le cadre coulombien standard, peut être réécrite sous la forme d’un potentiel effectif dans lequel \(\alpha\) intervient comme paramètre quantifiant la fraction de cohérence de surface mobilisable pour le couplage, sans entrer en contradiction avec la valeur numérique de \(-13{,}6\,\mathrm{eV}\) au rayon de Bohr.

\subsection{\texorpdfstring{\(\alpha\), \(\hbar\)}{alpha, hbar} et la granularité des couplages}

Dans le cadre de la description relationnelle, la constante de Planck réduite \(\hbar\) a déjà été interprétée comme la granularité minimale de l’action associée à un régime stationnaire logique élémentaire : chaque cycle de la pulsation interne transporte un quantum d’action de module \(\hbar\), et l’énergie propre d’une structure s’écrit \(E = \hbar \omega\) dès lors qu’une réalisation physique de cette pulsation est disponible.

Dans cette perspective, la constante de structure fine peut être interprétée de façon analogue, non comme un paramètre fondamental dépourvu d’origine, mais comme la mesure d’une granularité de couplage de surface : elle quantifie la fraction du « budget de cohérence interne » d’un proton susceptible de se manifester à sa surface lors d’un couplage avec un électron. Une première estimation, fondée sur une surface active de l’ordre de \(R_{\mathrm{surf}} \simeq 500\) relations, conduisait déjà à une valeur voisine de \(\alpha^{-1}\) via la relation
\[
\alpha^{-1} \simeq \frac{R_{\mathrm{surf}}^{2}}{\mu},
\]
où \(\mu = m_p/m_e\) désigne le rapport des masses proton/électron.

Un affinage de cette estimation met en évidence une structure plus détaillée. Le nombre total de relations stationnaires de valence dans le proton est \(r_{\mathrm{valence}} = 930\). En considérant un noyau de cohérence « moyen » associé à chaque quark de valence, il est naturel d’examiner le tiers de cette quantité, \(r_{\mathrm{valence}}/3 = 310\), comme mesure d’un contenu moyen attaché à un régime stationnaire local (quark de valence) dans la hiérarchie relationnelle. En multipliant ce tiers par le nombre d’or \(\varphi\), défini par \(\varphi = (1+\sqrt{5})/2 \approx 1{,}618\), on obtient
\[
R_{\mathrm{surf}}^{\left(\varphi\right)}
\simeq \varphi \,\frac{r_{\mathrm{valence}}}{3}
\approx \varphi \,\frac{930}{3}
\approx 501{,}59,
\]
ce qui remplace la valeur approximative \(R_{\mathrm{surf}} \simeq 500\) par une expression entièrement construite à partir de données intrinsèques à l’architecture protonique (\(930\)) et d’un facteur géométrique élémentaire.

En insérant cette valeur raffinée dans la relation
\[
\alpha^{-1} \simeq \frac{{R_{\mathrm{surf}}^{\left(\varphi\right)}}^{2}}{\mu}
\approx \frac{501{,}59^{2}}{1836{,}15}
\approx 137{,}022,
\]
avec \(\mu \simeq 1836{,}15\), on obtient numériquement \(\alpha^{-1} \simeq 137{,}022\). Cette valeur est en accord, à mieux que quelques \(10^{-4}\) en valeur relative, avec la valeur expérimentale \(\alpha^{-1} \simeq 137{,}036\), et ce sans recourir à un terme ad hoc. La constante de structure fine peut ainsi être reformulée sous la forme
\[
\alpha^{-1}
\simeq \frac{\varphi^{2}}{\mu}
\left(\frac{r_{\mathrm{valence}}}{3}\right)^{2},
\]
ce qui l’exprime entièrement en fonction du rapport de masse \(\mu\), du contenu relationnel de valence \(r_{\mathrm{valence}}\) et du facteur \(\varphi\).

Dans ce cadre, le nombre d’or n’est pas introduit comme un simple motif esthétique ou comme une construction numérologique extérieure. Il est envisagé comme un candidat à une constante d’optimisation géométrique relationnelle : un facteur quantifiant la configuration la plus efficiente de distribution du contenu de cohérence d’un noyau stationnaire moyen vers la surface active, sous les contraintes de fermeture hiérarchique et de minimisation des fuites. Autrement dit, \(\varphi\) serait susceptible de jouer le rôle d’un invariant de topologie logique, caractérisant une répartition optimale des relations de cohérence à la surface des régimes stationnaires, plutôt que celui d’un nombre « magique » imposé de manière extrinsèque.

Cette interprétation s’inscrit dans la continuité d’une intuition déjà présente dans la littérature, selon laquelle la constante de structure fine relie un rapport de masses adimensionnel \(\mu\) à un nombre fini de canaux de couplage entre entités. Dans le cadre conceptuel proposé ici, cette intuition se reformule naturellement en termes de surfaces actives et de topologie relationnelle. Dans cette optique, \(\alpha\) quantifie la fraction de la granularité d’action fixée par \(\hbar\) qui demeure disponible à la surface pour des couplages, une fois la fermeture interne hiérarchique assurée.

La relation usuelle
\[
E = \hbar \omega
\]
appliquée à la pulsation propre de l’électron dans un état lié à un proton peut alors être réinterprétée comme l’expression quantitative du coût de cohérence de surface nécessaire au maintien simultané :
\begin{itemize}
  \item du régime stationnaire interne du proton ;
  \item du régime stationnaire élémentaire de l’électron ;
  \item et du couplage hiérarchique les reliant via la surface logique active caractérisée par \(\alpha\).
\end{itemize}

Dans cette perspective, \(\alpha\) cesse d’apparaître comme un paramètre dépourvu d’origine. Elle devient l’indicateur de la manière dont un régime stationnaire composite répartit son budget de cohérence entre fermeture interne et couplages de surface, préparant ainsi la mise en relation avec les constantes de couplage gravitationnel introduites dans les sections ultérieures.

\subsection{\texorpdfstring{\(k_B\)}{kB} et granularité thermique}

Dans ce cadre, la constante de Planck réduite \(\hbar\) détermine la granularité minimale de l’action associée à un cycle de pulsation d’un régime stationnaire isolé : chaque cycle transporte un quantum d’action \(\hbar\), et l’énergie propre d’une structure s’écrit \(E = \hbar \omega\) dès lors qu’une réalisation physique de cette pulsation est disponible. 

La constante de Boltzmann \(k_B\) joue un rôle analogue dès lors que l’on ne considère plus une structure isolée, mais un ensemble statistique de structures hiérarchisées. Dans la formulation usuelle de la physique statistique, \(k_B\) relie l’énergie thermique moyenne par degré de liberté à la température via des relations du type \(E_{\mathrm{th}} \sim k_B T\). Dans une perspective relationnelle, cette énergie thermique peut être interprétée comme la fraction du budget de cohérence globale effectivement mobilisée, en moyenne, dans les excitations collectives d’un grand nombre de cycles de pulsation.

Dans ce cadre, la température cesse d’être une grandeur primitive ; elle quantifie, de façon statistique, la population effective de cycles de cohérence excités au sein d’un ensemble de structures (électrons, nucléons, atomes, etc.). Le rôle de \(k_B\) est alors de relier ce comptage statistique des cycles excités à une échelle d’énergie moyenne, de manière strictement analogue au rôle de \(\hbar\), qui relie un cycle élémentaire à un quantum d’action. \(k_B\) fixe ainsi l’échelle énergétique caractéristique d’une fluctuation thermique par degré de liberté, tandis que \(\hbar\) fixe l’échelle d’une fluctuation d’action par cycle.

Dans l’état actuel du modèle LDP, \(k_B\) n’est pas dérivé de façon univoque à partir des axiomes. Une expression candidate est toutefois proposée ci-dessous, construite uniquement à partir de l’architecture de l’électron et du proton, puis évaluée numériquement. Cette expression doit être considérée comme une hypothèse structurée, et non comme une dérivation définitive.

En thermodynamique classique, la constante de Boltzmann \(k_B\) assure la médiation entre l’échelle microscopique et les observables macroscopiques telles que la température et l’entropie. Elle intervient dans la formule de Boltzmann \(S = k_B \ln W\) ainsi que dans le lien \(E \sim k_B T\) qui associe un niveau moyen d’agitation à une énergie thermique caractéristique. Dans une lecture relationnelle, \(k_B\) doit pouvoir être interprétée comme la granularité minimale de la cohérence lorsqu’une structure logique est soumise à un ensemble de configurations possibles, plutôt que comme un simple facteur de proportionnalité empirique.

De manière parallèle, la construction suivante examine si la même architecture logique est en mesure de produire une échelle thermique compatible avec la valeur expérimentale de \(k_B\), à l’intérieur d’une marge d’erreur contrôlée.

Dans le cadre du formalisme LDP, la même pulsation minimale qui définit \(\hbar\) pour la structure \(\left(4,6\right)\), identifiée à l’électron, fournit une période de référence \(\tau_0\). Cette dernière est déterminée à partir de la longueur d’onde de Compton de l’électron \(\lambda_C\) et de la vitesse de propagation informationnelle \(c\) :
\[
\tau_0 = \frac{\lambda_C}{c},
\]
avec \(\lambda_C = \frac{h}{m_e c}\). Cette période peut être interprétée comme la durée logique d’un cycle de cohérence minimal de l’électron au repos.

Pour une structure composite telle que le proton, l’énergie de cohérence ne se trouve pas concentrée uniquement sur la structure stationnaire \(\left(4,6\right)\), mais se répartit au sein d’une architecture hiérarchique impliquant :
\begin{itemize}
  \item des sous-structures stationnaires associées aux quarks de valence, caractérisées par \(n_u=24\), \(n_d=28\), \(r_u=276\), \(r_d=378\), conduisant à un total de \(930\) relations pour les trois quarks ;
  \item une « mer » de relations dynamiques \(R_{\mathrm{mer}} = 10\,087\), pour un total \(R_{\mathrm{tot}} = 11\,017\), de sorte que la fraction stationnaire représente environ \(9\%\) et la fraction dynamique \(91\%\) ;
  \item une surface active de l’ordre de \(501{,}6\) relations, directement impliquées dans le couplage avec l’électron et déjà mobilisées dans la construction de la constante de structure fine \(\alpha\).
\end{itemize}

Dans cette perspective, la fraction dynamique \(R_{\mathrm{mer}}/R_{\mathrm{tot}}\) quantifie le degré de liberté relationnel accessible aux configurations thermiques locales, tandis que le taux de fuite \(\varepsilon\), introduit plus haut comme violation minimale de la cohérence triangulaire à l’interface proton–univers, mesure la part de cohérence qui se dissipe dans le fond relationnel global.

Une expression candidate pour la constante de Boltzmann peut, dans ce cadre, être formulée en reliant :
\begin{itemize}
  \item la tension rythmique minimale \(h/\tau_0\) associée à un cycle de pulsation de l’électron,
  \item le nombre de relations actives intervenant dans la cohérence logique du proton,
  \item le rapport de masse \(m_p/m_e\), qui fixe le rapport des fréquences propres entre les deux structures,
  \item la fraction dynamique de la « mer » et le taux de fuite \(\varepsilon\), qui quantifient la part de cohérence effectivement disponible sous forme d’agitation thermique.
\end{itemize}

Sous ces hypothèses, on obtient l’expression suivante, exempte de tout degré de liberté continu supplémentaire :
\[
k_B^{\mathrm{LDP}}
= \frac{h}{\tau_0}\,
  \frac{n_{\mathrm{rel}}^{e}}{\left(m_p/m_e\right)^{2}}\,
  \frac{R_{\mathrm{mer}}}{R_{\mathrm{tot}}}\,
  \frac{\varepsilon}{4},
\]
où :
\begin{itemize}
  \item \(h\) désigne la constante de Planck, 
  \item \(\tau_0\) la période de la pulsation propre de l’électron au repos,
  \item \(n_{\mathrm{rel}}^{e} = 6\) le nombre de relations de la structure \((4,6)\),
  \item \(m_p/m_e = 1836{,}15267343\) le rapport des masses proton/électron, 
  \item \(R_{\mathrm{mer}} = 10\,087\) et \(R_{\mathrm{tot}} = 11\,017\) le nombre de relations associé respectivement à la mer et au proton total,
  \item \(\varepsilon \approx 0{,}0060923\) le taux de fuite déterminé précédemment à partir de la cohérence du proton et réutilisé dans les expressions de \(G\) et de \(H_0\),
  \item le facteur \(1/4\) répartit le flux de fuite sur les quatre groupes logiques structurant le nucléon.
\end{itemize}

Cette formulation ne repose sur l’introduction d’aucun paramètre continu ajustable : l’ensemble des entiers intervenant est fixé par l’architecture relationnelle \((4,6)\) pour l’électron, \(24, 28, 10\,087, 11\,017\) pour le proton, tandis que \(\varepsilon\) correspond au même taux de fuite que celui apparaissant dans les expressions de la constante gravitationnelle et du taux d’expansion cosmique.

En substituant les valeurs numériques de référence \(h = 6{,}62607015\times 10^{-34}\,\mathrm{J\cdot s}\), \(m_e \approx 9{,}1093837\times 10^{-31}\,\mathrm{kg}\), \(m_p \approx 1{,}6726219\times 10^{-27}\,\mathrm{kg}\), \(c = 2{,}99792458\times 10^{8}\,\mathrm{m\cdot s^{-1}}\), etc., on obtient :
\[
k_B^{\mathrm{LDP}} \approx 1{,}3802\times 10^{-23}\,\mathrm{J\cdot K^{-1}},
\]
à comparer à la valeur CODATA
\[
k_B^{\mathrm{exp}} = 1{,}380649\times 10^{-23}\,\mathrm{J\cdot K^{-1}}, 
\]
qui est à présent définie comme une constante exacte. 

Il en résulte une erreur relative de l’ordre de
\[
\frac{\left|k_B^{\mathrm{LDP}} - k_B^{\mathrm{exp}}\right|}{k_B^{\mathrm{exp}}}
\approx 3\times 10^{-2}.
\]

Une telle précision demeure modeste au regard des déterminations expérimentales contemporaines, et l’expression proposée ne doit donc pas être interprétée comme une dérivation définitive de \(k_B\) dans le cadre LDP. Elle met néanmoins en évidence que, en mobilisant exclusivement des quantités déjà fixées par l’architecture \((4,6)\) et par la structure hiérarchique du proton (y compris le même taux de fuite \(\varepsilon\) qui intervient dans \(G\) et \(H_0\)), il devient possible de relier la granularité thermique classique au même fondationnel relationnel que la masse, la charge électrique et l’interaction gravitationnelle. Ce type de convergence, même partielle, suggère que la thermodynamique macroscopique pourrait être interprétée comme une manifestation statistique de contraintes de cohérence logique déjà à l’œuvre au niveau des structures nucléoniques.

Cette expression conserve un statut conjectural : elle repose, d’une part, sur l’interprétation des entiers \(6\), \(10\,087\), \(276\), \(378\) comme des nombres de relations dans l’architecture protonique, et, d’autre part, sur la réutilisation du même paramètre de fuite \(\varepsilon\) que celui employé pour \(G\) et \(H_0\). Néanmoins, le fait qu’une combinaison entièrement déterminée par la structure logique (c’est‑à‑dire sans introduction d’un paramètre continu libre supplémentaire) conduise à une valeur de \(k_B\) avec une erreur relative de l’ordre de \(3\times 10^{-2}\) étaye l’hypothèse selon laquelle la thermodynamique macroscopique peut être reliée au même substrat relationnel que la microphysique et la cosmologie.

\section{Proton, trous noirs et information de surface}

Dans le cadre de la physique contemporaine, deux objets a priori disjoints occupent les extrémités opposées des échelles de masse, de taille et d’intensité de champ : le proton, constituant élémentaire de la matière baryonique ordinaire, et le trou noir, solution gravitationnelle extrême de la relativité générale. Pourtant, ces deux systèmes peuvent être interprétés comme des configurations hautement saturées en information, dont le contenu physique pertinent est principalement associé à une surface finie, plutôt qu’à un volume interne traité comme homogène et indifférencié.

\subsection{Le proton comme noyau d’information hiérarchique}

Dans la description standard, le proton ne peut plus être modélisé comme un simple triplet de quarks ponctuels. La chromodynamique quantique (QCD) montre que sa masse et son énergie résultent majoritairement de la dynamique non linéaire du champ de gluons et de la « mer » de quarks et d’antiquarks virtuels, plutôt que des seules masses des trois quarks de valence. La structure interne du proton s’apparente, dès lors, à un réseau multi-échelle de degrés de liberté, parfois décrit de manière heuristique comme une organisation de type fractal, où les quarks de valence ne constituent que la partie émergente d’une structure beaucoup plus dense et corrélée.

Dans le cadre relationnel adopté ici, cette complexité interne est reformulée comme une architecture hiérarchique au sein de laquelle une surface active finie regroupe l’essentiel des canaux de couplage avec l’électron. Cette représentation permet de réinterpréter certaines constantes sans dimension comme des rapports caractérisant le degré de cohérence entre la dynamique interne et les propriétés effectives de surface.

\subsection{Trous noirs : l’horizon comme support d’entropie}

Sur le versant gravitationnel, les travaux de Bekenstein et Hawking ont établi que les trous noirs ne peuvent plus être décrits comme de simples régions causales fermées d’où aucune information ne peut s’échapper. L’attribution d’une entropie
\[
S_{\mathrm{BH}} = \frac{k c^{3} A}{4 G \hbar},
\]
proportionnelle à l’aire \(A\) de l’horizon, s’avère nécessaire pour garantir la validité de la seconde loi de la thermodynamique lorsqu’un système matériel est absorbé par un trou noir. Cette entropie est interprétée comme une entropie informationnelle : l’horizon supporte un contenu informationnel maximal, c’est-à-dire un nombre maximal de bits que peut contenir une région donnée de l’espace-temps, conformément à l’esprit de la borne de Bekenstein et du principe holographique.

Hawking a ensuite démontré que les trous noirs émettent un rayonnement de nature thermique, ce qui permet d’assigner une température bien définie à l’horizon et confirme que ces grandeurs ne constituent pas de simples constructions formelles : le trou noir satisfait des lois de la thermodynamique physiques, dans lesquelles l’aire de l’horizon joue effectivement le rôle de variable entropique. Dans plusieurs cadres théoriques contemporains (paradigme de la membrane, modèles d’« horizon comme dispositif de mesure », approches holographiques), les degrés de liberté pertinents sont explicitement localisés sur la surface de l’horizon, plutôt que dans le volume interne. La plus massive et la plus étendue des structures gravitationnelles connues se trouve ainsi décrite comme un objet dont le contenu informationnel effectif est concentré sur une surface finie, mesurée en unités de Planck.

\subsection{Un mini trou noir inversé}

À la lumière de ces deux descriptions, le rapprochement peut être formulé de manière plus systématique. Celles-ci convergent vers un même schéma conceptuel : une organisation extrêmement dense de degrés de liberté, dont la complexité interne se trouve, du point de vue d’un observateur externe, résumée par un nombre fini de paramètres définis sur une surface. Le parallèle peut alors être explicité de la manière suivante :
\begin{itemize}
  \item Dans le cas du trou noir, l’horizon encode la totalité de l’entropie accessible : l’information concernant tout ce qui a traversé l’horizon est, en un certain sens, « stockée » sur cette surface, et la géométrie extérieure ne dépend que de quelques charges globales (masse, charge électrique, moment cinétique).
  \item Dans le cas du proton relationnel, l’architecture interne hiérarchisée (quarks de valence saturés, « mer » relationnelle dynamique) fixe un budget de cohérence global, mais seule une surface active finie contribue effectivement aux couplages externes. C’est cette « peau » relationnelle qui se manifeste à travers des constantes sans dimension telles que \(\alpha\) ou des rapports de masse caractéristiques.
\end{itemize}

Dans ce cadre, il devient pertinent de décrire le proton comme un mini trou noir inversé. Alors que le trou noir astrophysique concentre l’information macroscopique d’une région spatiale sur un horizon géométrique, le proton concentre l’information d’une structure composite minimale sur une fermeture relationnelle interne, dont une surface logique active joue, vis-à-vis de l’électron et des autres champs, un rôle fonctionnel analogue à celui d’un horizon. Dans les deux situations, la dynamique observable depuis l’extérieur est contrainte par la structure d’information accessible sur une surface de dimension finie, plutôt que par les détails du « volume » interne des degrés de liberté sous-jacents.

\subsection{\emph{It from bit} : Hawking, Bekenstein et Wheeler}

Wheeler a synthétisé cette convergence conceptuelle dans la formule devenue canonique « it from bit » : l’hypothèse selon laquelle chaque entité du monde physique dérive, en dernière analyse, de degrés de liberté informationnels discrets, structurés en réponses binaires oui/non, plutôt que d’émerger d’une substance continue préexistante. Dans le cas des trous noirs, ce programme acquiert une réalisation quantitative avec l’entropie de Bekenstein–Hawking : l’horizon porte une information entropique qui doit être traitée comme une grandeur physique objective et mesurable.

Le cadre relationnel proposé ici transpose cette intuition au domaine hadronique. Le proton y est modélisé comme une organisation combinatoire de relations accord/désaccord, dont une partie est figée en régimes stationnaires locaux (les quarks de valence) tandis qu’une autre soutient une dynamique collective (la « mer relationnelle »). La surface active du proton, définie comme l’enveloppe minimale de relations par lesquelles cette structure peut entrer en interaction avec un électron décrit par \(\left(4,6\right)\), se voit alors attribuer un double rôle : d’une part, elle comprime la cohérence interne du proton en un nombre fini de canaux de couplage ; d’autre part, elle permet de réinterpréter certaines constantes sans dimension, telle la constante de structure fine, comme des rapports entre cohérence interne et surface d’interaction, plutôt que comme des paramètres fondamentaux dépourvus d’origine conceptuelle.

Dans cette perspective, la thèse « it from bit » reçoit une mise en œuvre opérationnelle à deux échelles distinctes : à l’échelle macroscopique, avec les trous noirs de Bekenstein–Hawking, et à l’échelle microscopique, avec le proton envisagé comme structure relationnelle composite minimale. Dans les deux cas, une organisation finie de l’information, localisée sur une surface, acquiert un statut ontologique équivalent à celui d’un objet physique : l’horizon pour le trou noir, la surface active pour le proton. Le parallèle n’implique pas une identification du proton à un trou noir au sens de la relativité générale ; il souligne plutôt que, du point de vue de la cohérence interne et de la structuration de l’information, la plus petite brique composite de la matière et la plus vaste structure gravitationnelle connue obéissent à une logique commune : la réalité phénoménale émerge là où une structure informationnelle atteint un seuil de saturation qui la rend irréductible à un simple formalisme descriptif.

\section{De la coexistence à une métrique émergente}

La quantification relationnelle des structures fermées, ainsi que l’interprétation du temps propre comme comptage discret de cycles de pulsation, caractérisent une structure hiérarchique considérée isolément. Néanmoins, la phénoménologie physique ne se réduit pas à des entités solitaires : électrons, protons et neutrons coexistent en grand nombre, se combinent en noyaux, en atomes, puis en corps macroscopiques, tout en préservant leurs identités structurelles.

Dans le cadre logique retenu, la coexistence de plusieurs structures hiérarchiques requiert que leurs régimes stationnaires respectifs puissent maintenir leur cohérence interne sans entrer dans un conflit destructif mutuel. Cette condition ne spécifie pas encore une dynamique d’interaction à proprement parler ; elle formule une contrainte de compatibilité : certaines configurations de fermetures hiérarchiques peuvent partager un même « voisinage logique », tandis que d’autres doivent être mutuellement exclues afin d’éviter l’apparition de contradictions structurelles.

Il en découle l’émergence d’un réseau de voisinages, au sein duquel chaque structure hiérarchique est directement corrélée à un nombre fini d’autres structures avec lesquelles elle est compatible au sens des contraintes de charge et de cohérence. Ce réseau n’implique a priori ni continuité, ni dimension fixée, ni géométrie préexistante ; il encode exclusivement les conditions de possibilité d’une coexistence stable.

Pour les structures composites, telles que les protons, cette compatibilité est directement modulée par la fuite minimale de cohérence associée à chaque hiérarchie. La fraction \(\epsilon\) de cohérence qui ne se referme pas intégralement à l’intérieur d’un proton doit être nécessairement prise en charge par le réseau de voisinages : elle se manifeste alors comme une contrainte additionnelle sur les modalités selon lesquelles les fermetures hiérarchiques peuvent partager un même voisinage logique. En conséquence, une configuration qui accroît localement la densité de protons modifie le schéma des fuites de cohérence et impose un réajustement des conditions de compatibilité dans son environnement immédiat.

Dans ce réseau, la notion de voisinage induit une topologie élémentaire : deux configurations sont dites voisines si leurs fermetures respectives peuvent être simultanément maintenues sans violation des contraintes de cohérence. L’« espace » qui en résulte ne constitue donc pas un contenant neutre préalablement donné, mais émerge comme une propriété collective de l’organisation des relations de coexistence entre structures stationnaires.

Sur ce fond, la métrique minimale introduite précédemment acquiert une signification opérationnelle précise : elle fournit un outil de comparaison des séparations entre structures hiérarchiques, non pas en termes de distances présupposées, mais en termes de coût différentiel de cohérence requis pour les maintenir dans un même voisinage logique. Une configuration sera qualifiée de « plus éloignée » qu’une autre si son maintien compatible avec l’ensemble des structures déjà présentes exige un réajustement plus important de la distribution de cohérence.

Le passage d’une simple relation de compatibilité binaire à une comparaison graduée des séparations constitue ainsi la première étape vers la définition d’une métrique émergente au sens physique. Il prépare l’interprétation de la gravitation, non comme une force additionnelle s’exerçant sur un espace préexistant, mais comme l’expression systématique des réajustements de cohérence induits par la modification de la distribution des structures hiérarchiques.

Il convient de souligner que cette métrique émergente ne se manifeste pas au niveau d’une structure minimale isolée de type \((4,6)\). Tant qu’une unique pulsation stationnaire élémentaire est réalisée, il n’existe ni réseau de voisinages, ni nécessité de comparer des séparations : aucune géométrie ne peut alors être définie au sens physique. Ce n’est qu’au-delà d’un certain seuil, lorsque des structures composites hiérarchiques coexistent en nombre, chacune associée à une fuite minimale \(\epsilon\), que le coût des réajustements de cohérence devient une grandeur collective, interprétable comme une courbure effective de la métrique émergente. La gravitation peut alors être comprise comme la traduction quantitative de ces réajustements, comme cela sera explicité dans la section suivante.

\section{Gravitation interprétée comme modulation du coût de cohérence}

Une structure hiérarchique fermée, telle qu’un électron ou un nucléon, se caractérise par un coût de cohérence interne, quantifiable par le nombre de relations nécessaires à son maintien ainsi que par la pulsation propre associée à sa masse au repos. Ce coût ne se superpose pas à un substrat neutre : il représente la fraction du budget global de cohérence que le cadre relationnel doit allouer à cette structure afin d’en garantir la stabilité dynamique.

Lorsque plusieurs structures hiérarchiques coexistent, leurs coûts de cohérence ne s’additionnent pas simplement sur un arrière-plan indifférent. Chacune de ces structures impose en effet des contraintes sur les modalités de distribution de la cohérence dans son voisinage logique. Plus la densité de structures fermées est élevée dans une région donnée, plus la part de cohérence disponible pour soutenir simultanément l’ensemble des régimes stationnaires devient restreinte et fortement contrainte.

Considérons deux nucléons initialement très éloignés au sens de la métrique émergente : leurs fermetures internes peuvent, au premier ordre, être approximées comme quasi indépendantes, et les réajustements de cohérence imposés au cadre relationnel demeurent alors faibles. À mesure que l’on rapproche ces nucléons dans le réseau de voisinages logiques, la compatibilité de leurs régimes stationnaires requiert un réaménagement de plus en plus important du budget de cohérence dans la région commune. Le coût de cohérence total attribué à cette région s’en trouve accru.

Cette augmentation ne résulte pas d’une injection d’énergie externe, mais d’une redistribution de la cohérence préexistante : le cadre relationnel doit concentrer une quantité accrue de ressources afin de maintenir, dans une même région, plusieurs fermetures hiérarchiques mutuellement compatibles. Dans une description métrique, cette concentration se manifeste sous la forme d’une courbure effective : les trajectoires qui minimisent les réajustements de cohérence cessent alors d’être rectilignes et sont déviées au voisinage des zones où le coût de cohérence atteint son maximum.

La fuite de cohérence minimale associée à chaque proton acquiert, dans ce contexte, un rôle central. Tant que la fermeture de cohérence demeure strictement assurée, comme dans le cas d’une structure élémentaire \((4,6)\) isolée, aucune contrainte n’est imposée à la métrique émergente : l’intégralité de la circulation de cohérence reste confinée au sein de la structure elle‑même. En revanche, dès qu’une hiérarchie composite présente une fuite \(\varepsilon\), une fraction de cette circulation doit être prise en charge par le cadre relationnel environnant. La gravitation peut dès lors être interprétée comme l’effet collectif de ces fuites minimales, lorsque de nombreux nucléons imposent simultanément leur contrainte au réseau de voisinages logiques.

Dans ce cadre conceptuel, la gravitation peut être comprise comme la manifestation de la tendance des structures hiérarchiques à évoluer le long de trajectoires correspondant à des chemins de moindre coût de cohérence. Ces structures ne « s’attirent » pas au sens d’une force newtonienne : elles évoluent dans un cadre où la métrique émergente est continuellement réajustée en fonction de la distribution spatiale et hiérarchique des coûts de cohérence, et leurs trajectoires s’alignent sur les directions qui minimisent ces réajustements.

Dans ce formalisme, la « masse gravitationnelle » est réinterprétée comme une mesure du coût de cohérence hiérarchique associé à une structure fermée, tandis que la courbure de la métrique émergente encode la distribution de ce coût au sein du réseau de voisinages logiques. La gravitation cesse ainsi d’apparaître comme un champ additionnel défini sur un espace préexistant ; elle émerge plutôt comme la manifestation macroscopique de la redistribution collective de la cohérence imposée par la coexistence et la persistance de structures hiérarchiques au sein du substrat relationnel.

\subsection*{Esquisse des 18 filtres de cohérence et de la gravitation résiduelle}

L’expression proposée pour la constante de couplage gravitationnel effective,
\[
G_{\mathrm{LDP}} \simeq \frac{2\,\hbar c}{m_p^{2}}\,
\frac{1}{\alpha^{18\left(1-\varepsilon\right)}},
\]
met en jeu un exposant discret \(18\), modulé par le facteur de fuite \((1-\varepsilon)\). Dans l’analyse précédente, ce nombre a été interprété comme le « nombre minimal d’itérations hiérarchiques » nécessaires pour que la fuite de cohérence associée à un proton se traduise, à grande échelle, par une déformation mesurable de la métrique émergente. La présente section propose une esquisse plus détaillée de ces dix-huit étapes de filtrage, en les reliant aux contraintes déjà introduites dans la construction de la brique \((4,6)\), des blocs de valence \(u/d\) et de l’architecture hiérarchique du proton.

Dans le cadre LDP, la cohérence n’est pas distribuée de manière continue sur un fond neutre, mais circule sous forme de « paquets » au sein de structures relationnelles tendant à se refermer sur elles-mêmes. À chaque franchissement d’un filtre de fermeture par un paquet de cohérence, une fraction de sa capacité de projection vers l’extérieur est dissipée au profit de la stabilisation interne du système ou de l’établissement de couplages de plus courte portée (par exemple électromagnétiques). La gravitation peut alors être interprétée comme l’effet résiduel de ce processus : ce qui subsiste, après un nombre fini de filtrages, se manifeste sous la forme d’une contrainte extrêmement diffuse, insuffisante pour organiser la matière à l’échelle locale, mais suffisante, par agrégation, pour induire une courbure de la métrique émergente.

Dans cette perspective, l’exposant \(18\) représente le nombre de filtres de cohérence pertinents que les paquets doivent traverser entre l’échelle de la brique minimale \((4,6)\) et l’échelle des agrégats macroscopiques à laquelle la constante \(G\) est mesurée.

Sans prétendre au caractère exclusif de cette décomposition, on peut proposer le regroupement suivant, fondé exclusivement sur les contraintes déjà imposées au modèle.

\paragraph{Six filtres élémentaires (brique \((4,6)\)).}

Ces conditions caractérisent la possibilité même d’un régime stationnaire minimal, identifié à l’électron :
\begin{itemize}
  \item existence d’une pulsation binaire interne, interprétée comme un 2-cycle logique portant sur les signes des arêtes ;
  \item cohérence triangulaire locale, imposant que les produits de signes soient égaux à \(+1\) sur l’ensemble des triangles présents ;
  \item complétude du graphe \((4,6)\) : tous les couples de sommets sont reliés par une arête ;
  \item connexité du bloc de triangles cohérents : au moins deux triangles partagent une arête, de manière à assurer la couverture de l’ensemble de la structure ;
  \item équivalence structurelle des sommets : absence de hiérarchie interne entre les sommets de \((4,6)\), tous soumis aux mêmes contraintes combinatoires ;
  \item minimalité : aucun sous-graphe propre de \((4,6)\) ne satisfait simultanément l’ensemble de ces propriétés.
\end{itemize}

Pris conjointement, ces six filtres garantissent que la cohérence logique demeure strictement confinée à l’intérieur d’une brique élémentaire, sans aucune propagation indue vers les structures externes.

\paragraph{Huit filtres hiérarchiques internes (proton).}

Ces filtres spécifient les contraintes nécessaires à l’agrégation de plusieurs briques en une structure composite hiérarchique (le proton), tout en maintenant la stabilité globale du système :
\begin{itemize}
  \item imposition d’une multiplicité égale à \(4\) pour le nombre d’entités par quark, formalisée comme une superposition de tétrades logiques ;
  \item exigence de fermeture interne de chaque quark, modélisée par un graphe complet \(K_{n_q}\), garantissant une cohérence triangulaire interne quasi saturée ;
  \item sélection des tailles \(n_u = 24\) et \(n_d = 28\) comme première combinaison satisfaisant simultanément (i) la complétude interne, (ii) une asymétrie minimale et (iii) la maximisation de la fonction d’optimisation \(\Omega\) ;
  \item réalisation d’une charge globale \(+1\) par la combinaison interne \((u,u,d)\), choisie de manière à être compatible avec la charge \(-1\) de la brique \((4,6)\) ;
  \item décomposition du contenu relationnel en un cœur stationnaire (\(930\) relations) et une mer dynamique (\(10\,087\) relations), avec une fraction relative d’environ \(9\%/91\%\) ;
  \item garantie d’une connectivité suffisante de la mer dynamique pour assurer la propagation de la cohérence entre les blocs de valence, sans atteindre un degré de saturation globale susceptible de compromettre la hiérarchie interne ;
  \item postulation de l’existence d’une surface active finie, sous-ensemble strict des relations de valence, apte à se coupler logiquement à l’électron au moyen de triangles mixtes cohérents ;
  \item maintien de la fermeture globale du proton malgré cette décomposition interne en blocs de valence, mer dynamique et surface active.
\end{itemize}

Ces huit filtres spécifient le mode de redistribution des paquets de cohérence au sein du proton : une fraction demeure confinée dans les fermetures locales associées aux quarks, une autre alimente la mer dynamique, tandis qu’une partie résiduelle reste disponible à la surface pour assurer des couplages externes.

\paragraph{Quatre filtres de projection vers l’environnement.}

Quatre filtres supplémentaires régissent la transition des paquets de cohérence résiduels hors du régime strictement interne, de sorte qu’ils puissent se manifester dans le cadre relationnel environnant :
\begin{itemize}
  \item \emph{Filtre de compatibilité logique proton–électron} : existence de triangles mixtes cohérents assurant un couplage de charges \(+1/-1\), sans engendrer de croissance divergente du nombre de triangles violés ;
  \item \emph{Filtre de dilution par les couplages de surface} : atténuation de la part de cohérence effectivement accessible via les couplages de surface, modélisée par la constante de structure fine \(\alpha\), qui fixe la résolution (« finesse ») maximale de ces échanges ;
  \item \emph{Filtre d’incomplétude logique de la fermeture} : impossibilité de réaliser une fermeture simultanément parfaite à tous les niveaux (briques élémentaires, blocs \(u/d\), mer, surface), impliquant l’existence d’une fuite minimale \(\varepsilon > 0\) hors des \(11\,017\) relations internes ;
  \item \emph{Filtre de propagation hiérarchique de la fuite} : transfert hiérarchique de cette fuite à travers les niveaux supérieurs d’agrégation (nucléons, noyaux, atomes, matière ordinaire, structures astrophysiques), jusqu’à ce qu’elle soit interprétable, dans une description métrique, comme une courbure effective paramétrée par la constante effective \(G_{\mathrm{LDP}}\).
\end{itemize}

Ces quatre filtres explicitent le transfert progressif des paquets de cohérence, depuis la sphère strictement interne (proton + électrons) vers le cadre relationnel collectif, puis vers l’échelle macroscopique. À chaque étape, la contribution directe de la cohérence est intégralement ou partiellement absorbée soit par les couplages électromagnétiques (via \(\alpha\) et la surface active), soit par les mécanismes de stabilisation interne de la hiérarchie. Seule une fraction résiduelle, extrêmement faible, subsiste sous la forme d’une contrainte de fond, que l’on identifie à la composante gravitationnelle.

L’identification proposée ci-dessus des dix-huit filtres doit être considérée comme une hypothèse structurée plutôt que comme une dérivation combinatoire définitivement établie. Une analyse plus approfondie pourrait en effet mettre en évidence des dépendances non triviales entre certains de ces filtres ou, au contraire, conduire à leur redécomposition en sous-contraintes plus élémentaires. Cette construction suggère néanmoins qu’il est légitime d’interpréter l’exposant \(18\) non comme un entier arbitraire, mais comme un codage discret du nombre effectif de couches de contraintes de cohérence que les paquets doivent franchir pour que la gravitation émerge en tant qu’effet résiduel de la clôture hiérarchique des formes.

Des travaux ultérieurs devront préciser, dans un cadre strictement combinatoire (formalisé en termes de graphes signés, de classes de triangles et d’orbites de la dynamique interne), la correspondance entre ces filtres et des invariants mathématiquement bien définis. Il s’agira également d’établir en quel sens l’exposant \(18\) est minimal pour reproduire simultanément (i) la structure interne du proton, (ii) la valeur de la constante de couplage \(\alpha\) et (iii) l’ordre de grandeur empirique de la constante gravitationnelle \(G\).

\section{Paramètre de fuite \texorpdfstring{\(\varepsilon\)}{epsilon} et amplification hiérarchique}

La modélisation du proton en tant que structure composite minimale met en évidence un phénomène de \emph{fuite de cohérence} : une fraction de la circulation relationnelle nécessaire au maintien de sa stabilité ne se referme pas intégralement au sein de ses relations internes, mais transite par le cadre relationnel environnant. Cette fuite ne doit pas être assimilée à une perte au sens énergétique du terme ; elle traduit plutôt le fait qu’une hiérarchie ne peut simultanément réaliser une clôture maximale de l’ensemble de ses sous-systèmes et une clôture globale parfaitement isolée.

Afin de quantifier cette imperfection contrôlée, on introduit un paramètre de fuite \(\mathbf{\varepsilon}\). À l’échelle du proton, \(\varepsilon\) mesure la déviation relative entre, d’une part, une clôture hiérarchique idéale, dans laquelle l’ensemble des relations se refermerait intégralement au sein de la structure, et, d’autre part, la clôture effective, qui laisse une faible fraction de la cohérence se projeter vers l’extérieur. Ce paramètre constitue ainsi une grandeur adimensionnée, intrinsèquement associée à la hiérarchie composite elle-même, et non à une métrique externe préalablement imposée.

Cette fuite locale n’est pas confinée à l’échelle hadronique. Lorsque de nombreux protons (ou noyaux) s’assemblent en structures de plus grande taille (atomes, solides, planètes, étoiles), les fuites de cohérence associées à chaque sous-système ne s’additionnent pas de manière strictement linéaire ; elles se combinent et s’amplifient selon une organisation hiérarchique. À chaque niveau de structuration, une fraction, fût-elle infime, de la cohérence interne échappe à la clôture parfaite et se transmet aux niveaux supérieurs, de sorte que l’effet cumulé peut devenir détectable à l’échelle macroscopique, où il se manifeste sous la forme d’une déformation de la métrique émergente.

Dans ce cadre conceptuel, le paramètre \(\varepsilon\) occupe une place centrale dans la tentative de reformulation du phénomène gravitationnel. À l’échelle d’un proton isolé, sa valeur est extrêmement faible et ne se traduit par aucune déviation directement observable. En revanche, lorsqu’on considère la propagation de cette « fuite » à travers plusieurs niveaux hiérarchiques — depuis la structure interne du proton jusqu’à la matière ordinaire, puis jusqu’aux configurations astrophysiques — ce même paramètre intervient dans un exposant effectif qui gouverne la manière dont la cohérence se trouve progressivement diluée à grande échelle. C’est précisément cette dépendance exponentielle qui réapparaît dans l’expression proposée pour une constante de couplage gravitationnel effective.

À ce stade de développement du modèle, il convient de souligner que la valeur numérique de \(\varepsilon\) n’est pas encore obtenue par une dérivation purement combinatoire. Elle est ajustée de telle sorte que l’expression de la constante de couplage gravitationnel reproduise la valeur expérimentale de la constante de gravitation newtonienne. Le paramètre \(\varepsilon\) doit par conséquent être interprété comme une quantité fixée par les données observationnelles, et non comme une grandeur entièrement prédite à partir des axiomes du modèle. Toute dérivation ultérieure de \(\varepsilon\) devra retrouver cette valeur comme conséquence de la structure hiérarchique postulée ; à défaut, la reformulation proposée de la gravitation demeurerait incomplète du point de vue théorique.

La section suivante exploite ce paramètre de fuite, conjointement avec la structure interne du proton et la constante de structure fine, afin d’établir une expression pour la constante de couplage gravitationnel. L’objectif est d’évaluer dans quelle mesure le cadre relationnel permet de rendre compte de l’ordre de grandeur de l’interaction gravitationnelle, sans introduire de nouvelle constante fondamentale indépendante.

\section{Vers une constante de couplage gravitationnel}

Si la gravitation est interprétée comme la manifestation d’une variation du coût de cohérence collective, il devient alors possible de définir une constante de couplage gravitationnel en tant que mesure de la sensibilité de la métrique émergente à ce coût. Une telle constante établirait une relation quantitative entre le contenu hiérarchique d’une structure fermée et la courbure qu’elle induit dans son voisinage.

Dans ce cadre, chaque nucléon peut être caractérisé par trois paramètres fondamentaux :
\begin{itemize}
  \item un nombre total de relations internes, quantifiant son coût de cohérence hiérarchique ;
  \item une longueur d’onde de Compton associée à sa masse au repos, fournissant une échelle minimale de fermeture spatiale ;
  \item une fréquence de Compton, déterminant la cadence de réaffirmation de sa cohérence interne, et par conséquent son temps propre élémentaire.
\end{itemize}

Ces quantités permettent d’introduire une densité de cohérence associée à un nucléon, interprétable comme un coût relationnel par unité de volume de fermeture minimale et par cycle de pulsation propre. Dans une configuration où la densité de nucléons est spatialement inhomogène, la densité de cohérence locale varie en corrélation, ce qui impose une réorganisation correspondante de la métrique émergente. Plus la densité de cohérence est élevée dans une région donnée, plus la métrique y présente une courbure prononcée, au sens où les trajectoires de coût minimal en cohérence y sont davantage défléchies.

La constante de couplage gravitationnel peut dès lors être définie de manière plus rigoureuse. Elle doit assurer la mise en relation, d’une part, de la variation spatiale de la densité de cohérence, pondérée par le paramètre de fuite minimale \(\varepsilon\) propre à chaque proton, et, d’autre part, de la courbure effective de la métrique émergente. Une formulation naturelle de cette relation consiste à combiner les éléments suivants :
\begin{itemize}
  \item l’action élémentaire \(\hbar\), déjà interprétée comme la granularité minimale de l’action associée à un régime stationnaire logique élémentaire ; 
  \item la vitesse limite \(c\), qui fixe l’échelle fondamentale de propagation de l’information. Dans le cadre relationnel, cette vitesse limite ne caractérise pas uniquement la propagation de la lumière, mais également la vitesse maximale à laquelle les réajustements de cohérence peuvent se propager au sein de la métrique émergente, ce qui en fait la borne supérieure de toute influence causale ; 
  \item la masse du proton \(m_p\), qui représente le coût global de cohérence hiérarchique associé à cette structure hadronique ; 
  \item la constante de structure fine \(\alpha\), qui encode la contribution des relations actives de surface. Dans le modèle LDP, ces relations actives de surface correspondent à un nombre fini de liens logiques demeurant disponibles pour des couplages externes après saturation de la cohérence interne du proton ; 
  \item un comptage détaillé de ces liens, fondé sur la hiérarchie quarks de valence / mer de quarks ainsi que sur la géométrie de la fermeture, conduit à une surface active d’environ \(502\) relations, permettant ainsi de réinterpréter \(\alpha\) comme un rapport de cohérence de surface ; 
  \item le paramètre de fuite \(\varepsilon\), qui quantifie la violation logique minimale de la fermeture hiérarchique. 
\end{itemize}

Après avoir clarifié le rôle de \(\hbar\) et de \(c\) au niveau du régime stationnaire logique élémentaire, puis celui de la constante de structure fine \(\alpha\) dans la structuration des interactions de surface, il devient possible d’examiner de façon systématique les conditions sous lesquelles une constante de couplage gravitationnel peut émerger de la fermeture hiérarchique imparfaite des hadrons.

Dans ce cadre, une expression possible pour la constante de couplage gravitationnel s’écrit :
\[
G_{\mathrm{LDP}} = \frac{2\,\hbar\, c}{m_p^2}\,\frac{1}{\alpha^{18\left(1-\epsilon\right)}}
\]
où l’exposant \(18\left(1-\epsilon\right)\) est interprété comme le nombre minimal d’itérations fractales nécessaires à la propagation, des échelles protoniques vers les échelles macroscopiques, de la déformation induite par la fuite de cohérence. Cette interprétation repose sur un schéma hiérarchique détaillé de la structure interne du proton, pour lequel le décompte combinatoire précis n’est pas présenté ici et devra faire l’objet d’une analyse spécifique dans un travail ultérieur.

Cette relation ne procède pas d’un ajustement libre à la valeur expérimentale de \(G\), mais d’une combinaison contrainte des paramètres mis en évidence dans les sections précédentes : \(\hbar\), qui quantifie la granularité de l’action par cycle ; \(c\), qui fixe la vitesse limite de propagation des réajustements de cohérence ; \(m_p\), qui caractérise le coût de cohérence hiérarchique associé au proton ; \(\alpha\), qui mesure la fraction de cohérence de surface disponible pour les couplages ; et \(\epsilon\), qui encode la fuite minimale imposée par la hiérarchie composite. Le facteur numérique \(18\) reflète ainsi le nombre minimal d’itérations hiérarchiques requis pour que la fuite associée à un proton se transmette jusqu’aux échelles macroscopiques ; son origine combinatoire détaillée dépasse le cadre du présent article et devra être explicitée dans un développement technique ultérieur.

Le facteur \(18\) ne constitue pas un paramètre d’ajustement libre : il représente le nombre minimal d’itérations hiérarchiques requises pour qu’une perturbation élémentaire de la cohérence, introduite par la fuite \(\epsilon\) au niveau d’un proton, se propage de manière itérative jusqu’aux échelles macroscopiques. Ce nombre résulte du dénombrement systématique des groupes de relations actives intervenant dans la fermeture hiérarchique interne du proton (quarks de valence, mer de quarks et surface active), tel que défini dans la construction détaillée du modèle.

La procédure de comptage, ainsi que la façon dont elle conduit à \(18\) itérations effectives, est présentée de manière exhaustive dans le développement complet du modèle LDP, auquel le présent texte renvoie pour l’ensemble des considérations de nature combinatoire.

En insérant les valeurs expérimentales de \(\hbar\), \(c\), \(m_p\) et \(\alpha\), et en prenant \(\epsilon \approx 0{,}0060923\), on obtient numériquement :
\[
G_{\mathrm{LDP}} \simeq 6{,}6743\times 10^{-11}\,\mathrm{m}^3\,\mathrm{kg}^{-1}\,\mathrm{s}^{-2}.
\]

La définition de \(G_{\mathrm{LDP}}\) doit être appréhendée avec précaution. Elle ne résulte pas, à ce stade, d’une dérivation intégrale fondée sur les seuls axiomes du modèle, mais procède d’une combinaison contrainte de grandeurs déjà introduites : \(\hbar\), interprété comme une granularité d’action par cycle ; \(c\), envisagée comme la vitesse limite de propagation des réajustements de cohérence ; \(m_p\), représentant le coût de cohérence hiérarchique associé au proton ; \(\alpha\), caractérisant le degré de cohérence de surface ; et \(\varepsilon\), définissant le paramètre de fuite minimale attaché à la hiérarchie composite.

Le facteur exponentiel \(18\left(1-\varepsilon\right)\) formalise l’hypothèse selon laquelle une fuite microscopique, initialement spécifiée à l’échelle du proton, se propage à travers un nombre fini d’itérations hiérarchiques jusqu’à acquérir une signature observable à l’échelle macroscopique. À l’état actuel du développement du modèle, la valeur \(18\) ne découle toutefois pas d’un comptage combinatoire exhaustif : elle synthétise un schéma hiérarchique encore partiellement spécifié, qui devra faire l’objet d’une dérivation technique rigoureuse et détaillée dans des travaux ultérieurs.

De façon analogue, la valeur numérique de \(\varepsilon\) utilisée dans l’expression de \(G_{\mathrm{LDP}}\) n’est pas, pour l’instant, obtenue de manière autonome. Elle est fixée en imposant que \(G_{\mathrm{LDP}}\) reproduise la constante de gravitation déterminée expérimentalement, ce qui confère à \(\varepsilon\) le statut de paramètre ajusté sur les données observationnelles, plutôt que celui d’une quantité dérivée de la structure interne du proton. Toute dérivation future cohérente du modèle devra retrouver cette valeur comme conséquence des axiomes et de l’architecture hiérarchique postulés ; à défaut, l’expression proposée de \(G_{\mathrm{LDP}}\) ne pourra être considérée que comme une reformulation phenomenologique cohérente des observations empiriques, et non comme une prédiction fondamentale issue du formalisme.

En synthèse, \(G_{\mathrm{LDP}}\) fournit principalement un test de cohérence conceptuelle : il met en évidence qu’une constante de couplage gravitationnel peut être reconstruite, avec le bon ordre de grandeur, à partir des seuls éléments relationnels déjà introduits. Toutefois, il ne constitue pas encore une dérivation complète ni une prédiction autonome de la constante gravitationnelle.

À ce stade, la valeur numérique de \(\varepsilon\) ne découle donc pas d’un calcul combinatoire entièrement déductif. En pratique, elle est déterminée en imposant que l’expression de \(G_{\mathrm{LDP}}\) reproduise la valeur expérimentale de la constante de gravitation. Ce procédé de calibrage doit être interprété comme un test de cohérence interne du cadre théorique : il contraint l’intervalle de valeurs admissibles pour \(\varepsilon\) et fixe une cible quantitative précise que toute dérivation microscopique ultérieure devra nécessairement retrouver.

L’écart relatif obtenu est de l’ordre de \(10^{-4}\), ce qui demeure compatible avec les incertitudes actuelles sur la valeur de \(G\), ainsi qu’avec les approximations inhérentes à la description hiérarchique du proton.

Cette concordance ne constitue pas une validation définitive du modèle, mais elle met en évidence la plausibilité qu’une constante effective de couplage gravitationnel émerge quantitativement à partir :
\begin{itemize}
  \item de la structure interne du proton ; 
  \item de la distribution de ses relations actives ; 
  \item et d’un taux minimal de perte de cohérence \(\epsilon\) ; 
\end{itemize}
sans qu’il soit nécessaire de postuler l’existence d’une interaction gravitationnelle fondamentale indépendante de la hiérarchie relationnelle.

Une constante de couplage gravitationnel peut, en principe, être définie comme le facteur de proportionnalité entre : 
\begin{itemize}
  \item la variation de la densité de cohérence induite par l’introduction d’une structure hiérarchique élémentaire (par exemple un proton) dans une région donnée ;
  \item et la variation de la courbure effective de la métrique émergente, telle qu’elle serait inférée à partir de la déviation des trajectoires de moindre coût de cohérence dans cette région.
\end{itemize}

Dans une formulation plus strictement adossée aux cadres théoriques établis, cette constante joue un rôle analogue à celui de la constante de gravitation newtonienne, ou plus généralement de la constante de couplage gravitationnel en relativité générale : elle assure la correspondance fonctionnelle entre une grandeur caractérisant le contenu matériel (en l’occurrence, le coût de cohérence hiérarchique) et une déformation mesurable de la métrique. La distinction fondamentale réside toutefois dans le fait que, dans le cadre relationnel considéré, cette constante n’est pas introduite comme paramètre fondamental a priori ; elle émerge au contraire de la structure interne des nucléons et de la manière dont leurs fermetures hiérarchiques s’agrègent et s’articulent au sein du réseau de voisinages logiques.

Une paramétrisation numérique cohérente peut alors, en principe, être obtenue en combinant le nombre total de relations internes d’un nucléon, le rayon associé à sa longueur d’onde de Compton et sa fréquence propre, de manière à construire une quantité ayant les dimensions d’une constante de couplage gravitationnel. Le fait qu’une telle construction conduise à un ordre de grandeur compatible avec la valeur mesurée de la constante de gravitation suggérerait que l’interprétation gravitationnelle du coût de cohérence hiérarchique n’est pas seulement conceptuellement défendable, mais également quantitativement plausible.

\subsection{Statut des paramètres internes du modèle}

La construction présentée ci-dessus fait intervenir plusieurs paramètres numériques internes. Avant de conclure, il est nécessaire d’en expliciter le statut épistémologique, dans la mesure où ces paramètres ne jouent pas un rôle homogène dans la structure du modèle : 
le paramètre de fuite \(\epsilon\) ; 
le facteur hiérarchique \(18\) apparaissant dans l’exposant de \(\alpha\) ; 
les entiers associés aux architectures locales \((4, 6, 24, 28, 500, 10087, 11017,\) etc.\().\)

Certains de ces nombres découlent directement des axiomes et de la théorie des graphes : 
la structure minimale \(\left(4,6\right)\) résulte des contraintes de reproductibilité, de stabilité de la référence, d’autonomie logique et d’absence de hiérarchie explicite ; 
la formule \(r\left(n\right) = n\left(n-1\right)/2\) provient du choix de graphes complets pour décrire les régimes stationnaires locaux ; 
la nécessité d’une fuite de cohérence non nulle dans les structures composites (nucléons) est une conséquence de l’impossibilité de réaliser simultanément une fermeture maximale à la fois locale et globale.

D’autres entiers sont introduits comme solutions minimales de problèmes combinatoires, sans que leur unicité n’ait, à ce stade, été démontrée : 
les valeurs \(n_u = 24\) et \(n_d = 28\) pour les quarks de valence \(u\) et \(d\),  
le nombre effectif de relations de surface \(R_{\mathrm{surf}} \simeq 502\),  
le nombre total de relations internes du proton \(r_{\mathrm{total}} = 11017\).

Ces choix ne sont cependant pas arbitraires : ils résultent d’une exploration systématique des architectures compatibles avec les axiomes (complétude minimale, fuite minimale, fraction stationnaire \(\sim 9\%\), fraction dynamique \(\sim 91\%\)) et permettent de reproduire correctement plusieurs rapports empiriques (notamment la prédominance de l’énergie dynamique dans la masse du proton, ainsi que la valeur de \(\alpha\) avec une précision meilleure que \(1\%\)).  
Ils doivent néanmoins être considérés comme des conjectures fortement contraintes, susceptibles d’être raffinées ou d’être justifiées de manière plus rigoureuse dans le cadre d’un développement combinatoire détaillé.

Enfin, certains paramètres ne sont pas encore obtenus par dérivation strictement déductive, mais par ajustement sur des constantes expérimentales mesurées :  
la valeur numérique \(\epsilon \simeq 0{,}0060923\) est choisie de manière à ce que \(G_{\mathrm{LDP}}\) coïncide avec la valeur mesurée de la constante gravitationnelle newtonienne ;  
le facteur \(18\) dans l’exposant de \(\alpha\) encode le nombre minimal d’itérations hiérarchiques requis pour propager la fuite de cohérence depuis l’échelle protonique jusqu’aux échelles macroscopiques ; son origine combinatoire détaillée reste à expliciter dans un développement ultérieur.

Dans l’état actuel du modèle, ces paramètres doivent donc être considérés comme semi-fondamentaux :  
leur nécessité qualitative (existence d’une fuite \(\epsilon > 0\), existence d’un facteur d’amplification hiérarchique) est imposée par la structure logique du formalisme ;  
leur valeur numérique précise est, en revanche, déterminée en imposant la compatibilité avec les constantes physiques observées (en particulier la constante gravitationnelle \(G\)).

L’achèvement empirique du modèle devra montrer comment ces valeurs émergent du comptage combinatoire détaillé des relations actives au sein des structures hadroniques, sans recours à un ajustement empirique direct. Dans le présent travail, le calibrage joue le rôle de test de cohérence : il met en évidence qu’un même cadre relationnel permet de reproduire, à partir d’une architecture interne des nucléons, à la fois la constante de structure fine \(\alpha\) et une constante de couplage gravitationnel \(G_{\mathrm{LDP}}\) numériquement compatible avec \(G\). En synthèse, le modèle LDP s’articule autour de trois niveaux : une base de logique déductive rigoureuse, une architecture relationnelle conjecturale mais fortement contrainte en son noyau, et un calibrage empirique qui fournit une cible quantitative précise à toute élaboration théorique complète ultérieure.

\section{Constante cosmologique et cohérence résiduelle}

Dans le cadre de la relativité générale, la constante cosmologique \(\Lambda\) apparaît comme un terme additionnel dans les équations d’Einstein, équivalent à une densité d’énergie du vide homogène responsable d’une accélération de l’expansion de l’Univers. Les observations cosmologiques contemporaines indiquent que cette composante, communément interprétée comme une forme d’énergie sombre, contribue de manière dominante au budget énergétique total de l’Univers, tout en présentant une valeur numérique extrêmement faible en unités naturelles. Le contraste marqué entre cette petitesse observée et les estimations naïves de l’énergie du vide issues de la théorie quantique des champs constitue le « problème de la constante cosmologique ». 

Dans le modèle LDP, l’espace-temps n’est pas une entité fondamentale : il émerge de la compatibilité des fermetures hiérarchiques et du réseau de voisinages logiques qu’elles induisent. La gravitation y est réinterprétée comme la trace, dans une description métrique effective, d’une fuite de cohérence minimale \(\varepsilon\) amplifiée au travers de plusieurs niveaux hiérarchiques. Dans cette perspective, il est naturel de s’interroger sur la possibilité d’identifier la constante cosmologique à une densité minimale de cohérence résiduelle, distribuée dans le réseau relationnel global indépendamment de toute structure hiérarchique spécifiée (nucléons, atomes, galaxies). 

On peut alors formuler la conjecture selon laquelle \(\Lambda\) traduirait, au niveau de la géométrie effective, la présence d’un fond de cohérence qui n’est capturé ni par les fermetures locales (particules, noyaux), ni par les fermetures intermédiaires (étoiles, galaxies), mais qui subsiste comme tension résiduelle du réseau aux plus grandes échelles. Cette densité de cohérence du « vide logique » se manifesterait, dans la description métrique, sous la forme d’une courbure moyenne non nulle, y compris en l’absence de matière ordinaire, reproduisant ainsi le rôle d’une constante cosmologique positive dans les modèles de type \(\Lambda\)CDM.

À ce stade de développement du modèle, cette proposition demeure de nature essentiellement qualitative. Aucune formulation quantitative de \(\Lambda\) n’est encore établie : le LDP ne fournit pas de dérivation de la valeur numérique de la constante cosmologique et ne résout pas le problème, bien connu, de la hiérarchie entre l’énergie du vide prédite par les théories des champs et la valeur observée de \(\Lambda\). La proposition se limite à avancer que, dans un cadre où l’espace-temps est émergent, la constante cosmologique effective pourrait être interprétée comme un indicateur du degré de cohérence résiduelle du réseau relationnel à grande échelle, plutôt que comme un paramètre fondamental supplémentaire.

Le même paramètre de fuite \(\varepsilon\), initialement introduit pour assurer la fermeture de la structure protonique, réapparaît dans les expressions de \(G\) puis de \(H_0\), conduisant à une estimation de \(H_0 \approx 64{,}7\,\mathrm{km\ s^{-1}\ Mpc^{-1}}\), compatible avec l’intervalle des mesures cosmologiques actuelles. Cette réutilisation d’un même taux de fuite, depuis l’échelle nucléonique jusqu’aux échelles cosmologiques, constitue l’un des aspects les plus déconcertants du cadre théorique, en ce qu’elle met en corrélation une exigence de cohérence locale avec une grandeur caractérisant l’extension globale de l’Univers.

\section{Systèmes liés et gravitation effective}

Une fois introduite, la constante de couplage gravitationnel permet de formuler, au niveau effectif, la dynamique par laquelle des ensembles de nucléons se comportent comme des sources collectives de courbure pour la métrique émergente. Les corps composés, depuis les noyaux atomiques jusqu’aux objets macroscopiques, peuvent alors être décrits comme des structures hiérarchisées constituées d’agrégats de coûts de cohérence élémentaires. 

La masse gravitationnelle d’un système étendu ne se réduit pas à la somme des masses au repos de ses constituants pris isolément. Chaque niveau de structuration modifie la distribution des coûts de cohérence : certaines contributions se superposent constructivement, tandis que d’autres se compensent partiellement, en fonction de l’organisation des fermetures hiérarchiques. La constante de couplage gravitationnel paramètre précisément l’écart entre la somme naïve des coûts élémentaires et l’effet collectif effectivement mesurable sur la métrique émergente. 

Dans le cadre relationnel, la « force » gravitationnelle ressentie par une structure donnée n’est qu’une description phénoménologique condensée de la déviation de sa trajectoire de moindre coût de cohérence lorsqu’elle traverse une région où la densité de cohérence est spatialement inhomogène. La loi de la chute libre peut alors être reformulée comme suit : pour un couplage fixé, toutes les structures suivent localement les mêmes trajectoires de moindre coût, indépendamment de leur propre coût interne de cohérence, tant que celui-ci demeure négligeable devant celui du milieu environnant.

\section{Matière ordinaire et organisation hiérarchique}

La matière ordinaire peut être interprétée comme le résultat d’une organisation hiérarchique itérative de structures fermées, incluant successivement électrons, nucléons, atomes, molécules et agrégats macroscopiques. Chacun de ces niveaux organisationnels correspond à une nouvelle modalité de répartition du budget de cohérence entre des sous-structures stationnaires et des régimes dynamiques collectifs. 

À chaque échelle, une même logique structurante se manifeste : 
\begin{itemize}
  \item certains sous-systèmes instaurent des régimes de fermeture locale quasi stationnaires ; 
  \item un environnement relationnel plus flexible supporte les interactions et assure la cohésion globale ; 
  \item la somme hiérarchisée des coûts de cohérence contribue à définir une masse gravitationnelle effective. 
\end{itemize}

Un exemple élémentaire est fourni par l’hydrogène, considéré comme première fermeture atomique. L’atome d’hydrogène constitue la première manifestation hiérarchique dans laquelle un régime stationnaire minimal (l’électron) s’organise autour d’une structure composite (le proton) sans altérer la cohérence interne propre à chacun des deux constituants. Dans cette configuration, la surface effective du proton – entendue comme l’ensemble des relations de surface encore disponibles pour des couplages externes – joue un rôle déterminant : c’est sur cette « peau » logique que vient se stabiliser le régime stationnaire électronique. 

La constante de structure fine \(\alpha\) peut alors être interprétée comme encodant le rapport entre la cohérence interne totale du proton et la fraction de cette cohérence qui se projette à sa surface et demeure accessible au couplage avec l’électron. Dans la description coulombienne classique, le potentiel de liaison au rayon de Bohr s’écrit \(V_{\mathrm{Coulomb}}(r) = -k e^{2}/r\) et conduit à une énergie de l’ordre de \(-4{,}36\times 10^{-18}\,\mathrm{J}\). Dans le cadre relationnel, on peut introduire un potentiel effectif \(V_{\mathrm{LDP}}(r) = -\alpha h c/r\), où \(\alpha\) quantifie cette fraction de cohérence de surface, tandis que \(h c / r\) fixe l’échelle d’action et de propagation associée ; évalué au rayon de Bohr, ce potentiel donne \(V_{\mathrm{LDP}}(r_{\mathrm{Bohr}}) \approx -4{,}35\times 10^{-18}\,\mathrm{J}\). Cette proximité d’ordre de grandeur indique que l’interprétation relationnelle de \(\alpha\) comme mesure d’une cohérence de surface demeure, au niveau de principe, compatible avec l’énergie de liaison de l’hydrogène. 

Dans cette perspective, la séparation conceptuelle entre « matière » et « géométrie » perd de sa rigidité. Ce qui est identifié comme matière correspond à la part du budget de cohérence concentrée dans des fermetures hiérarchiques persistantes ; ce qui est interprété comme géométrie renvoie à la manière dont la métrique émergente se trouve déformée sous l’effet de cette concentration. Les deux descriptions apparaissent ainsi comme des représentations complémentaires d’une même organisation relationnelle sous-jacente.

\section{Indications relatives à l’échelle cosmologique}

À des échelles cosmologiques, la distribution de la densité de cohérence ne se rapporte plus à un ensemble limité de structures isolées, mais doit être considérée comme une propriété globale associée à l’ensemble des fermetures hiérarchiques présentes dans l’univers observable. La métrique émergente devient alors un objet global, dont les caractéristiques reflètent à la fois la répartition statistique des coûts de cohérence et leur organisation hiérarchique.

Une constante de couplage gravitationnel, définie à partir de la structure élémentaire des nucléons, détermine quantitativement l’intensité de la courbure de l’espace-temps induite, pour une organisation hiérarchique donnée de la matière, au sein de cette métrique globale. Réciproquement, les observables cosmologiques à grande échelle (telles que le taux moyen d’expansion, les propriétés statistiques de la formation des structures et la courbure effective) peuvent être interprétées comme des contraintes supplémentaires sur les modalités de distribution, d’organisation et de stabilisation des fermetures hiérarchiques au cours de l’évolution en temps propre cosmologique.

L’objectif n’est pas ici de construire une théorie cosmologique complète, mais de montrer que le cadre relationnel admet, en principe, une extension cohérente depuis l’échelle des particules jusqu’à celle de l’univers. Les mêmes notions fondamentales (coût de cohérence, densité de cohérence, métrique émergente, couplage gravitationnel) demeurent pertinentes, quel que soit le niveau de complexité hiérarchique considéré.

\section{Hydrogène et matière ordinaire}

L’atome d’hydrogène constitue la première forme atomique stable issue de la combinaison d’une structure hiérarchique composite (le proton) et d’une structure élémentaire minimale (l’électron). Dans le cadre relationnel, il ne s’agit pas d’un simple agrégat de particules, mais du premier compromis cohérent entre deux régimes de fermeture distincts, condition nécessaire à l’émergence de la matière ordinaire.

\subsection{Proton, électron et premier système atomique}

Dans ce modèle, l’électron réalise la structure relationnelle minimale \(\left(4,6\right)\), caractérisée par une fermeture parfaite, dépourvue de hiérarchie interne et exempte de fuite de cohérence : l’ensemble de la pulsation se referme sur la structure elle-même, ce qui définit un temps propre élémentaire directement associé à sa fréquence de Compton. À l’inverse, le proton incarne une fermeture hiérarchique composite, comprenant un cœur de valence, une mer relationnelle dense et une fraction de liens résiduels localisée à la surface. Cette structure se décompose en une part stationnaire de l’ordre de quelques pourcents et une part dynamique d’environ \(91\%\), traduisant un coût de cohérence nettement plus élevé.

L’hydrogène associe ces deux régimes au sein d’un système unique : un proton hiérarchiquement fermé, présentant une fuite minimale \(\epsilon\), entouré d’un électron de type \(\left(4,6\right)\), qui doit préserver sa propre fermeture tout en demeurant compatible avec la surface logique active du proton. Dans ce cadre, l’atome d’hydrogène ne constitue pas seulement le système atomique le plus simple, mais représente la première solution non triviale aux contraintes de coexistence entre deux structures stationnaires de natures différentes.

\subsection{Pourquoi l’électron ne tombe pas sur le proton}

Dans le cadre de la mécanique classique, un électron en orbite autour d’un noyau devrait émettre un rayonnement électromagnétique continu, perdre progressivement son énergie mécanique et, par conséquent, spiraler vers le noyau jusqu’à s’y écraser. Cette prédiction conduit à une instabilité fondamentale de l’atome d’hydrogène, sauf à postuler l’existence d’orbites « non rayonnantes » ad hoc, comme dans le modèle semi-classique de Bohr.  

Dans le formalisme LDP, la problématique se reformule en termes de structures de fermeture. Un effondrement de l’électron sur le proton impliquerait la destruction de la fermeture minimale \(\left(4,6\right)\), c’est‑à‑dire la rupture de l’intégrité logique de la structure élémentaire, ce qui est incompatible avec les axiomes de persistance. La configuration liée électron–proton ne peut donc pas évoluer vers un état dans lequel la fermeture minimale serait annulée, car une telle évolution violerait les contraintes structurelles imposées par le cadre théorique.

La stabilité de l’hydrogène apparaît alors comme la conséquence d’un état d’équilibre structurel. D’une part, l’électron ne peut réduire indéfiniment sa distance au proton sans entrer en contradiction avec sa propre structure de cohérence (associée à sa fermeture minimale). D’autre part, le proton ne peut accroître sans limite l’extension effective de sa surface logique sans compromettre sa fermeture hiérarchique interne. L’« orbite fondamentale » ne doit donc pas être interprétée comme une trajectoire dynamiquement protégée de manière énigmatique, mais comme la distance à laquelle les deux fermetures peuvent coexister et se maintenir simultanément, sans violation de leurs contraintes respectives de cohérence et de persistance.

\subsection{Potentiel effectif, rayon de Bohr et rôle de \texorpdfstring{\(\alpha\)}{alpha}}

Dans le cadre de la physique standard, l’énergie de liaison de l’état fondamental de l’atome d’hydrogène est égale à \(-13{,}6\ \mathrm{eV}\), tandis que le rayon de Bohr \(a_0\) définit l’échelle de longueur caractéristique à laquelle la probabilité de présence de l’électron est maximale. Dans l’approche relationnelle, ces quantités se réinterprètent comme le coût de cohérence requis pour maintenir un régime stationnaire de type \(\left(4,6\right)\) à une distance finie d’une structure hiérarchique composite, en intégrant à la fois la dilution tridimensionnelle des liens et la surface logique active du proton.

Le texte met en évidence qu’en réécrivant le potentiel coulombien à partir des conditions de clôture tridimensionnelle et d’un nombre fini de relations de surface \(R_{\mathrm{surf}} \simeq 502\), la constante de structure fine \(\alpha\) émerge comme le rapport mesurant la fraction de cohérence de surface disponible pour ce couplage. L’évaluation de ce potentiel effectif au rayon de Bohr conduit à une énergie de liaison \(E_{\mathrm{LDP}} \simeq 4{,}35\cdot 10^{-18}\,\mathrm{J}\), en accord, à l’intérieur des approximations considérées, avec la valeur standard \(-4{,}36\cdot 10^{-18}\,\mathrm{J}\) (soit \(-13{,}6\ \mathrm{eV}\)). Cette concordance indique que l’interprétation de \(\alpha\) comme fraction de cohérence de surface reste compatible avec les données expérimentales relatives à l’hydrogène.

\subsection{Hydrogène, seuil de complexité et matière ordinaire}

L’existence de l’hydrogène ne résulte pas de l’introduction d’une interaction supplémentaire arbitraire, mais de la possibilité de stabiliser un couplage entre une structure minimale (l’électron) et une première structure hiérarchique composite (le proton), sans rompre la fermeture organisationnelle de l’un ou de l’autre. Ce seuil correspond à la première forme d’organisation dans laquelle la cohérence interne d’une structure (le proton) est en mesure de soutenir de manière durable l’ancrage d’une autre structure (l’électron) à distance finie, moyennant une redistribution partielle de son « budget » de cohérence de surface.

La matière ordinaire, dans ses formes plus complexes, peut alors être interprétée comme une généralisation de ce motif élémentaire : l’ensemble des architectures atomiques et moléculaires dérive de la répétition et de la composition d’un même schéma logique fondamental, dans lequel des fermetures hiérarchiques internes laissent subsister des surfaces suffisamment actives pour stabiliser des régimes stationnaires de type électronique. L’hydrogène apparaît dès lors comme la première interface durable entre deux niveaux de fermeture, et comme la condition minimale d’émergence d’une chimie, ainsi que d’un monde matériel structuré à grande échelle.

\section{Neutron : symétrie, fatigue et instabilité}

La description qui suit propose une interprétation qualitative de la désintégration du neutron dans le cadre LDP. Elle demeure conjecturale et ne constitue pas, à ce stade, une dérivation quantitative complète.

Le neutron partage avec le proton le statut de nucléon composite, mais il met en œuvre une architecture hiérarchique symétrique de type \(\left(u,d,d\right)\) qui le place à la limite de la stabilité : stable lorsqu’il est intégré dans certains noyaux, mais instable à l’état isolé. Dans le cadre relationnel, cette différence de comportement ne se réduit pas à une simple permutation de quarks ; elle résulte plutôt d’une répartition distincte du coût de cohérence et des violations logiques au sein de la structure hiérarchique.

\subsection{Architecture relationnelle du neutron}

À l’instar du proton, le neutron est modélisé au moyen d’une hiérarchie structurée en trois groupes stationnaires (correspondant aux quarks de valence) et un groupe dynamique (la mer de quarks), organisés sous les contraintes de complétude minimale et d’optimisation logique. La différence principale réside dans la symétrie interne : alors que le proton présente une configuration asymétrique \(\left(u,u,d\right)\), le neutron met en jeu une configuration symétrique \(\left(u,d,d\right)\), ce qui modifie à la fois la distribution des violations et les modalités de leur compensation par la mer de quarks. 

Dans ce cadre, les groupes associés au quark \(u\) et aux deux quarks \(d\) présentent des nombres d’entités et de relations comparables à ceux du proton (par exemple, 24 relations pour le groupe \(u\) et 28 pour chacun des groupes \(d\), tandis que la mer de quarks comporte de l’ordre de \(10^{4}\) relations dynamiques). Toutefois, la duplication du groupe \(d\) introduit une contrainte supplémentaire de symétrie interne, plus rigide. Le nombre total de relations internes atteint ainsi \(10\,992\), avec une fraction stationnaire d’environ \(9\%\) et une fraction dynamique de l’ordre de \(91\%\), comme dans le cas du proton, mais réparties de telle sorte que les quarks \(d\) occupent un rôle plus fortement contraint dans la fermeture relationnelle globale.

\subsection{Longueur d’onde de Compton et compression logique}

À l’instar du proton, la structure globale du neutron peut être caractérisée à partir de sa longueur d’onde de Compton, qui encode la pulsation propre associée à sa masse au repos. En adoptant la valeur expérimentale \(\lambda_{C,n} \simeq 1{,}31959\times 10^{-15}\,\mathrm{m}\), il est possible de définir un rayon logique associé à une onde de bouclage stationnaire, puis de le confronter à un rayon moyen « matériel » \(R_n \simeq 0{,}87\times 10^{-15}\,\mathrm{m}\), choisi sur la base des données issues de la physique nucléaire. 

Le calcul indique que le rayon déduit de la longueur d’onde de Compton du neutron est inférieur à ce rayon matériel, avec un écart relatif d’environ \(3{,}56\%\). Cette différence suggère que l’onde de fermeture impose une compression notable de la structure hiérarchique interne. Pour le proton, une analyse similaire conduit à un écart relatif de l’ordre de \(0{,}07\%\), soit plus de cinquante fois plus faible que dans le cas du neutron, ce qui suggère que la structure du neutron est soumise à une contrainte rythmique significativement plus forte.

\subsection{Fatigue de la fermeture et désintégration du neutron}

Dans le cadre LDP, cette compression logique se manifeste par une augmentation du coût de cohérence requis pour préserver la fermeture hiérarchique du neutron à chaque cycle de pulsation globale. La symétrie interne \(\left(u,d,d\right)\) redistribue les violations de manière moins favorable que la configuration asymétrique du proton, de sorte que la mer de quarks doit compenser une fraction plus importante des tensions internes, ce qui renforce la rigidité effective de la structure globale. 

Sur de longues échelles de temps propres, cette situation peut être interprétée comme un processus de fatigue de la fermeture. Les violations résiduelles, bien que minimisées à chaque cycle, ne peuvent être totalement annulées sans reconfiguration hiérarchique, ce qui rend la structure progressivement moins apte à maintenir sa cohérence interne. La désintégration d’un neutron libre \(\left(n \rightarrow p + e^{-} + \bar{\nu}_e\right)\) peut alors être comprise comme la résolution de cette impasse : la configuration symétrique cède et se réorganise en un proton plus stable, accompagné d’un électron et d’un antineutrino qui emportent la part de cohérence devenue incompatible avec la fermeture initiale.

\subsection{Neutron, noyaux et matière ordinaire}

Le caractère instable du neutron isolé contraste avec sa stabilité au sein d’un grand nombre de noyaux, où il peut subsister indéfiniment tant que la configuration globale demeure compatible avec les contraintes de cohérence structurelle. Dans le cadre relationnel, cette stabilisation s’interprète comme un phénomène collectif : la présence d’autres nucléons redistribue les contributions aux violations, de sorte que la compression intrinsèque du neutron est partiellement prise en charge par la métrique émergente et par les fermetures voisines, ce qui réduit la fatigue interne associée.

La différence de comportement entre proton et neutron ne se réduit donc pas à une simple variation de composition, mais reflète deux modalités distinctes d’organisation d’une fermeture hiérarchique composite sous les mêmes axiomes minimaux. À l’échelle macroscopique, cette asymétrie rend compte du fait que la matière ordinaire observable est dominée par les protons (qui, associés aux électrons, forment les atomes), tandis que les neutrons remplissent essentiellement une fonction de liant structurel au sein des noyaux, soumis à un compromis plus délicat entre stabilité locale et fatigue globale.

Le mécanisme proposé rend compte qualitativement de l’instabilité du neutron, mais les estimations actuelles de durée de vie qui en découlent demeurent très approximatives et ne sauraient se substituer aux calculs standards fondés sur l’interaction faible.

\section{Photon : triangle éphémère et transport d’information}

L’hypothèse d’identifier le photon à une structure minimale de type \(\left(3,3\right)\) constitue une proposition de travail hautement spéculative. Elle n’est, à ce stade, ni déduite formellement des axiomes du modèle, ni en mesure de reproduire l’ensemble des propriétés empiriques connues du photon, telles que la polarisation, le spin ou encore la description complète fournie par le cadre de l’électrodynamique quantique.

Dans le formalisme LDP, le photon occupe une position conceptuelle particulière : il ne se manifeste ni comme un volume fermé, à l’instar de l’électron, ni comme une structure hiérarchisée, comparable à celle des nucléons, mais plutôt comme le motif surfacique minimal apte à assurer le transport d’information et d’énergie entre deux fermetures. Son caractère intrinsèquement éphémère résulte directement des axiomes du modèle : un triangle logique ne saurait se stabiliser en une entité volumique autonome et ne peut exister que transitoirement, le temps d’assurer la transition entre deux configurations cohérentes distinctes.

\subsection{Triangle logique minimal}

Dans le cadre de l’architecture relationnelle, un système de trois entités en interaction constitue le premier motif de surface élémentaire : un triangle, apte à supporter une tension mais incapable d’induire, à lui seul, la clôture d’un volume stable. Alors que quatre entités (tétraèdre \(\left(4,6\right)\)) réalisent la condition minimale de fermeture volumique et, par conséquent, la formation d’une structure persistante (identifiée ici à l’électron), trois entités ne produisent qu’une configuration transitoire, intrinsèquement ouverte vers la dissipation.

L’hypothèse LDP propose d’identifier le photon à cette configuration temporaire de trois entités oscillant selon un triangle parfait, c’est-à-dire à « un instant de pure surface, dépourvu de toute structure de soutien permettant de porter durablement le poids du réel ». Ce triangle ne présente ni hiérarchie interne, ni dynamique de fond (ou « mer » interne), ni fermeture complète : il est entièrement soutenu par la tension triangulaire énoncée par l’axiome 2, et se révèle, de ce fait, intrinsèquement instable dès lors qu’il n’est plus continûment alimenté par la transition qui en a assuré la genèse.

\subsection{\texorpdfstring{\(E=\hbar\omega\)}{E = hbar omega} relu en LDP}

Dans la formulation standard de la physique quantique, la relation \(E = \hbar\omega\) associe à chaque photon une énergie proportionnelle à sa fréquence, sans que cette loi ne s’accompagne d’une interprétation structurelle plus approfondie. Dans le cadre LDP, cette relation est réinterprétée comme l’expression quantitative du coût de cohérence requis pour assembler et maintenir, pendant une durée finie, le triangle logique qui véhicule la transition entre deux fermetures.

Chaque photon est alors compris comme la manifestation d’une transition entre deux régimes stationnaires (par exemple, deux niveaux d’un électron lié ou deux états nucléaires), et la fréquence \(\omega\) encode le rythme de cette transition dans la métrique émergente. La constante de Planck réduite \(\hbar\), déjà interprétée comme la granularité minimale de l’action pour un régime stationnaire fermé, quantifie ici la quantité d’action par cycle allouée à la surface triangulaire ; l’énergie du photon se trouve ainsi reformulée comme le produit de cette granularité élémentaire par la cadence imposée au motif éphémère.

\subsection{Effet Compton : test de cohérence}

L’effet Compton constitue un banc d’essai direct pour cette interprétation relationnelle du photon. Dans la description standard, un photon de longueur d’onde incidente \(\lambda_{\mathrm{inc}}\) interagit avec un électron quasi libre et est diffusé avec une longueur d’onde augmentée \(\lambda_{\mathrm{diff}}\), suivant la relation
\[
\lambda_{\mathrm{diff}} - \lambda_{\mathrm{inc}} = \lambda_{C,e}\left(1-\cos{\theta}\right),
\]
où \(\lambda_{C,e}\) désigne la longueur d’onde de Compton de l’électron et \(\theta\) l’angle de diffusion.

Dans l’interprétation LDP, l’électron est modélisé comme un tétraèdre fermé \(\left(4,6\right)\) parfaitement cohérent (\(\epsilon = 0\)), supportant une onde stationnaire à la fréquence de Compton, tandis que le photon est représenté par un triangle rythmique qui vient effleurer la surface logique de cet électron. L’interaction du photon se traduit alors par une modulation locale de l’onde stationnaire de l’électron :
\begin{itemize}
  \item la violation triangulaire \(\nu\) augmente temporairement sur une région localisée de la surface (transition locale de \(\nu = 0\) à \(\nu > 0\)) ;
  \item une partie du rythme du photon est transférée à la pulsation interne de l’électron, qui enregistre ainsi une impulsion logique (interprétée comme le « recul ») ;
  \item le triangle réémerge avec une fréquence réduite \(\omega' < \omega\), correspondant à une longueur d’onde accrue, tandis que la fermeture globale de l’électron est restaurée (\(\nu \rightarrow 0\) après relaxation).
\end{itemize}

Les lois de conservation de la description standard (énergie et quantité de mouvement) se réinterprètent alors comme la conservation du budget d’action \(\hbar\) et de la cohérence de surface sur la sphère électronique : la variation de fréquence du photon est exactement égale à celle requise pour préserver l’intégrité de la structure \(\left(4,6\right)\) tout en accommodant la perturbation.

\subsection{Spectre électromagnétique et transitions}

Dans ce cadre, le spectre électromagnétique ne constitue pas une simple énumération arbitraire de fréquences possibles, mais représente l’ensemble des transitions compatibles avec les contraintes de cohérence imposées par les fermetures impliquées. Chaque photon observable est alors interprété comme la manifestation d’une transition entre deux structures stationnaires (atomiques, nucléaires ou de complexité supérieure) susceptibles de se reconfigurer sans altérer leurs fermetures internes, en externalisant la différence de cohérence sous la forme d’un triangle logique caractérisé par une fréquence \(\omega\). 

La relation \(E = \hbar\omega\) assure que toute différence d’énergie entre deux états stationnaires peut être prise en charge par une surface triangulaire unique, sans introduction de paramètre supplémentaire : la fréquence du photon est intégralement déterminée par la transition considérée, et la flexibilité du motif triangulaire permet, en principe, de rendre compte de l’ensemble des fréquences physiquement observables, des ondes radio jusqu’aux rayons gamma. Le photon apparaît alors comme l’unité élémentaire, mobile, du spectre électromagnétique, non pas en tant que particule dotée d’une extension volumique, mais comme un contrat transitoire établi entre deux fermetures, qui s’évanouit dès que la tension logique qu’il véhicule est absorbée ou réallouée au sein du réseau relationnel. 

Cette construction doit par conséquent être considérée comme une hypothèse de travail exploratoire plutôt que comme une théorie achevée : une formalisation complète devrait expliciter la manière dont les équations de Maxwell et la structure de l’électrodynamique quantique (EDQ) émergent effectivement de ce cadre relationnel.

\section{Constantes, effets mesurables et statut du LDP}

Dans le cadre du LDP, les constantes physiques ne sont pas considérées comme de simples paramètres externes ajustés a posteriori sur les données expérimentales, mais comme des invariants quantitatifs caractérisant différentes granularités de cohérence. Elles fonctionnent comme des \emph{signatures} de niveaux distincts de structuration des processus physiques : granularité d’action associée à un cycle de pulsation, granularité de couplage de surface caractérisant une structure composite, granularité thermique au sein d’un ensemble statistique, ou encore granularité de fuite gravitationnelle vers le fond relationnel global. 

Cette section recense les principales constantes mobilisées dans le texte et explicite leur statut conceptuel dans le cadre relationnel, en les interprétant systématiquement comme des marqueurs de cohérence et de structuration plutôt que comme des paramètres libres.

\subsection{\texorpdfstring{\(\hbar\)}{hbar} : granularité minimale de l’action}

Dans le cadre de la description relationnelle, la constante de Planck réduite \(\hbar\) est interprétée comme la granularité minimale de l’action associée à un régime stationnaire logique élémentaire. Chaque cycle de pulsation interne d’une structure minimale de type \(\left(4,6\right)\), identifiée à l’électron, transporte un quantum d’action d’amplitude \(\hbar\). L’énergie propre d’une telle structure s’exprime alors sous la forme \(E = \hbar\omega\) dès lors qu’une réalisation physique de cette pulsation est disponible.

Dans cette perspective, la relation \(E = \hbar\omega\) ne traduit pas l’ajout d’une quantité d’énergie à une entité préexistante, mais représente plutôt le coût de cohérence associé au maintien d’un régime stationnaire élémentaire. La constante \(\hbar\) détermine ainsi le pas minimal par lequel la cohérence interne se referme sur elle-même. Elle doit alors être comprise comme une constante de granularité logique, plutôt que comme un paramètre empirique introduit de manière ad hoc.

\subsection{\texorpdfstring{\(\alpha\) et \(\varphi\)}{alpha et phi} : cohérence de surface et optimisation géométrique}

Dans ce cadre théorique, la constante de structure fine \(\alpha\) peut être interprétée de manière analogue à une mesure de granularité de couplage de surface, plutôt que comme un paramètre fondamental dépourvu d’origine. Elle quantifie la fraction du budget de cohérence interne d’un proton susceptible de se manifester effectivement à sa surface lors de l’établissement d’un couplage stationnaire avec un électron.

Une première estimation, fondée sur une surface active de l’ordre de \(R_{\mathrm{surf}} \simeq 500\) relations, conduisait déjà à une valeur proche de \(\alpha^{-1}\) via la relation
\[
\alpha^{-1} \simeq \frac{R_{\mathrm{surf}}^{2}}{\mu},
\]
où \(\mu = m_p/m_e\) désigne le rapport de masse proton/électron. Un raffinement de cette estimation met en évidence une structure interne plus détaillée : le nombre total de relations stationnaires de valence au sein du proton est \(r_{\mathrm{valence}} = 930\), et le tiers de cette quantité, \(r_{\mathrm{valence}}/3 = 310\), joue le rôle de noyau moyen de cohérence associé à chaque quark de valence.

En multipliant ce tiers par le nombre d’or \(\varphi = (1+\sqrt{5})/2\), on obtient une surface active affinée
\[
R_{\mathrm{surf}}^{(\varphi)} \simeq \varphi \,\frac{r_{\mathrm{valence}}}{3} \approx 501{,}59,
\]
entièrement dérivée de paramètres internes à l’architecture protonique, complétés par un facteur géométrique simple. Insérée dans la relation \(\alpha^{-1} \simeq \left(R_{\mathrm{surf}}^{(\varphi)}\right)^{2}/\mu\), avec \(\mu \simeq 1836{,}15\), cette expression conduit à \(\alpha^{-1} \simeq 137{,}022\), en accord, à mieux que quelques \(10^{-4}\) en valeur relative, avec la valeur expérimentale \(\alpha^{-1} \simeq 137{,}036\), et ce sans introduction de terme ad hoc.

La constante de structure fine peut ainsi être réécrite sous la forme
\[
\alpha^{-1} \simeq \frac{\varphi^{2}}{\mu}\left(\frac{r_{\mathrm{valence}}}{3}\right)^{2},
\]
expression qui ne dépend que de \(\mu\), de \(r_{\mathrm{valence}}\) et du facteur \(\varphi\). Dans cette perspective, le nombre d’or n’est pas introduit comme un motif esthétique, mais comme un candidat à une constante d’optimisation géométrique relationnelle : il paramètre la répartition optimale du contenu de cohérence d’un noyau stationnaire moyen vers la surface active, sous des contraintes de fermeture hiérarchique et de minimisation des fuites de cohérence.

Cette interprétation s’inscrit dans une intuition déjà présente dans la littérature, selon laquelle la constante de structure fine relie un rapport de masses sans dimension \(\mu\) à un nombre fini de canaux de couplage entre entités. Dans le cadre proposé, cette idée se reformule naturellement en termes de surfaces actives et de topologie relationnelle : \(\alpha\) mesure la fraction de la granularité d’action fixée par \(\hbar\) qui demeure disponible à la surface pour les couplages, une fois assurée la fermeture interne hiérarchique de la structure protonique.

\subsection{\texorpdfstring{\(k_B\)}{kB} : granularité thermique et hiérarchie protonique}

La constante de Boltzmann \(k_B\) joue un rôle directement analogue dès lors que l’on ne considère plus une structure isolée, mais un ensemble statistique de structures hiérarchiques. Dans une lecture relationnelle, l’énergie thermique moyenne \(E_{\mathrm{th}} \sim k_B T\) quantifie la fraction du budget global de cohérence effectivement mobilisée, en moyenne, dans les excitations collectives d’un grand nombre de cycles de pulsation.

Dans ce cadre, la température cesse d’apparaître comme une grandeur primitive : elle mesure la population effective de cycles de cohérence excités au sein d’un ensemble de structures (électrons, nucléons, atomes). Le rôle de \(k_B\) est alors d’établir la correspondance entre ce comptage statistique et une échelle énergétique moyenne, de manière analogue à \(\hbar\), qui relie un cycle élémentaire à un quantum d’action. Autrement dit, \(k_B\) fixe l’échelle énergétique caractéristique d’une fluctuation thermique par degré de liberté.

Dans l’état actuel du LDP, \(k_B\) n’est pas encore dérivée de manière univoque, mais une expression candidate est proposée à partir de l’architecture conjointe de l’électron et du proton. Cette expression met en relation : la tension rythmique minimale \(h/\tau_0\) associée à un cycle de pulsation électronique, le nombre de relations actives de la structure \(\left(4,6\right)\), le rapport de masse \(\mu\), la fraction dynamique \(R_{\mathrm{mer}}/R_{\mathrm{tot}} \simeq 91\%\) de la mer de quarks, ainsi que le taux de fuite \(\varepsilon\), qui mesure la part de cohérence dissipée dans le fond relationnel global.

Sous ces hypothèses, on obtient une expression dépourvue de paramètre continu libre supplémentaire,
\[
k_B^{\mathrm{LDP}}
= \frac{h}{\tau_0}\,
  \frac{n_{\mathrm{rel}}^{e}}{\mu^{2}}\,
  \frac{R_{\mathrm{mer}}}{R_{\mathrm{tot}}}\,
  \frac{\varepsilon}{4}
\]
où chaque facteur renvoie à une propriété physiquement interprétable de l’architecture électronique et protonique. Numériquement, cette expression reproduit la valeur expérimentale de \(k_B\) avec une erreur relative contrôlée, ce qui suggère que la même hiérarchie relationnelle qui fixe \(\alpha\) pourrait également contraindre l’échelle thermique effective.

\subsection{\texorpdfstring{\(G\)}{G} et \texorpdfstring{\(\varepsilon\)}{epsilon} : couplage gravitationnel effectif}

Dans le cadre du LDP, la constante de gravitation \(G\) est interprétée comme la manifestation métrique d’un mécanisme de couplage plus fondamental : une fuite de cohérence minimale qui relie les structures hiérarchiques au fond relationnel global. La structure hiérarchique interne d’un proton ne peut, en effet, satisfaire simultanément la condition de fermeture maximale de l’ensemble de ses sous-systèmes et celle de fermeture globale de la structure considérée dans son intégralité. Cette contrainte structurelle induit une fuite de cohérence \(\varepsilon\), extrêmement faible mais strictement non nulle, vers le cadre relationnel environnant. 

Le paramètre \(\varepsilon\) est défini comme la déviation relative entre une configuration de fermeture hiérarchique idéale et la fermeture effectivement réalisée par la structure protonique, en tenant compte de ses \(11\,017\) relations internes ainsi que de la distinction entre fraction stationnaire et fraction dynamique. Il établit ainsi un lien quantitatif entre l’architecture interne du proton, la courbure de la métrique émergente et une constante de couplage gravitationnel effective, qui peut être interprétée comme une constante \(G_{\mathrm{LDP}}\) dérivée, plutôt que postulée a priori. 

L’expression de \(G_{\mathrm{LDP}}\) demeure à ce stade conjecturale et n’est que partiellement contrainte par des ordres de grandeur cosmologiques, par exemple via le paramètre de Hubble. Néanmoins, son statut conceptuel est bien défini : il ne s’agit pas d’un paramètre libre introduit ad hoc pour ajuster l’intensité des interactions gravitationnelles, mais de la reformulation, dans le langage métrique, d’une fuite logique minimale imposée par l’impossibilité de saturer simultanément l’ensemble des conditions de fermeture hiérarchique.

\subsection{Effets observables et domaine de validité du LDP}

Dans le cadre du LDP, les constantes fondamentales \(\hbar\), \(\alpha\), \(k_B\) et \(G\) reçoivent une interprétation unifiée en tant que manifestations distinctes d’un même « budget de cohérence ». Plus précisément, \(\hbar\) est comprise comme une granularité élémentaire de l’action, \(\alpha\) comme une fraction de cohérence allouée à la surface d’interaction électromagnétique, \(k_B\) comme une échelle moyenne d’excitation collective des degrés de liberté, et \(G\) comme la fuite minimale de cohérence vers le fond relationnel. Dans ce contexte, le nombre d’or \(\varphi\) apparaît comme un candidat à un invariant de topologie logique, caractérisant une organisation optimale des relations de cohérence entre un noyau stationnaire et une surface active, plutôt que comme un simple artefact numérologique. 

Sur le plan empirique, ces constantes demeurent, bien entendu, des quantités déterminées par la mesure expérimentale. Le formalisme LDP ne vise pas à supprimer la nécessité de la mesure, mais à proposer des expressions combinatoires susceptibles de reproduire leurs valeurs expérimentales avec des erreurs relatives typiquement comprises entre \(10^{-4}\) et \(10^{-2}\). La portée du modèle doit par conséquent être appréciée avec prudence : il ne fournit pas encore de dérivations strictement uniques, mais met en évidence que les mêmes entiers relationnels (\(4\), \(6\), \(24\), \(28\), \(930\), \(10\,087\), \(11\,017\)), le même rapport \(\mu\), le même facteur \(\varphi\) et le même paramètre de fuite \(\varepsilon\) peuvent organiser, de manière interne­ment cohérente, les constantes d’action, de couplage électromagnétique, de couplage gravitationnel et de granularité thermique.

\section{Perspectives et tests possibles}

Le cadre LDP n’a pas pour vocation de se substituer aux théories établies dans leurs domaines de validité respectifs. Il propose plutôt une relecture de leurs objets fondamentaux (particules, constantes, métrique) en tant que manifestations d’une logique minimale de cohérence. Dans cette perspective, la question centrale n’est pas de déterminer s’il reproduit d’emblée, à la précision expérimentale, l’intégralité des valeurs connues, mais d’évaluer s’il ouvre un programme de tests susceptible de falsifier ou de raffiner ses hypothèses combinatoires.

Les sections précédentes ont mis en évidence plusieurs résultats numériques notables : 
\begin{itemize}
  \item l’identification de la structure \(\left(4,6\right)\) à l’électron et de la hiérarchie \(\left(24,28,930,10087,11017\right)\) au proton ; 
  \item une relecture de la constante de structure fine \(\alpha\) à partir d’une surface active \(R_{\mathrm{surf}}^{\left(\varphi\right)} \simeq \varphi\, r_{\mathrm{valence}}/3\), conduisant à un accord meilleur que quelques \(10^{-4}\) en valeur relative ; 
  \item une expression candidate pour \(k_B\), construite à partir des mêmes entiers, de la fraction dynamique de la mer et du taux de fuite \(\varepsilon\), compatible avec la valeur mesurée à l’intérieur d’une incertitude contrôlée ; 
  \item une interprétation du couplage gravitationnel effectif en termes de fuite minimale de cohérence, ajustée sur les ordres de grandeur cosmologiques. 
\end{itemize}

Ces accords ne doivent pas être interprétés comme des prédictions exactes, mais comme des indicateurs de cohérence interne : ils suggèrent que les mêmes briques combinatoires peuvent organiser plusieurs constantes adimensionnées. Les écarts résiduels, typiquement de l’ordre de \(10^{-4}\) à \(10^{-2}\), ne constituent pas des détails négligeables ; ils offrent au contraire une voie de test empirique : s’ils se révèlent systématiques et corrélés, ils pourraient traduire des effets de fuite de cohérence ou de structure hiérarchique qui ne sont pas encore modélisés de manière suffisamment fine.

\subsection{Tests sur la structure \texorpdfstring{\(\left(4,6\right)\)}{(4,6)} et les hiérarchies internes}

Un premier axe d’investigation concerne la robustesse de la structure minimale \(\left(4,6\right)\) et de la hiérarchie protonique \(\left(24,28,930,10087,11017\right)\). Dans le cadre du LDP, ces entiers ne sont pas introduits de manière arbitraire, mais résultent de contraintes de complétude minimale, de cohérence triangulaire et d’optimisation logique. Toute mise en défaut empirique de cette architecture remettrait donc directement en cause la validité du cadre théorique sous-jacent.

Plusieurs orientations expérimentales et numériques peuvent être envisagées :
\begin{itemize}
  \item \textbf{Spectroscopie de très haute précision de systèmes liés simples} (atome d’hydrogène, hydrogène muonique, ions hydrogénoïdes) : si la constante de structure fine \(\alpha\) est effectivement reliée à \(r_{\mathrm{valence}}\), à \(\varphi\) et à \(\mu\), les corrections fines et hyperfines devraient pouvoir être réinterprétées comme des déformations de l’architecture de surface \(R_{\mathrm{surf}}^{\left(\varphi\right)}\), plutôt que comme des paramètres indépendants. De faibles écarts systématiques par rapport aux prédictions du modèle standard pourraient alors fournir une estimation effective de corrections hiérarchiques non prises en compte dans la construction actuelle.
  \item \textbf{Extensions aux nucléons et aux noyaux plus lourds} : l’analyse présentée s’est focalisée sur le proton en tant que structure composite minimale. Une généralisation naturelle consiste à appliquer la même méthodologie combinatoire à d’autres nucléons (par exemple le neutron) ainsi qu’à des noyaux légers, en recherchant des architectures hiérarchiques satisfaisant les axiomes du LDP tout en reproduisant les rapports de masses, les rayons ainsi que les fractions d’énergie stationnaire et dynamique. Si les entiers sélectionnés pour ces structures ne permettent pas de conserver le rôle central de \(\varphi\) et de \(\varepsilon\), le schéma proposé pour \(\alpha\), \(k_B\) et \(G\) devrait être remis en question.
  \item \textbf{Simulations combinatoires} : une exploration numérique systématique de l’espace des configurations \(\left(n_u,n_d,r_{\mathrm{valence}},R_{\mathrm{mer}}\right)\) sous les axiomes de cohérence et de minimalité peut être menée, afin de vérifier que la paire \(\left(24,28\right)\) et la valeur totale \(r_{\mathrm{valence}} = 930\) demeurent effectivement privilégiées, et que les architectures alternatives ne parviennent pas à reproduire simultanément la charge, la fraction dynamique et les rapports de masse observés.
\end{itemize}

\subsection{\texorpdfstring{\(\alpha\), \(\varphi\)}{alpha, phi} et les écarts résiduels}

La réinterprétation de \(\alpha^{-1}\) sous la forme
\[
\alpha^{-1} \simeq \frac{\varphi^{2}}{\mu}\left(\frac{r_{\mathrm{valence}}}{3}\right)^{2}
\]
repose sur un choix spécifique de surface active \(R_{\mathrm{surf}}^{\left(\varphi\right)}\) et sur l’hypothèse que \(\varphi\) joue le rôle d’un facteur d’optimisation géométrique de nature relationnelle. Les correspondances numériques obtenues sont quantitativement satisfaisantes, mais la présence d’un écart résiduel, même très faible, revêt une importance conceptuelle centrale : dans le cadre du LDP, cet écart n’est pas seulement acceptable, il est en principe interprétable.

Une interprétation possible consiste à considérer ces écarts comme la manifestation de corrections de surface induites par la fuite \(\varepsilon\) et par la dynamique de la « mer » de relations. Dans la construction actuelle, \(R_{\mathrm{surf}}^{\left(\varphi\right)}\) est formulée à partir d’entiers exacts \(\left(930\right)\) et de \(\varphi\), puis confrontée à une constante \(\alpha\) mesurée dans des configurations expérimentales spécifiques (états liés, vide quantique, procédures de renormalisation). Il est dès lors plausible que l’écart résiduel :
\begin{itemize}
  \item reflète des contributions provenant de niveaux hiérarchiques non explicitement intégrés au modèle (par exemple, des modes collectifs à plus grande échelle) ; 
  \item ou encode la manière dont la surface active est en réalité « rugueuse » ou sujette à des fluctuations, plutôt qu’idéalisée et lissée par le seul paramètre \(\varphi\).
\end{itemize}

Un programme de validation consiste alors à :
\begin{itemize}
  \item comparer systématiquement la valeur \(\alpha_{\mathrm{LDP}}\), déduite de \(R_{\mathrm{surf}}^{\left(\varphi\right)}\), à des déterminations expérimentales indépendantes de \(\alpha\), en fonction de l’échelle d’énergie, afin d’établir si les écarts reproduisent ou non la même dépendance en énergie (« course » de \(\alpha\)) que celle prédite par les schémas de renormalisation ; 
  \item analyser théoriquement comment des corrections de surface associées à \(\varepsilon\) ou à la fraction dynamique de la mer pourraient modifier \(R_{\mathrm{surf}}^{\left(\varphi\right)}\) d’une quantité quantitativement compatible avec les écarts observés.
\end{itemize}

Dans cette perspective, ce ne sont pas les valeurs centrales qui constituent le test le plus discriminant du LDP, mais la structure fine des écarts : leur amplitude, leur dépendance en échelle et leurs corrélations éventuelles avec d’autres constantes ou grandeurs dérivées.

\subsection{\texorpdfstring{\(k_B\)}{kB} : échelle thermique et statistiques relationnelles}

L’expression proposée pour \(k_B\) dans ce cadre relie de manière explicite :
\begin{itemize}
  \item la pulsation propre de l’électron (via \(h/\tau_0\)) ; 
  \item le nombre de relations caractérisant la structure minimale \(\left(4,6\right)\) ; 
  \item le rapport de masses \(\mu\) ; 
  \item la fraction dynamique \(R_{\mathrm{mer}}/R_{\mathrm{tot}}\) du proton ; 
  \item et le taux de fuite \(\varepsilon\), interprété comme la fraction de cohérence disponible pour l’agitation collective. 
\end{itemize}

Cette construction fournit une échelle thermique numériquement compatible avec la valeur expérimentale de \(k_B\), sans introduire de paramètre continu additionnel. Là encore, les écarts résiduels entre la valeur théorique et la valeur mesurée peuvent être exploités comme outil de test du modèle.

Plusieurs axes d’investigation se dégagent :

\begin{itemize}
  \item \textbf{Thermodynamique de systèmes simples} : réanalyser la capacité calorifique de systèmes modèles (gaz d’électrons faiblement couplés, gaz de nucléons dans un domaine de paramètres où les effets nucléaires sont contrôlés) en la formulant en termes de population de cycles de cohérence, et vérifier si l’expression \(k_B^{\mathrm{LDP}}\) rend correctement compte de la transition entre régimes quantique et classique.
  
  \item \textbf{Rôle de \(\mathbf{\varepsilon}\)} : dans la mesure où \(k_B^{\mathrm{LDP}}\) dépend explicitement de \(\varepsilon\), toute amélioration de l’estimation de ce taux de fuite (par exemple à partir de contraintes gravitationnelles ou cosmologiques) devrait induire une correction corrélée de la valeur de \(k_B\). Un désaccord systématique entre ces deux déterminations de \(\varepsilon\) constituerait ainsi un test direct de la cohérence du cadre théorique.
  
  \item \textbf{Fluctuations et entropie} : la relation usuelle \(S = k_B\ln{W}\) peut être reformulée en termes de nombre de configurations relationnelles accessibles à un ensemble donné de structures hiérarchiques. Une étude plus détaillée, fondée sur des modèles combinatoires explicites, permettrait de déterminer si l’échelle \(k_B^{\mathrm{LDP}}\) est effectivement adaptée pour relier le nombre de configurations relationnelles aux grandeurs thermodynamiques macroscopiques.
\end{itemize}

\subsection{Gravitation, fuite \texorpdfstring{\(\varepsilon\)}{epsilon} et tests cosmologiques}

Sur le plan gravitationnel, le LDP propose d’interpréter la courbure de la métrique émergente comme la trace, dans une description effective, de la distribution de densité relationnelle et du coût de cohérence qui lui est associé. La fuite minimale \(\varepsilon\), définie à partir des contraintes de fermeture du proton, joue dans ce cadre le rôle de paramètre de couplage entre les hiérarchies locales et le fond relationnel global.

Plusieurs catégories de tests potentiels peuvent être envisagées :
\begin{itemize}
  \item \textbf{Compatibilité avec la relativité générale} : dans les régimes où la relativité générale est validée de manière particulièrement précise (tests au sein du système solaire, pulsars binaires, signaux d’ondes gravitationnelles), il s’agit de vérifier que la reformulation relationnelle fournit au minimum un « relèvement » conceptuel cohérent, sans entrer en conflit avec les prédictions numériques établies. Toute divergence systématique impliquerait la nécessité de corriger la manière dont \(\varepsilon\) est relié à la courbure effective.
  \item \textbf{Ordres de grandeur cosmologiques} : le même paramètre \(\varepsilon\) intervient dans la tentative de relier l’architecture protonique à des échelles cosmologiques (par exemple via la densité d’énergie effective ou le paramètre de Hubble). Si ces corrélations se trouvaient confirmées, l’ajustement d’un unique \(\varepsilon\) devrait, en principe, permettre de rendre compte simultanément de l’intensité du couplage gravitationnel local et de certains ordres de grandeur cosmologiques.
  \item \textbf{Écarts résiduels sur \(\mathbf{G}\)} : la valeur de la constante de gravitation \(G\) est connue avec une précision significativement moindre que celle de constantes telles que \(\alpha\) ou \(k_B\). Le LDP incite à considérer les dispersions expérimentales non seulement comme des difficultés métrologiques, mais également comme une éventuelle fenêtre sur la structure relationnelle fine : des variations apparentes pourraient traduire des différences de densité relationnelle effective ou de composition hiérarchique des systèmes utilisés pour les mesures.
\end{itemize}

\subsection{Statut du LDP face aux observations}

Le LDP se situe encore au stade de cadre de fondation : plusieurs constructions restent conjecturales, certaines expressions (en particulier pour \(k_B\) et \(G\)) sont des hypothèses structurées plutôt que des dérivations définitives, et la prise en compte des phénomènes hors équilibre ou de la cosmologie détaillée demeure partielle. Il n’en propose pas moins un ensemble de tests convergents :
\begin{itemize}
  \item vérifier que les mêmes entiers relationnels et le même facteur \(\varphi\) peuvent organiser \(\alpha\), \(k_B\) et \(G\) dans des domaines d’observation distincts ; 
  \item analyser les écarts résiduels comme des estimations indépendantes de la fuite \(\varepsilon\) et des corrections hiérarchiques, plutôt que comme des erreurs arbitraires ; 
  \item étendre la construction à d’autres structures que l’électron et le proton, en cherchant à maintenir la cohérence du rôle assigné à \(\varphi\) comme invariant d’optimisation géométrique relationnelle.
\end{itemize}

De ce point de vue, la valeur du LDP ne se juge pas uniquement à la proximité de ses nombres avec les constantes mesurées, mais à la manière dont il transforme chaque écart, même minime, en question précise sur la structure du réel : quelle partie de la hiérarchie, de la surface active ou de la fuite de cohérence n’a pas encore été prise en compte ? 

\section{Genèse conjecturale des fermetures}

Le cadre LDP a été élaboré de manière délibérément acosmologique : il spécifie les conditions logico-structurelles d’existence de fermetures, sans présupposer de scénario cosmologique particulier ni d’histoire déterminée de l’univers. Dans cette perspective, les formes stables (électron, proton, noyaux, etc.) sont interprétées comme des solutions d’optimisation sous contraintes, susceptibles d’apparaître partout où un fond relationnel d’une richesse suffisante est disponible. Il demeure toutefois nécessaire d’esquisser un scénario dans lequel ces fermetures ne seraient pas posées comme données primitives, mais dériveraient de processus de contrainte à grande échelle.

On peut, de manière purement conjecturale, postuler qu’au-dessous de toute structuration métrique de la réalité se trouve un champ logique pré‑physique, assimilable à une matrice d’éléments informationnels pulsatoires, dans lequel ne subsistent que les axiomes de pulsation minimale et de cohérence. Au sein de ce champ, les pulsations élémentaires ne sont encore porteuses d’aucune organisation hiérarchique ; elles réalisent uniquement la possibilité qu’« il se produise quelque chose plutôt que rien ». Tant que la densité relationnelle demeure faible, les configurations restent essentiellement transitoires et ne donnent pas lieu à des fermetures durables. Ce n’est que lorsque certaines régions du champ sont soumises à des contraintes collectives d’intensité suffisante que des compositions plus structurées émergent : la plupart d’entre elles demeurent instables, mais certaines atteignent la fermeture minimale (photons, quarks, puis électrons et protons) et se maintiennent.

Dans cette perspective, ce que la cosmologie métrique standard décrit comme un Big Bang peut être réinterprété comme un épisode au cours duquel une région finie du champ logique pré‑physique franchit un seuil de contrainte, imposant la sélection rapide d’architectures fermées capables de soutenir une réalité ontologiquement stable. Il n’est pas nécessaire de postuler l’unicité d’un tel événement : rien, dans le cadre du LDP, n’exclut que plusieurs régions du fond relationnel subissent, à des « moments » distincts au sens émergent, des transitions analogues vers des familles de formes fermées. Du point de vue interne d’un observateur immergé dans l’une de ces branches, l’accès empirique serait limité à la branche cosmologique à laquelle appartiennent ses propres fermetures. 

Les trous noirs peuvent alors être reconsidérés, de façon tout aussi conjecturale, comme des dispositifs de contrainte extrême dans lesquels les hiérarchies relationnelles stabilisées cessent de pouvoir se maintenir sous forme de structures composites. À l’approche d’une densité relationnelle critique, la hiérarchie interne des formes se désagrège : la fermeture globale devient incompatible avec la cohérence triangulaire, et le « budget de cohérence » se libère sous forme d’éléments pulsatoires redevenant faiblement structurés. Dans un tel schéma, ce qui apparaît, dans le cadre de notre description métrique interne, comme la « disparition » de la matière dans un trou noir pourrait correspondre, dans le langage du LDP, à un processus de recyclage de la cohérence vers le champ pré‑physique, potentiellement susceptible d’alimenter, dans une autre région du fond relationnel, un nouvel épisode de type Big Bang. 

Une représentation plus radicale consiste à envisager que ce processus global forme une boucle dépourvue de commencement et de fin au sens strict. Dans le cadre du LDP, le temps n’existe qu’en tant que traduction métrique du comptage des cycles de cohérence au sein des structures fermées. En dehors de ces structures, c’est‑à‑dire dans le champ logique pré‑physique, il n’y aurait ni « avant » ni « après », mais uniquement des régimes de contrainte et de cohérence. Dans cette optique, les expressions « avant le Big Bang » ou « après l’évaporation d’un trou noir » ne possèdent de pertinence que relativement à des régimes stationnaires déjà capables de supporter une notion émergente de temps propre ; le fond relationnel, pour sa part, demeurerait atemporel.

Ce scénario de genèse et de recyclage des fermetures doit être considéré comme hautement conjectural. Il n’a pas vocation à se substituer aux cadres cosmologiques établis, ni à fournir une description quantitative de l’évolution de l’univers observable. Son ambition est plus limitée, mais potentiellement complémentaire : il vise à montrer que le LDP autorise, de manière naturelle, une représentation dans laquelle les formes ne se ferment pas de façon « spontanée », mais émergent sous l’effet d’événements de contrainte au sein d’un champ pulsatoire plus primitif, et dans laquelle les singularités gravitationnelles extrêmes pourraient intervenir dans des processus de recyclage de la cohérence. Cette construction cosmologique doit donc être comprise non comme une prédiction, mais comme une proposition exploratoire invitant à étudier comment les outils conceptuels du LDP pourraient, à terme, être articulés à des modèles plus élaborés de l’origine et de l’évolution des structures physiques.

La présente section ne se présente pas comme un modèle cosmologique quantitatif, mais comme une possible relecture de scénarios déjà discutés en cosmologie (univers cycliques, univers émergeant de trous noirs), reformulés dans le langage du LDP. Si cette représentation s’avérait heuristiquement féconde, l’architecture relationnelle commune qui organise \(\hbar\), \(\alpha\), \(k_B\) et \(G\) pourrait également fournir des conditions de cohérence et des contraintes sur la manière dont des épisodes de type Big Bang et des singularités gravitationnelles extrêmes se manifestent et se codent dans le fond relationnel.

\section*{Conclusion}
\addcontentsline{toc}{section}{Conclusion}

Le Logo Dynamique Projectif a été élaboré à partir d’une hypothèse de travail à la fois simple dans son énoncé et exigeante dans ses implications : reconstruire la possibilité même d’une réalité physique à partir d’axiomes purement relationnels, sans postuler a priori l’existence d’un espace, d’un temps ou de particules. Dans ce cadre, l’existence n’est plus associée à une substance ou à un support ontologique préalables, mais à la capacité de certaines structures de produire et de maintenir une distinction stable au sein d’un champ de pulsations élémentaires. La structure minimale 46, déduite des axiomes de pulsation, de cohérence triangulaire, de complétude et d’optimisation, joue alors le rôle de brique élémentaire : elle réalise la première fermeture logique admissible et fournit un prototype pour une interprétation strictement relationnelle de l’électron.

Sur cette base, le texte a montré comment des structures plus complexes peuvent émerger par superposition et hiérarchisation de fermetures minimales. Le proton, en particulier, a été décrit comme une architecture composite minimale, combinant trois régimes stationnaires locaux (quarks de valence), une mer relationnelle dynamique, une surface active finie et une fuite de cohérence \(\varepsilon\) vers le fond relationnel global. Les entiers caractérisant cette architecture, 24289301008711017, ne sont pas introduits comme paramètres libres : ils sont sélectionnés par la fonction d’optimisation logique sous des contraintes combinatoires strictes (multiples de 4, graphes complets locaux, fraction dynamique dominante, charge nette +1). Ils offrent par conséquent une base concrète pour réinterpréter plusieurs constantes fondamentales comme des signatures de cohérence interne, plutôt que comme des données exogènes à la structure.

Dans ce cadre, la constante de Planck réduite \(\hbar\) cesse d’être considérée comme un simple paramètre empirique dépourvu de fondement conceptuel : elle quantifie la granularité minimale de l’action associée à un cycle de pulsation d’un régime stationnaire élémentaire de type 46. La relation \(E=\hbar\omega\) peut alors être réinterprétée comme le coût de cohérence requis pour maintenir un tel régime. De manière analogue, la constante de structure fine \(\alpha\) reçoit une interprétation définie à l’échelle de la surface protonique : elle mesure la fraction du budget de cohérence interne qui demeure disponible, à la surface, pour des couplages externes. L’affinement introduit par le nombre d’or, via la définition de la surface active \(R_{\operatorname{surf}}^{\left(\varphi\right)} \simeq \varphi \, r_{valence}/3\), conduit à l’expression \(\alpha^{-1} \simeq \varphi^{2}(r_{valence}/3)^{2}/\mu\), laquelle reproduit la valeur expérimentale de \(\alpha\) avec une précision meilleure que quelques \(10^{-4}\) en valeur relative, et ce sans recours à des ajustements ad hoc. Dans ce cadre, le nombre d’or apparaît ainsi comme un candidat à un invariant d’optimisation géométrique relationnelle : il joue le rôle d’un facteur décrivant la répartition optimale de la cohérence entre le noyau stationnaire et la surface active, plutôt que celui d’un motif numérologique imposé a posteriori.

La même architecture logique permet de formuler des expressions candidates pour la constante de Boltzmann \(k_B\) et pour un couplage gravitationnel effectif \(G_{LDP}\). Dans l’interprétation relationnelle, \(k_B\) mesure une granularité thermique : il fixe l’échelle d’énergie moyenne associée à l’excitation statistique d’un grand nombre de cycles de cohérence, construite à partir de la pulsation propre de l’électron, du rapport de masses \(\mu\), de la fraction dynamique de la mer de quarks et du taux de fuite \(\varepsilon\). De son côté, le couplage gravitationnel émerge comme la traduction métrique de la fuite minimale de cohérence imposée par l’impossibilité, pour une hiérarchie composite, de saturer simultanément l’ensemble des fermetures locales et la fermeture globale. Dans tous ces cas, les concordances numériques obtenues restent approximatives, mais elles reposent systématiquement sur les mêmes ingrédients combinatoires : la structure 46, l’architecture protonique, le rapport \(\mu\), le facteur \(\varphi\) et le paramètre \(\varepsilon\).

Le présent texte ne vise pas, à ce stade, à proposer une théorie complète et définitivement structurée des constantes fondamentales, de la gravitation ou des phénomènes thermiques. Plusieurs constructions, en particulier celles qui concernent \(k_B\), \(G_{LDP}\) et la cosmologie, demeurent conjecturales ou seulement partiellement ajustées aux données expérimentales, et l’espace des configurations combinatoires n’a été exploré que de manière encore restreinte. La valeur épistémique du LDP doit donc être appréciée moins à l’aune de la précision numérique de ses prédictions qu’à celle de la cohérence structurelle du schème qu’il propose : un cadre conceptuel dans lequel les constantes sans dimension, les structures de particules, la métrique émergente et le temps propre se laissent interpréter comme des manifestations complémentaires d’une même logique minimale de cohérence. Les écarts résiduels par rapport aux mesures expérimentales, loin de constituer des détails négligeables, sont au contraire envisagés comme des marges de manœuvre permettant de tester et de raffiner le rôle attribué à la surface active, à la hiérarchie interne et au paramètre de fuite \(\varepsilon\). 

La section consacrée à la genèse conjecturale des fermetures esquisse, en outre, une possible articulation cosmologique du LDP. Dans ce cadre, un champ logique pré-physique de pulsations élémentaires fournit un substrat atemporel à partir duquel, sous l’effet de contraintes globales de cohérence, émergent localement des structures fermées stables, éventuellement via des épisodes de type Big Bang. Les trous noirs y sont réinterprétés comme des dispositifs de contrainte extrême au sein desquels les hiérarchies stabilisées se désagrègent, recyclant la cohérence vers le fond relationnel plutôt que vers une singularité géométrique non interprétable. Cette proposition cosmologique demeure délibérément spéculative : elle ne possède pas le statut d’un modèle quantitatif susceptible de se substituer aux scénarios établis, mais celui d’une interprétation possible de leurs grandes lignes dans le langage formel du LDP, en accord avec une lecture relationnelle des constantes.

Nous disposons désormais d’une traduction formelle des axiomes du Logo Dynamique Projectif (LDP) dans un langage de graphes signés et de dynamiques discrètes, au sein duquel la structure \(\left(4,6\right)\) émerge comme première structure stationnaire élémentaire. Dans ce cadre, les blocs \(u\), \(d\) ainsi que la structure protonique apparaissent comme des fermetures hiérarchiques sélectionnées par les mêmes principes de cohérence triangulaire, de complétude et d’optimisation. La combinatoire détaillée des blocs internes — en particulier la justification précise de la minimalité des entiers 24, 28, 276 et 378 — ainsi que la dérivation rigoureuse de certains rapports sans dimension demeurent toutefois des problématiques ouvertes, qui requerront des investigations mathématiques et numériques plus systématiques. Ces limitations doivent être interprétées comme des invitations à approfondir et à mettre à l’épreuve le cadre proposé, plutôt que comme des faiblesses définitives.

Le LDP doit ainsi être considéré comme une proposition de cadre fondationnel en cours de formalisation. Il fournit une méthode explicite pour poser la question de l’existence, pour dériver une structure minimale de cohérence, et pour articuler, à partir de celle-ci, les notions de particule, de constante adimensionnée, de temps propre, de gravitation effective et de genèse cosmologique. Les développements ultérieurs consisteront à raffiner les constructions présentées, à explorer de manière plus systématique l’espace des architectures hiérarchiques, et surtout à multiplier les points de contact quantitatifs avec les données empiriques, de sorte que le LDP puisse être confirmé, modifié ou invalidé par les phénomènes mêmes qu’il ambitionne de réinterpréter.

\end{document}